\documentclass[12pt, ukrainian]{article}
\usepackage[a4paper]{geometry}
\setlength\parindent{0mm}
%\setlength\parskip{4mm}
\geometry{top=1.3cm,bottom=1.3cm,left=1.5cm,right=1.5cm}
\pagestyle{empty}
\usepackage[T2A]{fontenc}
\usepackage[utf8]{inputenc}
\usepackage[ukrainian]{babel}
\usepackage{verbatim, listings, clrscode3e, fancyvrb}
\def\code#1{\Verb!#1!}
\begin{document}
%begin
{{\begin {scriptsize}\par
{ \center  Київський національний університет імені Тараса Шевченка \par}
{\parbox {5 cm}{Освітня програма \hfill}{\parbox {7 cm}{Інформатика\hfill}}\parbox {6 cm}{\hfill Семестр 2}}\par
{\parbox {5 cm} {Навчальний предмет \hfill}\parbox {7 cm}{Програмування \hfill}\parbox {6cm}{\hfill}\par}\vspace{-15pt} \end{scriptsize}}{\center Екзаменаційний білет \No \textbf { 1 }\par}}
{%beginitem
1. Вкажіть основну \textbf{семантичну} відмінність між конструктором копії та конструктором переміщення.%enditem
\par
\vspace{2mm}
%beginitem
2. Максимально змістовно повно сформулюйте назву оголошеного в наведеному фрагменті методу класу.
\vspace{-3mm}
\begin{verbatim}class A{
   ...
   ~A();
};\end{verbatim}
\vspace{-3mm}
%enditem
\par
\vspace{2mm}
%beginitem
3. Вкажіть, що не так з наведеним кодом.
\vspace{-3mm}
\begin{verbatim}int *pa=new int(2), *pb=*pa; delete pa;\end{verbatim}
\vspace{-3mm}
%enditem
\par
\vspace{2mm}
%beginitem
4. За наявності означень
\vspace{-3mm}
\begin{verbatim}
class A{ };
class B:public A{ };
class C:public A{B b;};\end{verbatim}
\vspace{-3mm}
вкажіть послідовність викликів конструкторів за виконання коду \code{\{C c;\}} 

(Якщо конструктор викликається кілька разів поспіль, то має зазначатися кожен його виклик, наприклад, A, A, A.)
%enditem

\vspace{2mm}
%beginitem
5. Яку часову оцінку складності має алгоритм швидкого сортування?
%enditem
\par
\vspace{2mm}
%beginitem
6. Модифікувати клас 

\qquad\code{class Session\{ int DM, MA, Alg, Prog;\};} 


додавши до нього оператор перетворення на тип \code{bool}: об'єкт класу є істинним тоді і тільки тоді, коли значення всіх атрибутів не менші за 60. Означення оператора теж має бути записане.%enditem
\par
\vspace{2mm}
%beginitem
7. 
Реалізувати операцію різниці множин з повтореннями. Множина з повтореннями задається переліком своїх елементів, записаних за неспаданням у текстовому файлі через білі символи. Результат у такому ж самому форматі (і теж за неспаданням) має бути записаний у текстовий файл. Питома функція має отримувати три шляхи до файлів. Розмір вхідних множин такий, що вони не можуть бути завантажені повністю в пам’ять програми.

Крім правильності та відповідності умові оцінюється якість коду, наявність витоків ресурсів та коректність роботи з пам'яттю. 

%enditem
}
{\smallskip\smallskip \begin {scriptsize} Затверджено на засіданні кафедри теоретичної кібернетики, протокол \No 4  від   25.01.2024 р.\par\smallskip\smallskip\smallskip\smallskip
{\parbox {6.5 cm} {Зав. кафедри \dotfill Ю.В.~Крак}}\hspace {0.7cm} {\hbox to 6.5 cm { Екзаменатор \dotfill Т.О.~Карнаух}} \end{scriptsize}}
%end
\newpage
%begin
{{\begin {scriptsize}\par
{ \center  Київський національний університет імені Тараса Шевченка \par}
{\parbox {5 cm}{Освітня програма \hfill}{\parbox {7 cm}{Інформатика\hfill}}\parbox {6 cm}{\hfill Семестр 2}}\par
{\parbox {5 cm} {Навчальний предмет \hfill}\parbox {7 cm}{Програмування \hfill}\parbox {6cm}{\hfill}\par}\vspace{-15pt} \end{scriptsize}}{\center Екзаменаційний білет \No \textbf { 2 }\par}}
{%beginitem
1. Що важливо забезпечити при роботі з динамічними змінними?%enditem
\par
\vspace{2mm}
%beginitem
2. Максимально змістовно повно сформулюйте назву оголошеного в наведеному фрагменті методу класу.
\vspace{-3mm}
\begin{verbatim}class A{
   ...
   A(A&&);
};\end{verbatim}
\vspace{-3mm}
%enditem
\par
\vspace{2mm}
%beginitem
3. Вкажіть, що не так з наведеним кодом.
\vspace{-3mm}
\begin{verbatim}int **ppi=new int*(new int(3)); delete ppi;\end{verbatim}
\vspace{-3mm}
%enditem
\par
\vspace{2mm}
%beginitem
4. За наявності означень
\vspace{-3mm}
\begin{verbatim}
class A{ };
class B{ };
class C{B b;A a;};\end{verbatim}
\vspace{-3mm}
вкажіть послідовність викликів конструкторів за виконання коду \code{\{C c;\}} 

(Якщо конструктор викликається кілька разів поспіль, то має зазначатися кожен його виклик, наприклад, A, A, A.)
%enditem

\vspace{2mm}
%beginitem
5. Яку часову оцінку складності має алгоритм сортування вставками?
%enditem
\par
\vspace{2mm}
%beginitem
6. Модифікувати клас 

\qquad\code{class Session\{ int DM, MA, Alg, Prog;\};} 


додавши до нього оператор порівняння \code{<}: один об'єкт є меншим за інший тоді і тільки тоді, коли значення всіх атрибутів об'єктів не менші за 60 і сума атрибутів першого менша за суму атрибутів другого. Означення оператора теж має бути записане.%enditem
\par
\vspace{2mm}
%beginitem
7. 
Задано клас A (структуру, оператори > та == десь коректно означено) та клас вузла списку Node.

struct A\{

\quad  якісь коректні означення 

\quad  bool operator>(const A\&); 

\quad  bool operator==(const A\&); 

\quad A\& operator =(const A\&)=delete; 

\quad  A\& operator =(A\&\&)=delete;\};

struct Node\{Node *next; A a;\};

 Написати функцію, яка вилучає зі списку всі вузли, елемент в яких менший елемента в попередньому вузлі списку або дорівнює елементу в наступному вузлі. Список задається вказівником на перший вузол списку. Наприклад, зі списку, що містить послідовність 8, 5, 6, 9, 4, мають бути вилучені вузли, що містять 5 та 4. У цьому прикладі 6 залишається в списку.

Крім правильності та відповідності умові оцінюється якість коду, наявність витоків ресурсів та коректність роботи з пам'яттю. 

%enditem
}
{\smallskip\smallskip \begin {scriptsize} Затверджено на засіданні кафедри теоретичної кібернетики, протокол \No 4  від   25.01.2024 р.\par\smallskip\smallskip\smallskip\smallskip
{\parbox {6.5 cm} {Зав. кафедри \dotfill Ю.В.~Крак}}\hspace {0.7cm} {\hbox to 6.5 cm { Екзаменатор \dotfill Т.О.~Карнаух}} \end{scriptsize}}
%end
\newpage
%begin
{{\begin {scriptsize}\par
{ \center  Київський національний університет імені Тараса Шевченка \par}
{\parbox {5 cm}{Освітня програма \hfill}{\parbox {7 cm}{Інформатика\hfill}}\parbox {6 cm}{\hfill Семестр 2}}\par
{\parbox {5 cm} {Навчальний предмет \hfill}\parbox {7 cm}{Програмування \hfill}\parbox {6cm}{\hfill}\par}\vspace{-15pt} \end{scriptsize}}{\center Екзаменаційний білет \No \textbf { 3 }\par}}
{%beginitem
1. Вкажіть основну відмінність між відкритими (public) та прихованими (private) методами класу.%enditem
\par
\vspace{2mm}
%beginitem
2. Максимально змістовно повно сформулюйте назву оголошеного в наведеному фрагменті методу класу.
\vspace{-3mm}
\begin{verbatim}class A{
   ...
   A& operator =(A&&);
};\end{verbatim}
\vspace{-3mm}
%enditem
\par
\vspace{2mm}
%beginitem
3. Вкажіть, що не так з наведеним кодом.
\vspace{-3mm}
\begin{verbatim}int *pa=new int(2), *pb=pa; delete pb; delete pa;\end{verbatim}
\vspace{-3mm}
%enditem
\par
\vspace{2mm}
%beginitem
4. За наявності означень
\vspace{-3mm}
\begin{verbatim}
class A{ };
class B:public A{ };
class C:public A{B b;};\end{verbatim}
\vspace{-3mm}
вкажіть послідовність викликів деструкторів за виконання коду \code{\{C c;\}} 

(Якщо деструктор викликається кілька разів поспіль, то має зазначатися кожен його виклик, наприклад, $\sim$A, $\sim$A, $\sim$A.)
%enditem

\vspace{2mm}
%beginitem
5. Яку часову оцінку складності має алгоритм сортування купою (він же деревом, пірамідою)?
%enditem
\par
\vspace{2mm}
%beginitem
6. Модифікувати клас точок простору 

\qquad\code{class Point\{ double x, y, z;\};}


додавши до нього оператор перетворення на тип \code{bool}: точка є істинною тоді і тільки тоді, коли містить хоча б одну ненульову координату. Означення оператора теж має бути записане.%enditem
\par
\vspace{2mm}
%beginitem
7. 
Задано клас A (структуру, оператори < та == десь коректно означено) та клас вузла списку Node.

struct A\{

\quad  якісь коректні означення 

\quad  bool operator<(const A\&); 

\quad A\& operator =(const A\&)=delete; 

\quad  A\& operator =(A\&\&)=delete;\};

struct Node\{Node *prev; Node *next; A a;\};

Написати функцію, яка вилучає зі списку всі вузли, елемент в яких менший елемента в попередньому вузлі списку або дорівнює елементу в наступному вузлі. Список задається вказівником на перший вузол списку. Наприклад, зі списку, що містить послідовність 8, 5, 6, 9, 4, мають бути вилучені вузли, що містять 5 та 4. У цьому прикладі 6 залишається в списку.

Крім правильності та відповідності умові оцінюється якість коду, наявність витоків ресурсів та коректність роботи з пам'яттю. 

%enditem
}
{\smallskip\smallskip \begin {scriptsize} Затверджено на засіданні кафедри теоретичної кібернетики, протокол \No 4  від   25.01.2024 р.\par\smallskip\smallskip\smallskip\smallskip
{\parbox {6.5 cm} {Зав. кафедри \dotfill Ю.В.~Крак}}\hspace {0.7cm} {\hbox to 6.5 cm { Екзаменатор \dotfill Т.О.~Карнаух}} \end{scriptsize}}
%end
\newpage
%begin
{{\begin {scriptsize}\par
{ \center  Київський національний університет імені Тараса Шевченка \par}
{\parbox {5 cm}{Освітня програма \hfill}{\parbox {7 cm}{Інформатика\hfill}}\parbox {6 cm}{\hfill Семестр 2}}\par
{\parbox {5 cm} {Навчальний предмет \hfill}\parbox {7 cm}{Програмування \hfill}\parbox {6cm}{\hfill}\par}\vspace{-15pt} \end{scriptsize}}{\center Екзаменаційний білет \No \textbf { 4 }\par}}
{%beginitem
1. Вкажіть основну \textbf{синтаксичну} відмінність між конструктором копії та конструктором переміщення.%enditem
\par
\vspace{2mm}
%beginitem
2. Максимально змістовно повно сформулюйте назву оголошеного в наведеному фрагменті методу класу.
\vspace{-3mm}
\begin{verbatim}class A{
   ...
   A(const A&);
};\end{verbatim}
\vspace{-3mm}
%enditem
\par
\vspace{2mm}
%beginitem
3. Вкажіть, що не так з наведеним кодом.
\vspace{-3mm}
\begin{verbatim}int a, *p=&a; delete a;\end{verbatim}
\vspace{-3mm}
%enditem
\par
\vspace{2mm}
%beginitem
4. За наявності означень
\vspace{-3mm}
\begin{verbatim}
class A{ };
class B:public A{ };
class C:public B{ };\end{verbatim}
\vspace{-3mm}
вкажіть послідовність викликів конструкторів за виконання коду \code{\{C c;\}} 

(Якщо конструктор викликається кілька разів поспіль, то має зазначатися кожен його виклик, наприклад, A, A, A.)
%enditem

\vspace{2mm}
%beginitem
5. Яку часову оцінку складності має алгоритм сортування злиттям?
%enditem
\par
\vspace{2mm}
%beginitem
6. Модифікувати клас 

\qquad\code{class Session\{ int DM, MA, Alg, Prog;\};} 


додавши до нього оператор перетворення на тип \code{int}: якщо хоча б один атрибут менший за 60, то має повертатись 0, інакше -- сума атрибутів.%enditem
\par
\vspace{2mm}
%beginitem
7. 
Реалізувати операцію симетричної різниці множин з повтореннями. Множина з повтореннями задається переліком своїх елементів, записаних за неспаданням у текстовому файлі через білі символи. Результат у такому ж самому форматі (і теж за неспаданням) має бути записаний у текстовий файл. Питома функція має отримувати три шляхи до файлів. Розмір вхідних множин такий, що вони не можуть бути завантажені повністю в пам’ять програми.

Крім правильності та відповідності умові оцінюється якість коду, наявність витоків ресурсів та коректність роботи з пам'яттю. 

%enditem
}
{\smallskip\smallskip \begin {scriptsize} Затверджено на засіданні кафедри теоретичної кібернетики, протокол \No 4  від   25.01.2024 р.\par\smallskip\smallskip\smallskip\smallskip
{\parbox {6.5 cm} {Зав. кафедри \dotfill Ю.В.~Крак}}\hspace {0.7cm} {\hbox to 6.5 cm { Екзаменатор \dotfill Т.О.~Карнаух}} \end{scriptsize}}
%end
\newpage
%begin
{{\begin {scriptsize}\par
{ \center  Київський національний університет імені Тараса Шевченка \par}
{\parbox {5 cm}{Освітня програма \hfill}{\parbox {7 cm}{Інформатика\hfill}}\parbox {6 cm}{\hfill Семестр 2}}\par
{\parbox {5 cm} {Навчальний предмет \hfill}\parbox {7 cm}{Програмування \hfill}\parbox {6cm}{\hfill}\par}\vspace{-15pt} \end{scriptsize}}{\center Екзаменаційний білет \No \textbf { 5 }\par}}
{%beginitem
1. Вкажіть основну відмінність між класами (class) та структурами (struct).%enditem
\par
\vspace{2mm}
%beginitem
2. Максимально змістовно повно сформулюйте назву оголошеного в наведеному фрагменті методу класу.
\vspace{-3mm}
\begin{verbatim}class A{
   ...
   A();
};\end{verbatim}
\vspace{-3mm}
%enditem
\par
\vspace{2mm}
%beginitem
3. Вкажіть, що не так з наведеним кодом.
\vspace{-3mm}
\begin{verbatim}int *pa=new int(2), *pb=pa; delete pb; *pa=3;\end{verbatim}
\vspace{-3mm}
%enditem
\par
\vspace{2mm}
%beginitem
4. За наявності означень
\vspace{-3mm}
\begin{verbatim}
class A{ };
class B:public A{ };
class C:public B{ };\end{verbatim}
\vspace{-3mm}
вкажіть послідовність викликів деструкторів за виконання коду \code{\{C c;\}} 

(Якщо деструктор викликається кілька разів поспіль, то має зазначатися кожен його виклик, наприклад, $\sim$A, $\sim$A, $\sim$A.)
%enditem

\vspace{2mm}
%beginitem
5. Яку часову оцінку складності має алгоритм швидкого сортування?
%enditem
\par
\vspace{2mm}
%beginitem
6. Модифікувати клас точок простору 

\qquad\code{class Point\{ double x, y, z;\};}


додавши до нього оператор перетворення на рядок типу \code{std::string}, що містить рядкове зображення точки простору (у круглих дужках координати через кому), наприклад, (1.01, 2.35, 0.5), та його означення.%enditem
\par
\vspace{2mm}
%beginitem
7. 
Реалізувати операцію злиття множин з повтореннями. Множина з повтореннями задається переліком своїх елементів, записаних за неспаданням у текстовому файлі через білі символи. Результат у такому ж самому форматі (і теж за неспаданням) має бути записаний у текстовий файл. Питома функція має отримувати три шляхи до файлів. Розмір вхідних множин такий, що вони не можуть бути завантажені повністю в пам’ять програми.

Крім правильності та відповідності умові оцінюється якість коду, наявність витоків ресурсів та коректність роботи з пам'яттю. 

%enditem
}
{\smallskip\smallskip \begin {scriptsize} Затверджено на засіданні кафедри теоретичної кібернетики, протокол \No 4  від   25.01.2024 р.\par\smallskip\smallskip\smallskip\smallskip
{\parbox {6.5 cm} {Зав. кафедри \dotfill Ю.В.~Крак}}\hspace {0.7cm} {\hbox to 6.5 cm { Екзаменатор \dotfill Т.О.~Карнаух}} \end{scriptsize}}
%end
\newpage
%begin
{{\begin {scriptsize}\par
{ \center  Київський національний університет імені Тараса Шевченка \par}
{\parbox {5 cm}{Освітня програма \hfill}{\parbox {7 cm}{Інформатика\hfill}}\parbox {6 cm}{\hfill Семестр 2}}\par
{\parbox {5 cm} {Навчальний предмет \hfill}\parbox {7 cm}{Програмування \hfill}\parbox {6cm}{\hfill}\par}\vspace{-15pt} \end{scriptsize}}{\center Екзаменаційний білет \No \textbf { 6 }\par}}
{%beginitem
1. Вкажіть основну відмінність між віртуальними методами та тими, що не є віртуальними.%enditem
\par
\vspace{2mm}
%beginitem
2. Максимально змістовно повно сформулюйте назву оголошеного в наведеному фрагменті методу класу.
\vspace{-3mm}
\begin{verbatim}class A{
   ...
   A();
};\end{verbatim}
\vspace{-3mm}
%enditem
\par
\vspace{2mm}
%beginitem
3. Вкажіть, що не так з наведеним кодом.
\vspace{-3mm}
\begin{verbatim}int *pa=new int(2), *pb=new int(3); pa=pb;\end{verbatim}
\vspace{-3mm}
%enditem
\par
\vspace{2mm}
%beginitem
4. За наявності означень
\vspace{-3mm}
\begin{verbatim}
class A{ };
class B{ };
class C:public A{B b[2];};\end{verbatim}
\vspace{-3mm}
вкажіть послідовність викликів конструкторів за виконання коду \code{\{C c;\}} 

(Якщо конструктор викликається кілька разів поспіль, то має зазначатися кожен його виклик, наприклад, A, A, A.)
%enditem

\vspace{2mm}
%beginitem
5. Яку часову оцінку складності має алгоритм сортування бульбашкою?
%enditem
\par
\vspace{2mm}
%beginitem
6. Модифікувати клас 

\qquad\code{class Session\{ int DM, MA, Alg, Prog;\};} 


додавши до нього оператор перетворення на тип \code{bool}: об'єкт класу є істинним тоді і тільки тоді, коли значення всіх атрибутів не менші за 60. Означення оператора теж має бути записане.%enditem
\par
\vspace{2mm}
%beginitem
7. 
Задано клас A (структуру, оператори < та == десь коректно означено) та клас вузла списку Node.

struct A\{

\quad  якісь коректні означення 

\quad  bool operator<(const A\&); 

\quad  bool operator==(const A\&); 

\quad A\& operator =(const A\&)=delete; 

\quad  A\& operator =(A\&\&)=delete;\};

struct Node\{Node *prev; Node *next; A a;\};

Написати функцію, яка вилучає зі списку всі вузли, елемент в яких більший елемента в попередньому вузлі списку або дорівнює елементу в наступному вузлі. Список є циклічним та задається вказівником на перший вузол списку. Наприклад, зі списку, що містить послідовність 8, 10, 9, 4, мають бути вилучені вузли, що містять 8 (бо список циклічний і $8>4$) та 10. У цьому прикладі 9 залишається в списку.

Крім правильності та відповідності умові оцінюється якість коду, наявність витоків ресурсів та коректність роботи з пам'яттю. 

%enditem
}
{\smallskip\smallskip \begin {scriptsize} Затверджено на засіданні кафедри теоретичної кібернетики, протокол \No 4  від   25.01.2024 р.\par\smallskip\smallskip\smallskip\smallskip
{\parbox {6.5 cm} {Зав. кафедри \dotfill Ю.В.~Крак}}\hspace {0.7cm} {\hbox to 6.5 cm { Екзаменатор \dotfill Т.О.~Карнаух}} \end{scriptsize}}
%end
\newpage
%begin
{{\begin {scriptsize}\par
{ \center  Київський національний університет імені Тараса Шевченка \par}
{\parbox {5 cm}{Освітня програма \hfill}{\parbox {7 cm}{Інформатика\hfill}}\parbox {6 cm}{\hfill Семестр 2}}\par
{\parbox {5 cm} {Навчальний предмет \hfill}\parbox {7 cm}{Програмування \hfill}\parbox {6cm}{\hfill}\par}\vspace{-15pt} \end{scriptsize}}{\center Екзаменаційний білет \No \textbf { 7 }\par}}
{%beginitem
1. Вкажіть основну відмінність між відкритими (public) та прихованими (private) методами класу.%enditem
\par
\vspace{2mm}
%beginitem
2. Максимально змістовно повно сформулюйте назву оголошеного в наведеному фрагменті методу класу.
\vspace{-3mm}
\begin{verbatim}class A{
   ...
   A& operator =(A&&);
};\end{verbatim}
\vspace{-3mm}
%enditem
\par
\vspace{2mm}
%beginitem
3. Вкажіть, що не так з наведеним кодом.
\vspace{-3mm}
\begin{verbatim}int *pa=new int(2), *pb=&pa; delete pa;\end{verbatim}
\vspace{-3mm}
%enditem
\par
\vspace{2mm}
%beginitem
4. За наявності означень
\vspace{-3mm}
\begin{verbatim}
class A{ };
class B{ };
class C:public A{B b;A a;};\end{verbatim}
\vspace{-3mm}
вкажіть послідовність викликів деструкторів за виконання коду \code{\{C c;\}} 

(Якщо деструктор викликається кілька разів поспіль, то має зазначатися кожен його виклик, наприклад, $\sim$A, $\sim$A, $\sim$A.)
%enditem

\vspace{2mm}
%beginitem
5. Яку часову оцінку складності має алгоритм сортування вибором?
%enditem
\par
\vspace{2mm}
%beginitem
6. Модифікувати клас точок простору 

\qquad\code{class Point\{ double x, y, z;\};}


додавши до нього оператор перетворення на рядок типу \code{std::string}, що містить рядкове зображення точки простору (у круглих дужках координати через кому), наприклад, (1.01, 2.35, 0.5), та його означення.%enditem
\par
\vspace{2mm}
%beginitem
7. 
Задано клас A (структуру, оператори < та == десь коректно означено) та клас вузла списку Node.

struct A\{

\quad  якісь коректні означення 

\quad  bool operator<(const A\&); 

\quad  bool operator==(const A\&); 

\quad A\& operator =(const A\&)=delete; 

\quad  A\& operator =(A\&\&)=delete;\};

struct Node\{Node *next; A a;\};

 Написати функцію, яка вилучає зі списку всі вузли, елемент в яких більший елемента в попередньому вузлі списку або дорівнює елементу в наступному вузлі. Список є циклічним та задається вказівником на перший вузол списку. Наприклад, зі списку, що містить послідовність 8, 10, 9, 4, мають бути вилучені вузли, що містять 8 (бо список циклічний і $8>4$) та 10. У цьому прикладі 9 залишається в списку.

Крім правильності та відповідності умові оцінюється якість коду, наявність витоків ресурсів та коректність роботи з пам'яттю. 

%enditem
}
{\smallskip\smallskip \begin {scriptsize} Затверджено на засіданні кафедри теоретичної кібернетики, протокол \No 4  від   25.01.2024 р.\par\smallskip\smallskip\smallskip\smallskip
{\parbox {6.5 cm} {Зав. кафедри \dotfill Ю.В.~Крак}}\hspace {0.7cm} {\hbox to 6.5 cm { Екзаменатор \dotfill Т.О.~Карнаух}} \end{scriptsize}}
%end
\newpage
%begin
{{\begin {scriptsize}\par
{ \center  Київський національний університет імені Тараса Шевченка \par}
{\parbox {5 cm}{Освітня програма \hfill}{\parbox {7 cm}{Інформатика\hfill}}\parbox {6 cm}{\hfill Семестр 2}}\par
{\parbox {5 cm} {Навчальний предмет \hfill}\parbox {7 cm}{Програмування \hfill}\parbox {6cm}{\hfill}\par}\vspace{-15pt} \end{scriptsize}}{\center Екзаменаційний білет \No \textbf { 8 }\par}}
{%beginitem
1. Вкажіть основну \textbf{синтаксичну} відмінність між конструктором копії та конструктором переміщення.%enditem
\par
\vspace{2mm}
%beginitem
2. Максимально змістовно повно сформулюйте назву оголошеного в наведеному фрагменті методу класу.
\vspace{-3mm}
\begin{verbatim}class A{
   ...
   A(const A&);
};\end{verbatim}
\vspace{-3mm}
%enditem
\par
\vspace{2mm}
%beginitem
3. Вкажіть, що не так з наведеним кодом.
\vspace{-3mm}
\begin{verbatim}int a, *p=&a; delete a;\end{verbatim}
\vspace{-3mm}
%enditem
\par
\vspace{2mm}
%beginitem
4. За наявності означень
\vspace{-3mm}
\begin{verbatim}
class A{ };
class B{ };
class C:public A{B b[2];};\end{verbatim}
\vspace{-3mm}
вкажіть послідовність викликів деструкторів за виконання коду \code{\{C c;\}} 

(Якщо деструктор викликається кілька разів поспіль, то має зазначатися кожен його виклик, наприклад, $\sim$A, $\sim$A, $\sim$A.)
%enditem

\vspace{2mm}
%beginitem
5. Яку часову оцінку складності має алгоритм сортування вставками?
%enditem
\par
\vspace{2mm}
%beginitem
6. Модифікувати клас 

\qquad\code{class Session\{ int DM, MA, Alg, Prog;\};} 


додавши до нього оператор перетворення на тип \code{int}: якщо хоча б один атрибут менший за 60, то має повертатись 0, інакше -- сума атрибутів.%enditem
\par
\vspace{2mm}
%beginitem
7. 
Реалізувати операцію перетину множин з повтореннями. Множина з повтореннями задається переліком своїх елементів, записаних за неспаданням у текстовому файлі через білі символи. Результат у такому ж самому форматі (і теж за неспаданням) має бути записаний у текстовий файл. Питома функція має отримувати три шляхи до файлів. Розмір вхідних множин такий, що вони не можуть бути завантажені повністю в пам’ять програми.

Крім правильності та відповідності умові оцінюється якість коду, наявність витоків ресурсів та коректність роботи з пам'яттю. 

%enditem
}
{\smallskip\smallskip \begin {scriptsize} Затверджено на засіданні кафедри теоретичної кібернетики, протокол \No 4  від   25.01.2024 р.\par\smallskip\smallskip\smallskip\smallskip
{\parbox {6.5 cm} {Зав. кафедри \dotfill Ю.В.~Крак}}\hspace {0.7cm} {\hbox to 6.5 cm { Екзаменатор \dotfill Т.О.~Карнаух}} \end{scriptsize}}
%end
\newpage
%begin
{{\begin {scriptsize}\par
{ \center  Київський національний університет імені Тараса Шевченка \par}
{\parbox {5 cm}{Освітня програма \hfill}{\parbox {7 cm}{Інформатика\hfill}}\parbox {6 cm}{\hfill Семестр 2}}\par
{\parbox {5 cm} {Навчальний предмет \hfill}\parbox {7 cm}{Програмування \hfill}\parbox {6cm}{\hfill}\par}\vspace{-15pt} \end{scriptsize}}{\center Екзаменаційний білет \No \textbf { 9 }\par}}
{%beginitem
1. Вкажіть основну відмінність між віртуальними методами та тими, що не є віртуальними.%enditem
\par
\vspace{2mm}
%beginitem
2. Максимально змістовно повно сформулюйте назву оголошеного в наведеному фрагменті методу класу.
\vspace{-3mm}
\begin{verbatim}class A{
   ...
   ~A();
};\end{verbatim}
\vspace{-3mm}
%enditem
\par
\vspace{2mm}
%beginitem
3. Вкажіть, що не так з наведеним кодом.
\vspace{-3mm}
\begin{verbatim}int **ppi=new int*(new int(3)); delete ppi;\end{verbatim}
\vspace{-3mm}
%enditem
\par
\vspace{2mm}
%beginitem
4. За наявності означень
\vspace{-3mm}
\begin{verbatim}
class A{ };
class B{ };
class C:public A{B b;};\end{verbatim}
\vspace{-3mm}
вкажіть послідовність викликів конструкторів за виконання коду \code{\{C c;\}} 

(Якщо конструктор викликається кілька разів поспіль, то має зазначатися кожен його виклик, наприклад, A, A, A.)
%enditem

\vspace{2mm}
%beginitem
5. Яку часову оцінку складності має алгоритм швидкого сортування?
%enditem
\par
\vspace{2mm}
%beginitem
6. Модифікувати клас 

\qquad\code{class Session\{ int DM, MA, Alg, Prog;\};} 


додавши до нього оператор порівняння \code{<}: один об'єкт є меншим за інший тоді і тільки тоді, коли значення всіх атрибутів об'єктів не менші за 60 і сума атрибутів першого менша за суму атрибутів другого. Означення оператора теж має бути записане.%enditem
\par
\vspace{2mm}
%beginitem
7. 
Задано клас A (структуру, оператори > та == десь коректно означено) та клас вузла списку Node.

struct A\{

\quad  якісь коректні означення 

\quad  bool operator>(const A\&); 

\quad  bool operator==(const A\&); 

\quad A\& operator =(const A\&)=delete; 

\quad  A\& operator =(A\&\&)=delete;\};

struct Node\{Node *next; A a;\};

 Написати функцію, яка вилучає зі списку всі вузли, елемент в яких менший елемента в попередньому вузлі списку або дорівнює елементу в наступному вузлі. Список задається вказівником на перший вузол списку. Наприклад, зі списку, що містить послідовність 8, 5, 6, 9, 4, мають бути вилучені вузли, що містять 5 та 4. У цьому прикладі 6 залишається в списку.

Крім правильності та відповідності умові оцінюється якість коду, наявність витоків ресурсів та коректність роботи з пам'яттю. 

%enditem
}
{\smallskip\smallskip \begin {scriptsize} Затверджено на засіданні кафедри теоретичної кібернетики, протокол \No 4  від   25.01.2024 р.\par\smallskip\smallskip\smallskip\smallskip
{\parbox {6.5 cm} {Зав. кафедри \dotfill Ю.В.~Крак}}\hspace {0.7cm} {\hbox to 6.5 cm { Екзаменатор \dotfill Т.О.~Карнаух}} \end{scriptsize}}
%end
\newpage
%begin
{{\begin {scriptsize}\par
{ \center  Київський національний університет імені Тараса Шевченка \par}
{\parbox {5 cm}{Освітня програма \hfill}{\parbox {7 cm}{Інформатика\hfill}}\parbox {6 cm}{\hfill Семестр 2}}\par
{\parbox {5 cm} {Навчальний предмет \hfill}\parbox {7 cm}{Програмування \hfill}\parbox {6cm}{\hfill}\par}\vspace{-15pt} \end{scriptsize}}{\center Екзаменаційний білет \No \textbf { 10 }\par}}
{%beginitem
1. Вкажіть основну \textbf{семантичну} відмінність між конструктором копії та конструктором переміщення.%enditem
\par
\vspace{2mm}
%beginitem
2. Максимально змістовно повно сформулюйте назву оголошеного в наведеному фрагменті методу класу.
\vspace{-3mm}
\begin{verbatim}class A{
   ...
   A(A&&);
};\end{verbatim}
\vspace{-3mm}
%enditem
\par
\vspace{2mm}
%beginitem
3. Вкажіть, що не так з наведеним кодом.
\vspace{-3mm}
\begin{verbatim}int *pa=new int(2), *pb=new int(3); pa=pb;\end{verbatim}
\vspace{-3mm}
%enditem
\par
\vspace{2mm}
%beginitem
4. За наявності означень
\vspace{-3mm}
\begin{verbatim}
class A{ };
class B{ };
class C{B b;A a;};\end{verbatim}
\vspace{-3mm}
вкажіть послідовність викликів деструкторів за виконання коду \code{\{C c;\}} 

(Якщо деструктор викликається кілька разів поспіль, то має зазначатися кожен його виклик, наприклад, $\sim$A, $\sim$A, $\sim$A.)
%enditem

\vspace{2mm}
%beginitem
5. Яку часову оцінку складності має алгоритм сортування вибором?
%enditem
\par
\vspace{2mm}
%beginitem
6. Модифікувати клас точок простору 

\qquad\code{class Point\{ double x, y, z;\};}


додавши до нього оператор перетворення на тип \code{bool}: точка є істинною тоді і тільки тоді, коли містить хоча б одну ненульову координату. Означення оператора теж має бути записане.%enditem
\par
\vspace{2mm}
%beginitem
7. 
Задано клас A (структуру, оператори < та == десь коректно означено) та клас вузла списку Node.

struct A\{

\quad  якісь коректні означення 

\quad  bool operator<(const A\&); 

\quad  bool operator==(const A\&); 

\quad A\& operator =(const A\&)=delete; 

\quad  A\& operator =(A\&\&)=delete;\};

struct Node\{Node *next; A a;\};

 Написати функцію, яка вилучає зі списку всі вузли, елемент в яких більший елемента в попередньому вузлі списку або дорівнює елементу в наступному вузлі. Список є циклічним та задається вказівником на перший вузол списку. Наприклад, зі списку, що містить послідовність 8, 10, 9, 4, мають бути вилучені вузли, що містять 8 (бо список циклічний і $8>4$) та 10. У цьому прикладі 9 залишається в списку.

Крім правильності та відповідності умові оцінюється якість коду, наявність витоків ресурсів та коректність роботи з пам'яттю. 

%enditem
}
{\smallskip\smallskip \begin {scriptsize} Затверджено на засіданні кафедри теоретичної кібернетики, протокол \No 4  від   25.01.2024 р.\par\smallskip\smallskip\smallskip\smallskip
{\parbox {6.5 cm} {Зав. кафедри \dotfill Ю.В.~Крак}}\hspace {0.7cm} {\hbox to 6.5 cm { Екзаменатор \dotfill Т.О.~Карнаух}} \end{scriptsize}}
%end
\newpage
%begin
{{\begin {scriptsize}\par
{ \center  Київський національний університет імені Тараса Шевченка \par}
{\parbox {5 cm}{Освітня програма \hfill}{\parbox {7 cm}{Інформатика\hfill}}\parbox {6 cm}{\hfill Семестр 2}}\par
{\parbox {5 cm} {Навчальний предмет \hfill}\parbox {7 cm}{Програмування \hfill}\parbox {6cm}{\hfill}\par}\vspace{-15pt} \end{scriptsize}}{\center Екзаменаційний білет \No \textbf { 11 }\par}}
{%beginitem
1. Що важливо забезпечити при роботі з динамічними змінними?%enditem
\par
\vspace{2mm}
%beginitem
2. Максимально змістовно повно сформулюйте назву оголошеного в наведеному фрагменті методу класу.
\vspace{-3mm}
\begin{verbatim}class A{
   ...
   A(const A&);
};\end{verbatim}
\vspace{-3mm}
%enditem
\par
\vspace{2mm}
%beginitem
3. Вкажіть, що не так з наведеним кодом.
\vspace{-3mm}
\begin{verbatim}int *pa=new int(2), *pb=pa; delete pb; delete pa;\end{verbatim}
\vspace{-3mm}
%enditem
\par
\vspace{2mm}
%beginitem
4. За наявності означень
\vspace{-3mm}
\begin{verbatim}
class A{ };
class B{ };
class C:public A{B b;A a;};\end{verbatim}
\vspace{-3mm}
вкажіть послідовність викликів конструкторів за виконання коду \code{\{C c;\}} 

(Якщо конструктор викликається кілька разів поспіль, то має зазначатися кожен його виклик, наприклад, A, A, A.)
%enditem

\vspace{2mm}
%beginitem
5. Яку часову оцінку складності має алгоритм швидкого сортування?
%enditem
\par
\vspace{2mm}
%beginitem
6. Модифікувати клас точок простору 

\qquad\code{class Point\{ double x, y, z;\};}


додавши до нього оператор перетворення на тип \code{bool}: точка є істинною тоді і тільки тоді, коли містить хоча б одну ненульову координату. Означення оператора теж має бути записане.%enditem
\par
\vspace{2mm}
%beginitem
7. 
Реалізувати операцію симетричної різниці множин з повтореннями. Множина з повтореннями задається переліком своїх елементів, записаних за неспаданням у текстовому файлі через білі символи. Результат у такому ж самому форматі (і теж за неспаданням) має бути записаний у текстовий файл. Питома функція має отримувати три шляхи до файлів. Розмір вхідних множин такий, що вони не можуть бути завантажені повністю в пам’ять програми.

Крім правильності та відповідності умові оцінюється якість коду, наявність витоків ресурсів та коректність роботи з пам'яттю. 

%enditem
}
{\smallskip\smallskip \begin {scriptsize} Затверджено на засіданні кафедри теоретичної кібернетики, протокол \No 4  від   25.01.2024 р.\par\smallskip\smallskip\smallskip\smallskip
{\parbox {6.5 cm} {Зав. кафедри \dotfill Ю.В.~Крак}}\hspace {0.7cm} {\hbox to 6.5 cm { Екзаменатор \dotfill Т.О.~Карнаух}} \end{scriptsize}}
%end
\newpage
%begin
{{\begin {scriptsize}\par
{ \center  Київський національний університет імені Тараса Шевченка \par}
{\parbox {5 cm}{Освітня програма \hfill}{\parbox {7 cm}{Інформатика\hfill}}\parbox {6 cm}{\hfill Семестр 2}}\par
{\parbox {5 cm} {Навчальний предмет \hfill}\parbox {7 cm}{Програмування \hfill}\parbox {6cm}{\hfill}\par}\vspace{-15pt} \end{scriptsize}}{\center Екзаменаційний білет \No \textbf { 12 }\par}}
{%beginitem
1. Вкажіть основну відмінність між класами (class) та структурами (struct).%enditem
\par
\vspace{2mm}
%beginitem
2. Максимально змістовно повно сформулюйте назву оголошеного в наведеному фрагменті методу класу.
\vspace{-3mm}
\begin{verbatim}class A{
   ...
   A(A&&);
};\end{verbatim}
\vspace{-3mm}
%enditem
\par
\vspace{2mm}
%beginitem
3. Вкажіть, що не так з наведеним кодом.
\vspace{-3mm}
\begin{verbatim}int *pa=new int(2), *pb=*pa; delete pa;\end{verbatim}
\vspace{-3mm}
%enditem
\par
\vspace{2mm}
%beginitem
4. За наявності означень
\vspace{-3mm}
\begin{verbatim}
class A{ };
class B{ };
class C:public A{B b;};\end{verbatim}
\vspace{-3mm}
вкажіть послідовність викликів деструкторів за виконання коду \code{\{C c;\}} 

(Якщо деструктор викликається кілька разів поспіль, то має зазначатися кожен його виклик, наприклад, $\sim$A, $\sim$A, $\sim$A.)
%enditem

\vspace{2mm}
%beginitem
5. Яку часову оцінку складності має алгоритм сортування купою (він же деревом, пірамідою)?
%enditem
\par
\vspace{2mm}
%beginitem
6. Модифікувати клас 

\qquad\code{class Session\{ int DM, MA, Alg, Prog;\};} 


додавши до нього оператор перетворення на тип \code{bool}: об'єкт класу є істинним тоді і тільки тоді, коли значення всіх атрибутів не менші за 60. Означення оператора теж має бути записане.%enditem
\par
\vspace{2mm}
%beginitem
7. 
Реалізувати операцію різниці множин з повтореннями. Множина з повтореннями задається переліком своїх елементів, записаних за неспаданням у текстовому файлі через білі символи. Результат у такому ж самому форматі (і теж за неспаданням) має бути записаний у текстовий файл. Питома функція має отримувати три шляхи до файлів. Розмір вхідних множин такий, що вони не можуть бути завантажені повністю в пам’ять програми.

Крім правильності та відповідності умові оцінюється якість коду, наявність витоків ресурсів та коректність роботи з пам'яттю. 

%enditem
}
{\smallskip\smallskip \begin {scriptsize} Затверджено на засіданні кафедри теоретичної кібернетики, протокол \No 4  від   25.01.2024 р.\par\smallskip\smallskip\smallskip\smallskip
{\parbox {6.5 cm} {Зав. кафедри \dotfill Ю.В.~Крак}}\hspace {0.7cm} {\hbox to 6.5 cm { Екзаменатор \dotfill Т.О.~Карнаух}} \end{scriptsize}}
%end
\newpage
%begin
{{\begin {scriptsize}\par
{ \center  Київський національний університет імені Тараса Шевченка \par}
{\parbox {5 cm}{Освітня програма \hfill}{\parbox {7 cm}{Інформатика\hfill}}\parbox {6 cm}{\hfill Семестр 2}}\par
{\parbox {5 cm} {Навчальний предмет \hfill}\parbox {7 cm}{Програмування \hfill}\parbox {6cm}{\hfill}\par}\vspace{-15pt} \end{scriptsize}}{\center Екзаменаційний білет \No \textbf { 13 }\par}}
{%beginitem
1. Вкажіть основну \textbf{синтаксичну} відмінність між конструктором копії та конструктором переміщення.%enditem
\par
\vspace{2mm}
%beginitem
2. Максимально змістовно повно сформулюйте назву оголошеного в наведеному фрагменті методу класу.
\vspace{-3mm}
\begin{verbatim}class A{
   ...
   ~A();
};\end{verbatim}
\vspace{-3mm}
%enditem
\par
\vspace{2mm}
%beginitem
3. Вкажіть, що не так з наведеним кодом.
\vspace{-3mm}
\begin{verbatim}int *pa=new int(2), *pb=pa; delete pb; *pa=3;\end{verbatim}
\vspace{-3mm}
%enditem
\par
\vspace{2mm}
%beginitem
4. За наявності означень
\vspace{-3mm}
\begin{verbatim}
class A{ };
class B:public A{ };
class C:public B{ };\end{verbatim}
\vspace{-3mm}
вкажіть послідовність викликів деструкторів за виконання коду \code{\{C c;\}} 

(Якщо деструктор викликається кілька разів поспіль, то має зазначатися кожен його виклик, наприклад, $\sim$A, $\sim$A, $\sim$A.)
%enditem

\vspace{2mm}
%beginitem
5. Яку часову оцінку складності має алгоритм сортування бульбашкою?
%enditem
\par
\vspace{2mm}
%beginitem
6. Модифікувати клас точок простору 

\qquad\code{class Point\{ double x, y, z;\};}


додавши до нього оператор перетворення на рядок типу \code{std::string}, що містить рядкове зображення точки простору (у круглих дужках координати через кому), наприклад, (1.01, 2.35, 0.5), та його означення.%enditem
\par
\vspace{2mm}
%beginitem
7. 
Задано клас A (структуру, оператори < та == десь коректно означено) та клас вузла списку Node.

struct A\{

\quad  якісь коректні означення 

\quad  bool operator<(const A\&); 

\quad A\& operator =(const A\&)=delete; 

\quad  A\& operator =(A\&\&)=delete;\};

struct Node\{Node *prev; Node *next; A a;\};

Написати функцію, яка вилучає зі списку всі вузли, елемент в яких менший елемента в попередньому вузлі списку або дорівнює елементу в наступному вузлі. Список задається вказівником на перший вузол списку. Наприклад, зі списку, що містить послідовність 8, 5, 6, 9, 4, мають бути вилучені вузли, що містять 5 та 4. У цьому прикладі 6 залишається в списку.

Крім правильності та відповідності умові оцінюється якість коду, наявність витоків ресурсів та коректність роботи з пам'яттю. 

%enditem
}
{\smallskip\smallskip \begin {scriptsize} Затверджено на засіданні кафедри теоретичної кібернетики, протокол \No 4  від   25.01.2024 р.\par\smallskip\smallskip\smallskip\smallskip
{\parbox {6.5 cm} {Зав. кафедри \dotfill Ю.В.~Крак}}\hspace {0.7cm} {\hbox to 6.5 cm { Екзаменатор \dotfill Т.О.~Карнаух}} \end{scriptsize}}
%end
\newpage
%begin
{{\begin {scriptsize}\par
{ \center  Київський національний університет імені Тараса Шевченка \par}
{\parbox {5 cm}{Освітня програма \hfill}{\parbox {7 cm}{Інформатика\hfill}}\parbox {6 cm}{\hfill Семестр 2}}\par
{\parbox {5 cm} {Навчальний предмет \hfill}\parbox {7 cm}{Програмування \hfill}\parbox {6cm}{\hfill}\par}\vspace{-15pt} \end{scriptsize}}{\center Екзаменаційний білет \No \textbf { 14 }\par}}
{%beginitem
1. Що важливо забезпечити при роботі з динамічними змінними?%enditem
\par
\vspace{2mm}
%beginitem
2. Максимально змістовно повно сформулюйте назву оголошеного в наведеному фрагменті методу класу.
\vspace{-3mm}
\begin{verbatim}class A{
   ...
   A& operator =(A&&);
};\end{verbatim}
\vspace{-3mm}
%enditem
\par
\vspace{2mm}
%beginitem
3. Вкажіть, що не так з наведеним кодом.
\vspace{-3mm}
\begin{verbatim}int *pa=new int(2), *pb=&pa; delete pa;\end{verbatim}
\vspace{-3mm}
%enditem
\par
\vspace{2mm}
%beginitem
4. За наявності означень
\vspace{-3mm}
\begin{verbatim}
class A{ };
class B{ };
class C:public A{B b;};\end{verbatim}
\vspace{-3mm}
вкажіть послідовність викликів конструкторів за виконання коду \code{\{C c;\}} 

(Якщо конструктор викликається кілька разів поспіль, то має зазначатися кожен його виклик, наприклад, A, A, A.)
%enditem

\vspace{2mm}
%beginitem
5. Яку часову оцінку складності має алгоритм сортування злиттям?
%enditem
\par
\vspace{2mm}
%beginitem
6. Модифікувати клас 

\qquad\code{class Session\{ int DM, MA, Alg, Prog;\};} 


додавши до нього оператор перетворення на тип \code{int}: якщо хоча б один атрибут менший за 60, то має повертатись 0, інакше -- сума атрибутів.%enditem
\par
\vspace{2mm}
%beginitem
7. 
Реалізувати операцію перетину множин з повтореннями. Множина з повтореннями задається переліком своїх елементів, записаних за неспаданням у текстовому файлі через білі символи. Результат у такому ж самому форматі (і теж за неспаданням) має бути записаний у текстовий файл. Питома функція має отримувати три шляхи до файлів. Розмір вхідних множин такий, що вони не можуть бути завантажені повністю в пам’ять програми.

Крім правильності та відповідності умові оцінюється якість коду, наявність витоків ресурсів та коректність роботи з пам'яттю. 

%enditem
}
{\smallskip\smallskip \begin {scriptsize} Затверджено на засіданні кафедри теоретичної кібернетики, протокол \No 4  від   25.01.2024 р.\par\smallskip\smallskip\smallskip\smallskip
{\parbox {6.5 cm} {Зав. кафедри \dotfill Ю.В.~Крак}}\hspace {0.7cm} {\hbox to 6.5 cm { Екзаменатор \dotfill Т.О.~Карнаух}} \end{scriptsize}}
%end
\newpage
%begin
{{\begin {scriptsize}\par
{ \center  Київський національний університет імені Тараса Шевченка \par}
{\parbox {5 cm}{Освітня програма \hfill}{\parbox {7 cm}{Інформатика\hfill}}\parbox {6 cm}{\hfill Семестр 2}}\par
{\parbox {5 cm} {Навчальний предмет \hfill}\parbox {7 cm}{Програмування \hfill}\parbox {6cm}{\hfill}\par}\vspace{-15pt} \end{scriptsize}}{\center Екзаменаційний білет \No \textbf { 15 }\par}}
{%beginitem
1. Вкажіть основну відмінність між відкритими (public) та прихованими (private) методами класу.%enditem
\par
\vspace{2mm}
%beginitem
2. Максимально змістовно повно сформулюйте назву оголошеного в наведеному фрагменті методу класу.
\vspace{-3mm}
\begin{verbatim}class A{
   ...
   A();
};\end{verbatim}
\vspace{-3mm}
%enditem
\par
\vspace{2mm}
%beginitem
3. Вкажіть, що не так з наведеним кодом.
\vspace{-3mm}
\begin{verbatim}int *pa=new int(2), *pb=*pa; delete pa;\end{verbatim}
\vspace{-3mm}
%enditem
\par
\vspace{2mm}
%beginitem
4. За наявності означень
\vspace{-3mm}
\begin{verbatim}
class A{ };
class B{ };
class C:public A{B b;A a;};\end{verbatim}
\vspace{-3mm}
вкажіть послідовність викликів деструкторів за виконання коду \code{\{C c;\}} 

(Якщо деструктор викликається кілька разів поспіль, то має зазначатися кожен його виклик, наприклад, $\sim$A, $\sim$A, $\sim$A.)
%enditem

\vspace{2mm}
%beginitem
5. Яку часову оцінку складності має алгоритм сортування злиттям?
%enditem
\par
\vspace{2mm}
%beginitem
6. Модифікувати клас 

\qquad\code{class Session\{ int DM, MA, Alg, Prog;\};} 


додавши до нього оператор порівняння \code{<}: один об'єкт є меншим за інший тоді і тільки тоді, коли значення всіх атрибутів об'єктів не менші за 60 і сума атрибутів першого менша за суму атрибутів другого. Означення оператора теж має бути записане.%enditem
\par
\vspace{2mm}
%beginitem
7. 
Задано клас A (структуру, оператори < та == десь коректно означено) та клас вузла списку Node.

struct A\{

\quad  якісь коректні означення 

\quad  bool operator<(const A\&); 

\quad  bool operator==(const A\&); 

\quad A\& operator =(const A\&)=delete; 

\quad  A\& operator =(A\&\&)=delete;\};

struct Node\{Node *prev; Node *next; A a;\};

Написати функцію, яка вилучає зі списку всі вузли, елемент в яких більший елемента в попередньому вузлі списку або дорівнює елементу в наступному вузлі. Список є циклічним та задається вказівником на перший вузол списку. Наприклад, зі списку, що містить послідовність 8, 10, 9, 4, мають бути вилучені вузли, що містять 8 (бо список циклічний і $8>4$) та 10. У цьому прикладі 9 залишається в списку.

Крім правильності та відповідності умові оцінюється якість коду, наявність витоків ресурсів та коректність роботи з пам'яттю. 

%enditem
}
{\smallskip\smallskip \begin {scriptsize} Затверджено на засіданні кафедри теоретичної кібернетики, протокол \No 4  від   25.01.2024 р.\par\smallskip\smallskip\smallskip\smallskip
{\parbox {6.5 cm} {Зав. кафедри \dotfill Ю.В.~Крак}}\hspace {0.7cm} {\hbox to 6.5 cm { Екзаменатор \dotfill Т.О.~Карнаух}} \end{scriptsize}}
%end
\newpage
%begin
{{\begin {scriptsize}\par
{ \center  Київський національний університет імені Тараса Шевченка \par}
{\parbox {5 cm}{Освітня програма \hfill}{\parbox {7 cm}{Інформатика\hfill}}\parbox {6 cm}{\hfill Семестр 2}}\par
{\parbox {5 cm} {Навчальний предмет \hfill}\parbox {7 cm}{Програмування \hfill}\parbox {6cm}{\hfill}\par}\vspace{-15pt} \end{scriptsize}}{\center Екзаменаційний білет \No \textbf { 16 }\par}}
{%beginitem
1. Вкажіть основну відмінність між класами (class) та структурами (struct).%enditem
\par
\vspace{2mm}
%beginitem
2. Максимально змістовно повно сформулюйте назву оголошеного в наведеному фрагменті методу класу.
\vspace{-3mm}
\begin{verbatim}class A{
   ...
   A& operator =(A&&);
};\end{verbatim}
\vspace{-3mm}
%enditem
\par
\vspace{2mm}
%beginitem
3. Вкажіть, що не так з наведеним кодом.
\vspace{-3mm}
\begin{verbatim}int a, *p=&a; delete a;\end{verbatim}
\vspace{-3mm}
%enditem
\par
\vspace{2mm}
%beginitem
4. За наявності означень
\vspace{-3mm}
\begin{verbatim}
class A{ };
class B{ };
class C:public A{B b;};\end{verbatim}
\vspace{-3mm}
вкажіть послідовність викликів деструкторів за виконання коду \code{\{C c;\}} 

(Якщо деструктор викликається кілька разів поспіль, то має зазначатися кожен його виклик, наприклад, $\sim$A, $\sim$A, $\sim$A.)
%enditem

\vspace{2mm}
%beginitem
5. Яку часову оцінку складності має алгоритм сортування вибором?
%enditem
\par
\vspace{2mm}
%beginitem
6. Модифікувати клас точок простору 

\qquad\code{class Point\{ double x, y, z;\};}


додавши до нього оператор перетворення на тип \code{bool}: точка є істинною тоді і тільки тоді, коли містить хоча б одну ненульову координату. Означення оператора теж має бути записане.%enditem
\par
\vspace{2mm}
%beginitem
7. 
Реалізувати операцію злиття множин з повтореннями. Множина з повтореннями задається переліком своїх елементів, записаних за неспаданням у текстовому файлі через білі символи. Результат у такому ж самому форматі (і теж за неспаданням) має бути записаний у текстовий файл. Питома функція має отримувати три шляхи до файлів. Розмір вхідних множин такий, що вони не можуть бути завантажені повністю в пам’ять програми.

Крім правильності та відповідності умові оцінюється якість коду, наявність витоків ресурсів та коректність роботи з пам'яттю. 

%enditem
}
{\smallskip\smallskip \begin {scriptsize} Затверджено на засіданні кафедри теоретичної кібернетики, протокол \No 4  від   25.01.2024 р.\par\smallskip\smallskip\smallskip\smallskip
{\parbox {6.5 cm} {Зав. кафедри \dotfill Ю.В.~Крак}}\hspace {0.7cm} {\hbox to 6.5 cm { Екзаменатор \dotfill Т.О.~Карнаух}} \end{scriptsize}}
%end
\newpage
%begin
{{\begin {scriptsize}\par
{ \center  Київський національний університет імені Тараса Шевченка \par}
{\parbox {5 cm}{Освітня програма \hfill}{\parbox {7 cm}{Інформатика\hfill}}\parbox {6 cm}{\hfill Семестр 2}}\par
{\parbox {5 cm} {Навчальний предмет \hfill}\parbox {7 cm}{Програмування \hfill}\parbox {6cm}{\hfill}\par}\vspace{-15pt} \end{scriptsize}}{\center Екзаменаційний білет \No \textbf { 17 }\par}}
{%beginitem
1. Вкажіть основну \textbf{семантичну} відмінність між конструктором копії та конструктором переміщення.%enditem
\par
\vspace{2mm}
%beginitem
2. Максимально змістовно повно сформулюйте назву оголошеного в наведеному фрагменті методу класу.
\vspace{-3mm}
\begin{verbatim}class A{
   ...
   A(const A&);
};\end{verbatim}
\vspace{-3mm}
%enditem
\par
\vspace{2mm}
%beginitem
3. Вкажіть, що не так з наведеним кодом.
\vspace{-3mm}
\begin{verbatim}int **ppi=new int*(new int(3)); delete ppi;\end{verbatim}
\vspace{-3mm}
%enditem
\par
\vspace{2mm}
%beginitem
4. За наявності означень
\vspace{-3mm}
\begin{verbatim}
class A{ };
class B:public A{ };
class C:public A{B b;};\end{verbatim}
\vspace{-3mm}
вкажіть послідовність викликів конструкторів за виконання коду \code{\{C c;\}} 

(Якщо конструктор викликається кілька разів поспіль, то має зазначатися кожен його виклик, наприклад, A, A, A.)
%enditem

\vspace{2mm}
%beginitem
5. Яку часову оцінку складності має алгоритм сортування купою (він же деревом, пірамідою)?
%enditem
\par
\vspace{2mm}
%beginitem
6. Модифікувати клас 

\qquad\code{class Session\{ int DM, MA, Alg, Prog;\};} 


додавши до нього оператор порівняння \code{<}: один об'єкт є меншим за інший тоді і тільки тоді, коли значення всіх атрибутів об'єктів не менші за 60 і сума атрибутів першого менша за суму атрибутів другого. Означення оператора теж має бути записане.%enditem
\par
\vspace{2mm}
%beginitem
7. 
Задано клас A (структуру, оператори < та == десь коректно означено) та клас вузла списку Node.

struct A\{

\quad  якісь коректні означення 

\quad  bool operator<(const A\&); 

\quad  bool operator==(const A\&); 

\quad A\& operator =(const A\&)=delete; 

\quad  A\& operator =(A\&\&)=delete;\};

struct Node\{Node *prev; Node *next; A a;\};

Написати функцію, яка вилучає зі списку всі вузли, елемент в яких більший елемента в попередньому вузлі списку або дорівнює елементу в наступному вузлі. Список є циклічним та задається вказівником на перший вузол списку. Наприклад, зі списку, що містить послідовність 8, 10, 9, 4, мають бути вилучені вузли, що містять 8 (бо список циклічний і $8>4$) та 10. У цьому прикладі 9 залишається в списку.

Крім правильності та відповідності умові оцінюється якість коду, наявність витоків ресурсів та коректність роботи з пам'яттю. 

%enditem
}
{\smallskip\smallskip \begin {scriptsize} Затверджено на засіданні кафедри теоретичної кібернетики, протокол \No 4  від   25.01.2024 р.\par\smallskip\smallskip\smallskip\smallskip
{\parbox {6.5 cm} {Зав. кафедри \dotfill Ю.В.~Крак}}\hspace {0.7cm} {\hbox to 6.5 cm { Екзаменатор \dotfill Т.О.~Карнаух}} \end{scriptsize}}
%end
\newpage
%begin
{{\begin {scriptsize}\par
{ \center  Київський національний університет імені Тараса Шевченка \par}
{\parbox {5 cm}{Освітня програма \hfill}{\parbox {7 cm}{Інформатика\hfill}}\parbox {6 cm}{\hfill Семестр 2}}\par
{\parbox {5 cm} {Навчальний предмет \hfill}\parbox {7 cm}{Програмування \hfill}\parbox {6cm}{\hfill}\par}\vspace{-15pt} \end{scriptsize}}{\center Екзаменаційний білет \No \textbf { 18 }\par}}
{%beginitem
1. Вкажіть основну відмінність між віртуальними методами та тими, що не є віртуальними.%enditem
\par
\vspace{2mm}
%beginitem
2. Максимально змістовно повно сформулюйте назву оголошеного в наведеному фрагменті методу класу.
\vspace{-3mm}
\begin{verbatim}class A{
   ...
   ~A();
};\end{verbatim}
\vspace{-3mm}
%enditem
\par
\vspace{2mm}
%beginitem
3. Вкажіть, що не так з наведеним кодом.
\vspace{-3mm}
\begin{verbatim}int *pa=new int(2), *pb=&pa; delete pa;\end{verbatim}
\vspace{-3mm}
%enditem
\par
\vspace{2mm}
%beginitem
4. За наявності означень
\vspace{-3mm}
\begin{verbatim}
class A{ };
class B{ };
class C:public A{B b[2];};\end{verbatim}
\vspace{-3mm}
вкажіть послідовність викликів конструкторів за виконання коду \code{\{C c;\}} 

(Якщо конструктор викликається кілька разів поспіль, то має зазначатися кожен його виклик, наприклад, A, A, A.)
%enditem

\vspace{2mm}
%beginitem
5. Яку часову оцінку складності має алгоритм швидкого сортування?
%enditem
\par
\vspace{2mm}
%beginitem
6. Модифікувати клас 

\qquad\code{class Session\{ int DM, MA, Alg, Prog;\};} 


додавши до нього оператор перетворення на тип \code{int}: якщо хоча б один атрибут менший за 60, то має повертатись 0, інакше -- сума атрибутів.%enditem
\par
\vspace{2mm}
%beginitem
7. 
Реалізувати операцію симетричної різниці множин з повтореннями. Множина з повтореннями задається переліком своїх елементів, записаних за неспаданням у текстовому файлі через білі символи. Результат у такому ж самому форматі (і теж за неспаданням) має бути записаний у текстовий файл. Питома функція має отримувати три шляхи до файлів. Розмір вхідних множин такий, що вони не можуть бути завантажені повністю в пам’ять програми.

Крім правильності та відповідності умові оцінюється якість коду, наявність витоків ресурсів та коректність роботи з пам'яттю. 

%enditem
}
{\smallskip\smallskip \begin {scriptsize} Затверджено на засіданні кафедри теоретичної кібернетики, протокол \No 4  від   25.01.2024 р.\par\smallskip\smallskip\smallskip\smallskip
{\parbox {6.5 cm} {Зав. кафедри \dotfill Ю.В.~Крак}}\hspace {0.7cm} {\hbox to 6.5 cm { Екзаменатор \dotfill Т.О.~Карнаух}} \end{scriptsize}}
%end
\newpage
%begin
{{\begin {scriptsize}\par
{ \center  Київський національний університет імені Тараса Шевченка \par}
{\parbox {5 cm}{Освітня програма \hfill}{\parbox {7 cm}{Інформатика\hfill}}\parbox {6 cm}{\hfill Семестр 2}}\par
{\parbox {5 cm} {Навчальний предмет \hfill}\parbox {7 cm}{Програмування \hfill}\parbox {6cm}{\hfill}\par}\vspace{-15pt} \end{scriptsize}}{\center Екзаменаційний білет \No \textbf { 19 }\par}}
{%beginitem
1. Що важливо забезпечити при роботі з динамічними змінними?%enditem
\par
\vspace{2mm}
%beginitem
2. Максимально змістовно повно сформулюйте назву оголошеного в наведеному фрагменті методу класу.
\vspace{-3mm}
\begin{verbatim}class A{
   ...
   A(A&&);
};\end{verbatim}
\vspace{-3mm}
%enditem
\par
\vspace{2mm}
%beginitem
3. Вкажіть, що не так з наведеним кодом.
\vspace{-3mm}
\begin{verbatim}int *pa=new int(2), *pb=pa; delete pb; delete pa;\end{verbatim}
\vspace{-3mm}
%enditem
\par
\vspace{2mm}
%beginitem
4. За наявності означень
\vspace{-3mm}
\begin{verbatim}
class A{ };
class B:public A{ };
class C:public A{B b;};\end{verbatim}
\vspace{-3mm}
вкажіть послідовність викликів деструкторів за виконання коду \code{\{C c;\}} 

(Якщо деструктор викликається кілька разів поспіль, то має зазначатися кожен його виклик, наприклад, $\sim$A, $\sim$A, $\sim$A.)
%enditem

\vspace{2mm}
%beginitem
5. Яку часову оцінку складності має алгоритм швидкого сортування?
%enditem
\par
\vspace{2mm}
%beginitem
6. Модифікувати клас точок простору 

\qquad\code{class Point\{ double x, y, z;\};}


додавши до нього оператор перетворення на рядок типу \code{std::string}, що містить рядкове зображення точки простору (у круглих дужках координати через кому), наприклад, (1.01, 2.35, 0.5), та його означення.%enditem
\par
\vspace{2mm}
%beginitem
7. 
Задано клас A (структуру, оператори < та == десь коректно означено) та клас вузла списку Node.

struct A\{

\quad  якісь коректні означення 

\quad  bool operator<(const A\&); 

\quad A\& operator =(const A\&)=delete; 

\quad  A\& operator =(A\&\&)=delete;\};

struct Node\{Node *prev; Node *next; A a;\};

Написати функцію, яка вилучає зі списку всі вузли, елемент в яких менший елемента в попередньому вузлі списку або дорівнює елементу в наступному вузлі. Список задається вказівником на перший вузол списку. Наприклад, зі списку, що містить послідовність 8, 5, 6, 9, 4, мають бути вилучені вузли, що містять 5 та 4. У цьому прикладі 6 залишається в списку.

Крім правильності та відповідності умові оцінюється якість коду, наявність витоків ресурсів та коректність роботи з пам'яттю. 

%enditem
}
{\smallskip\smallskip \begin {scriptsize} Затверджено на засіданні кафедри теоретичної кібернетики, протокол \No 4  від   25.01.2024 р.\par\smallskip\smallskip\smallskip\smallskip
{\parbox {6.5 cm} {Зав. кафедри \dotfill Ю.В.~Крак}}\hspace {0.7cm} {\hbox to 6.5 cm { Екзаменатор \dotfill Т.О.~Карнаух}} \end{scriptsize}}
%end
\newpage
%begin
{{\begin {scriptsize}\par
{ \center  Київський національний університет імені Тараса Шевченка \par}
{\parbox {5 cm}{Освітня програма \hfill}{\parbox {7 cm}{Інформатика\hfill}}\parbox {6 cm}{\hfill Семестр 2}}\par
{\parbox {5 cm} {Навчальний предмет \hfill}\parbox {7 cm}{Програмування \hfill}\parbox {6cm}{\hfill}\par}\vspace{-15pt} \end{scriptsize}}{\center Екзаменаційний білет \No \textbf { 20 }\par}}
{%beginitem
1. Вкажіть основну \textbf{синтаксичну} відмінність між конструктором копії та конструктором переміщення.%enditem
\par
\vspace{2mm}
%beginitem
2. Максимально змістовно повно сформулюйте назву оголошеного в наведеному фрагменті методу класу.
\vspace{-3mm}
\begin{verbatim}class A{
   ...
   A();
};\end{verbatim}
\vspace{-3mm}
%enditem
\par
\vspace{2mm}
%beginitem
3. Вкажіть, що не так з наведеним кодом.
\vspace{-3mm}
\begin{verbatim}int *pa=new int(2), *pb=pa; delete pb; *pa=3;\end{verbatim}
\vspace{-3mm}
%enditem
\par
\vspace{2mm}
%beginitem
4. За наявності означень
\vspace{-3mm}
\begin{verbatim}
class A{ };
class B{ };
class C{B b;A a;};\end{verbatim}
\vspace{-3mm}
вкажіть послідовність викликів деструкторів за виконання коду \code{\{C c;\}} 

(Якщо деструктор викликається кілька разів поспіль, то має зазначатися кожен його виклик, наприклад, $\sim$A, $\sim$A, $\sim$A.)
%enditem

\vspace{2mm}
%beginitem
5. Яку часову оцінку складності має алгоритм сортування бульбашкою?
%enditem
\par
\vspace{2mm}
%beginitem
6. Модифікувати клас 

\qquad\code{class Session\{ int DM, MA, Alg, Prog;\};} 


додавши до нього оператор перетворення на тип \code{bool}: об'єкт класу є істинним тоді і тільки тоді, коли значення всіх атрибутів не менші за 60. Означення оператора теж має бути записане.%enditem
\par
\vspace{2mm}
%beginitem
7. 
Реалізувати операцію злиття множин з повтореннями. Множина з повтореннями задається переліком своїх елементів, записаних за неспаданням у текстовому файлі через білі символи. Результат у такому ж самому форматі (і теж за неспаданням) має бути записаний у текстовий файл. Питома функція має отримувати три шляхи до файлів. Розмір вхідних множин такий, що вони не можуть бути завантажені повністю в пам’ять програми.

Крім правильності та відповідності умові оцінюється якість коду, наявність витоків ресурсів та коректність роботи з пам'яттю. 

%enditem
}
{\smallskip\smallskip \begin {scriptsize} Затверджено на засіданні кафедри теоретичної кібернетики, протокол \No 4  від   25.01.2024 р.\par\smallskip\smallskip\smallskip\smallskip
{\parbox {6.5 cm} {Зав. кафедри \dotfill Ю.В.~Крак}}\hspace {0.7cm} {\hbox to 6.5 cm { Екзаменатор \dotfill Т.О.~Карнаух}} \end{scriptsize}}
%end
\newpage
%begin
{{\begin {scriptsize}\par
{ \center  Київський національний університет імені Тараса Шевченка \par}
{\parbox {5 cm}{Освітня програма \hfill}{\parbox {7 cm}{Інформатика\hfill}}\parbox {6 cm}{\hfill Семестр 2}}\par
{\parbox {5 cm} {Навчальний предмет \hfill}\parbox {7 cm}{Програмування \hfill}\parbox {6cm}{\hfill}\par}\vspace{-15pt} \end{scriptsize}}{\center Екзаменаційний білет \No \textbf { 21 }\par}}
{%beginitem
1. Вкажіть основну \textbf{семантичну} відмінність між конструктором копії та конструктором переміщення.%enditem
\par
\vspace{2mm}
%beginitem
2. Максимально змістовно повно сформулюйте назву оголошеного в наведеному фрагменті методу класу.
\vspace{-3mm}
\begin{verbatim}class A{
   ...
   A& operator =(A&&);
};\end{verbatim}
\vspace{-3mm}
%enditem
\par
\vspace{2mm}
%beginitem
3. Вкажіть, що не так з наведеним кодом.
\vspace{-3mm}
\begin{verbatim}int *pa=new int(2), *pb=new int(3); pa=pb;\end{verbatim}
\vspace{-3mm}
%enditem
\par
\vspace{2mm}
%beginitem
4. За наявності означень
\vspace{-3mm}
\begin{verbatim}
class A{ };
class B:public A{ };
class C:public B{ };\end{verbatim}
\vspace{-3mm}
вкажіть послідовність викликів конструкторів за виконання коду \code{\{C c;\}} 

(Якщо конструктор викликається кілька разів поспіль, то має зазначатися кожен його виклик, наприклад, A, A, A.)
%enditem

\vspace{2mm}
%beginitem
5. Яку часову оцінку складності має алгоритм сортування вставками?
%enditem
\par
\vspace{2mm}
%beginitem
6. Модифікувати клас точок простору 

\qquad\code{class Point\{ double x, y, z;\};}


додавши до нього оператор перетворення на рядок типу \code{std::string}, що містить рядкове зображення точки простору (у круглих дужках координати через кому), наприклад, (1.01, 2.35, 0.5), та його означення.%enditem
\par
\vspace{2mm}
%beginitem
7. 
Задано клас A (структуру, оператори < та == десь коректно означено) та клас вузла списку Node.

struct A\{

\quad  якісь коректні означення 

\quad  bool operator<(const A\&); 

\quad  bool operator==(const A\&); 

\quad A\& operator =(const A\&)=delete; 

\quad  A\& operator =(A\&\&)=delete;\};

struct Node\{Node *next; A a;\};

 Написати функцію, яка вилучає зі списку всі вузли, елемент в яких більший елемента в попередньому вузлі списку або дорівнює елементу в наступному вузлі. Список є циклічним та задається вказівником на перший вузол списку. Наприклад, зі списку, що містить послідовність 8, 10, 9, 4, мають бути вилучені вузли, що містять 8 (бо список циклічний і $8>4$) та 10. У цьому прикладі 9 залишається в списку.

Крім правильності та відповідності умові оцінюється якість коду, наявність витоків ресурсів та коректність роботи з пам'яттю. 

%enditem
}
{\smallskip\smallskip \begin {scriptsize} Затверджено на засіданні кафедри теоретичної кібернетики, протокол \No 4  від   25.01.2024 р.\par\smallskip\smallskip\smallskip\smallskip
{\parbox {6.5 cm} {Зав. кафедри \dotfill Ю.В.~Крак}}\hspace {0.7cm} {\hbox to 6.5 cm { Екзаменатор \dotfill Т.О.~Карнаух}} \end{scriptsize}}
%end
\newpage
%begin
{{\begin {scriptsize}\par
{ \center  Київський національний університет імені Тараса Шевченка \par}
{\parbox {5 cm}{Освітня програма \hfill}{\parbox {7 cm}{Інформатика\hfill}}\parbox {6 cm}{\hfill Семестр 2}}\par
{\parbox {5 cm} {Навчальний предмет \hfill}\parbox {7 cm}{Програмування \hfill}\parbox {6cm}{\hfill}\par}\vspace{-15pt} \end{scriptsize}}{\center Екзаменаційний білет \No \textbf { 22 }\par}}
{%beginitem
1. Вкажіть основну відмінність між відкритими (public) та прихованими (private) методами класу.%enditem
\par
\vspace{2mm}
%beginitem
2. Максимально змістовно повно сформулюйте назву оголошеного в наведеному фрагменті методу класу.
\vspace{-3mm}
\begin{verbatim}class A{
   ...
   ~A();
};\end{verbatim}
\vspace{-3mm}
%enditem
\par
\vspace{2mm}
%beginitem
3. Вкажіть, що не так з наведеним кодом.
\vspace{-3mm}
\begin{verbatim}int *pa=new int(2), *pb=&pa; delete pa;\end{verbatim}
\vspace{-3mm}
%enditem
\par
\vspace{2mm}
%beginitem
4. За наявності означень
\vspace{-3mm}
\begin{verbatim}
class A{ };
class B{ };
class C:public A{B b;A a;};\end{verbatim}
\vspace{-3mm}
вкажіть послідовність викликів конструкторів за виконання коду \code{\{C c;\}} 

(Якщо конструктор викликається кілька разів поспіль, то має зазначатися кожен його виклик, наприклад, A, A, A.)
%enditem

\vspace{2mm}
%beginitem
5. Яку часову оцінку складності має алгоритм швидкого сортування?
%enditem
\par
\vspace{2mm}
%beginitem
6. Модифікувати клас 

\qquad\code{class Session\{ int DM, MA, Alg, Prog;\};} 


додавши до нього оператор порівняння \code{<}: один об'єкт є меншим за інший тоді і тільки тоді, коли значення всіх атрибутів об'єктів не менші за 60 і сума атрибутів першого менша за суму атрибутів другого. Означення оператора теж має бути записане.%enditem
\par
\vspace{2mm}
%beginitem
7. 
Реалізувати операцію перетину множин з повтореннями. Множина з повтореннями задається переліком своїх елементів, записаних за неспаданням у текстовому файлі через білі символи. Результат у такому ж самому форматі (і теж за неспаданням) має бути записаний у текстовий файл. Питома функція має отримувати три шляхи до файлів. Розмір вхідних множин такий, що вони не можуть бути завантажені повністю в пам’ять програми.

Крім правильності та відповідності умові оцінюється якість коду, наявність витоків ресурсів та коректність роботи з пам'яттю. 

%enditem
}
{\smallskip\smallskip \begin {scriptsize} Затверджено на засіданні кафедри теоретичної кібернетики, протокол \No 4  від   25.01.2024 р.\par\smallskip\smallskip\smallskip\smallskip
{\parbox {6.5 cm} {Зав. кафедри \dotfill Ю.В.~Крак}}\hspace {0.7cm} {\hbox to 6.5 cm { Екзаменатор \dotfill Т.О.~Карнаух}} \end{scriptsize}}
%end
\newpage
%begin
{{\begin {scriptsize}\par
{ \center  Київський національний університет імені Тараса Шевченка \par}
{\parbox {5 cm}{Освітня програма \hfill}{\parbox {7 cm}{Інформатика\hfill}}\parbox {6 cm}{\hfill Семестр 2}}\par
{\parbox {5 cm} {Навчальний предмет \hfill}\parbox {7 cm}{Програмування \hfill}\parbox {6cm}{\hfill}\par}\vspace{-15pt} \end{scriptsize}}{\center Екзаменаційний білет \No \textbf { 23 }\par}}
{%beginitem
1. Вкажіть основну відмінність між класами (class) та структурами (struct).%enditem
\par
\vspace{2mm}
%beginitem
2. Максимально змістовно повно сформулюйте назву оголошеного в наведеному фрагменті методу класу.
\vspace{-3mm}
\begin{verbatim}class A{
   ...
   A(const A&);
};\end{verbatim}
\vspace{-3mm}
%enditem
\par
\vspace{2mm}
%beginitem
3. Вкажіть, що не так з наведеним кодом.
\vspace{-3mm}
\begin{verbatim}int *pa=new int(2), *pb=new int(3); pa=pb;\end{verbatim}
\vspace{-3mm}
%enditem
\par
\vspace{2mm}
%beginitem
4. За наявності означень
\vspace{-3mm}
\begin{verbatim}
class A{ };
class B{ };
class C{B b;A a;};\end{verbatim}
\vspace{-3mm}
вкажіть послідовність викликів конструкторів за виконання коду \code{\{C c;\}} 

(Якщо конструктор викликається кілька разів поспіль, то має зазначатися кожен його виклик, наприклад, A, A, A.)
%enditem

\vspace{2mm}
%beginitem
5. Яку часову оцінку складності має алгоритм сортування злиттям?
%enditem
\par
\vspace{2mm}
%beginitem
6. Модифікувати клас 

\qquad\code{class Session\{ int DM, MA, Alg, Prog;\};} 


додавши до нього оператор перетворення на тип \code{int}: якщо хоча б один атрибут менший за 60, то має повертатись 0, інакше -- сума атрибутів.%enditem
\par
\vspace{2mm}
%beginitem
7. 
Задано клас A (структуру, оператори > та == десь коректно означено) та клас вузла списку Node.

struct A\{

\quad  якісь коректні означення 

\quad  bool operator>(const A\&); 

\quad  bool operator==(const A\&); 

\quad A\& operator =(const A\&)=delete; 

\quad  A\& operator =(A\&\&)=delete;\};

struct Node\{Node *next; A a;\};

 Написати функцію, яка вилучає зі списку всі вузли, елемент в яких менший елемента в попередньому вузлі списку або дорівнює елементу в наступному вузлі. Список задається вказівником на перший вузол списку. Наприклад, зі списку, що містить послідовність 8, 5, 6, 9, 4, мають бути вилучені вузли, що містять 5 та 4. У цьому прикладі 6 залишається в списку.

Крім правильності та відповідності умові оцінюється якість коду, наявність витоків ресурсів та коректність роботи з пам'яттю. 

%enditem
}
{\smallskip\smallskip \begin {scriptsize} Затверджено на засіданні кафедри теоретичної кібернетики, протокол \No 4  від   25.01.2024 р.\par\smallskip\smallskip\smallskip\smallskip
{\parbox {6.5 cm} {Зав. кафедри \dotfill Ю.В.~Крак}}\hspace {0.7cm} {\hbox to 6.5 cm { Екзаменатор \dotfill Т.О.~Карнаух}} \end{scriptsize}}
%end
\newpage
%begin
{{\begin {scriptsize}\par
{ \center  Київський національний університет імені Тараса Шевченка \par}
{\parbox {5 cm}{Освітня програма \hfill}{\parbox {7 cm}{Інформатика\hfill}}\parbox {6 cm}{\hfill Семестр 2}}\par
{\parbox {5 cm} {Навчальний предмет \hfill}\parbox {7 cm}{Програмування \hfill}\parbox {6cm}{\hfill}\par}\vspace{-15pt} \end{scriptsize}}{\center Екзаменаційний білет \No \textbf { 24 }\par}}
{%beginitem
1. Вкажіть основну відмінність між віртуальними методами та тими, що не є віртуальними.%enditem
\par
\vspace{2mm}
%beginitem
2. Максимально змістовно повно сформулюйте назву оголошеного в наведеному фрагменті методу класу.
\vspace{-3mm}
\begin{verbatim}class A{
   ...
   A(A&&);
};\end{verbatim}
\vspace{-3mm}
%enditem
\par
\vspace{2mm}
%beginitem
3. Вкажіть, що не так з наведеним кодом.
\vspace{-3mm}
\begin{verbatim}int a, *p=&a; delete a;\end{verbatim}
\vspace{-3mm}
%enditem
\par
\vspace{2mm}
%beginitem
4. За наявності означень
\vspace{-3mm}
\begin{verbatim}
class A{ };
class B{ };
class C:public A{B b[2];};\end{verbatim}
\vspace{-3mm}
вкажіть послідовність викликів деструкторів за виконання коду \code{\{C c;\}} 

(Якщо деструктор викликається кілька разів поспіль, то має зазначатися кожен його виклик, наприклад, $\sim$A, $\sim$A, $\sim$A.)
%enditem

\vspace{2mm}
%beginitem
5. Яку часову оцінку складності має алгоритм сортування купою (він же деревом, пірамідою)?
%enditem
\par
\vspace{2mm}
%beginitem
6. Модифікувати клас точок простору 

\qquad\code{class Point\{ double x, y, z;\};}


додавши до нього оператор перетворення на тип \code{bool}: точка є істинною тоді і тільки тоді, коли містить хоча б одну ненульову координату. Означення оператора теж має бути записане.%enditem
\par
\vspace{2mm}
%beginitem
7. 
Реалізувати операцію різниці множин з повтореннями. Множина з повтореннями задається переліком своїх елементів, записаних за неспаданням у текстовому файлі через білі символи. Результат у такому ж самому форматі (і теж за неспаданням) має бути записаний у текстовий файл. Питома функція має отримувати три шляхи до файлів. Розмір вхідних множин такий, що вони не можуть бути завантажені повністю в пам’ять програми.

Крім правильності та відповідності умові оцінюється якість коду, наявність витоків ресурсів та коректність роботи з пам'яттю. 

%enditem
}
{\smallskip\smallskip \begin {scriptsize} Затверджено на засіданні кафедри теоретичної кібернетики, протокол \No 4  від   25.01.2024 р.\par\smallskip\smallskip\smallskip\smallskip
{\parbox {6.5 cm} {Зав. кафедри \dotfill Ю.В.~Крак}}\hspace {0.7cm} {\hbox to 6.5 cm { Екзаменатор \dotfill Т.О.~Карнаух}} \end{scriptsize}}
%end
\newpage
%begin
{{\begin {scriptsize}\par
{ \center  Київський національний університет імені Тараса Шевченка \par}
{\parbox {5 cm}{Освітня програма \hfill}{\parbox {7 cm}{Інформатика\hfill}}\parbox {6 cm}{\hfill Семестр 2}}\par
{\parbox {5 cm} {Навчальний предмет \hfill}\parbox {7 cm}{Програмування \hfill}\parbox {6cm}{\hfill}\par}\vspace{-15pt} \end{scriptsize}}{\center Екзаменаційний білет \No \textbf { 25 }\par}}
{%beginitem
1. Вкажіть основну відмінність між класами (class) та структурами (struct).%enditem
\par
\vspace{2mm}
%beginitem
2. Максимально змістовно повно сформулюйте назву оголошеного в наведеному фрагменті методу класу.
\vspace{-3mm}
\begin{verbatim}class A{
   ...
   A();
};\end{verbatim}
\vspace{-3mm}
%enditem
\par
\vspace{2mm}
%beginitem
3. Вкажіть, що не так з наведеним кодом.
\vspace{-3mm}
\begin{verbatim}int *pa=new int(2), *pb=pa; delete pb; delete pa;\end{verbatim}
\vspace{-3mm}
%enditem
\par
\vspace{2mm}
%beginitem
4. За наявності означень
\vspace{-3mm}
\begin{verbatim}
class A{ };
class B{ };
class C:public A{B b[2];};\end{verbatim}
\vspace{-3mm}
вкажіть послідовність викликів деструкторів за виконання коду \code{\{C c;\}} 

(Якщо деструктор викликається кілька разів поспіль, то має зазначатися кожен його виклик, наприклад, $\sim$A, $\sim$A, $\sim$A.)
%enditem

\vspace{2mm}
%beginitem
5. Яку часову оцінку складності має алгоритм швидкого сортування?
%enditem
\par
\vspace{2mm}
%beginitem
6. Модифікувати клас 

\qquad\code{class Session\{ int DM, MA, Alg, Prog;\};} 


додавши до нього оператор перетворення на тип \code{bool}: об'єкт класу є істинним тоді і тільки тоді, коли значення всіх атрибутів не менші за 60. Означення оператора теж має бути записане.%enditem
\par
\vspace{2mm}
%beginitem
7. 
Реалізувати операцію симетричної різниці множин з повтореннями. Множина з повтореннями задається переліком своїх елементів, записаних за неспаданням у текстовому файлі через білі символи. Результат у такому ж самому форматі (і теж за неспаданням) має бути записаний у текстовий файл. Питома функція має отримувати три шляхи до файлів. Розмір вхідних множин такий, що вони не можуть бути завантажені повністю в пам’ять програми.

Крім правильності та відповідності умові оцінюється якість коду, наявність витоків ресурсів та коректність роботи з пам'яттю. 

%enditem
}
{\smallskip\smallskip \begin {scriptsize} Затверджено на засіданні кафедри теоретичної кібернетики, протокол \No 4  від   25.01.2024 р.\par\smallskip\smallskip\smallskip\smallskip
{\parbox {6.5 cm} {Зав. кафедри \dotfill Ю.В.~Крак}}\hspace {0.7cm} {\hbox to 6.5 cm { Екзаменатор \dotfill Т.О.~Карнаух}} \end{scriptsize}}
%end
\newpage
%begin
{{\begin {scriptsize}\par
{ \center  Київський національний університет імені Тараса Шевченка \par}
{\parbox {5 cm}{Освітня програма \hfill}{\parbox {7 cm}{Інформатика\hfill}}\parbox {6 cm}{\hfill Семестр 2}}\par
{\parbox {5 cm} {Навчальний предмет \hfill}\parbox {7 cm}{Програмування \hfill}\parbox {6cm}{\hfill}\par}\vspace{-15pt} \end{scriptsize}}{\center Екзаменаційний білет \No \textbf { 26 }\par}}
{%beginitem
1. Вкажіть основну \textbf{семантичну} відмінність між конструктором копії та конструктором переміщення.%enditem
\par
\vspace{2mm}
%beginitem
2. Максимально змістовно повно сформулюйте назву оголошеного в наведеному фрагменті методу класу.
\vspace{-3mm}
\begin{verbatim}class A{
   ...
   ~A();
};\end{verbatim}
\vspace{-3mm}
%enditem
\par
\vspace{2mm}
%beginitem
3. Вкажіть, що не так з наведеним кодом.
\vspace{-3mm}
\begin{verbatim}int *pa=new int(2), *pb=pa; delete pb; *pa=3;\end{verbatim}
\vspace{-3mm}
%enditem
\par
\vspace{2mm}
%beginitem
4. За наявності означень
\vspace{-3mm}
\begin{verbatim}
class A{ };
class B{ };
class C:public A{B b;};\end{verbatim}
\vspace{-3mm}
вкажіть послідовність викликів конструкторів за виконання коду \code{\{C c;\}} 

(Якщо конструктор викликається кілька разів поспіль, то має зазначатися кожен його виклик, наприклад, A, A, A.)
%enditem

\vspace{2mm}
%beginitem
5. Яку часову оцінку складності має алгоритм сортування бульбашкою?
%enditem
\par
\vspace{2mm}
%beginitem
6. Модифікувати клас точок простору 

\qquad\code{class Point\{ double x, y, z;\};}


додавши до нього оператор перетворення на тип \code{bool}: точка є істинною тоді і тільки тоді, коли містить хоча б одну ненульову координату. Означення оператора теж має бути записане.%enditem
\par
\vspace{2mm}
%beginitem
7. 
Реалізувати операцію злиття множин з повтореннями. Множина з повтореннями задається переліком своїх елементів, записаних за неспаданням у текстовому файлі через білі символи. Результат у такому ж самому форматі (і теж за неспаданням) має бути записаний у текстовий файл. Питома функція має отримувати три шляхи до файлів. Розмір вхідних множин такий, що вони не можуть бути завантажені повністю в пам’ять програми.

Крім правильності та відповідності умові оцінюється якість коду, наявність витоків ресурсів та коректність роботи з пам'яттю. 

%enditem
}
{\smallskip\smallskip \begin {scriptsize} Затверджено на засіданні кафедри теоретичної кібернетики, протокол \No 4  від   25.01.2024 р.\par\smallskip\smallskip\smallskip\smallskip
{\parbox {6.5 cm} {Зав. кафедри \dotfill Ю.В.~Крак}}\hspace {0.7cm} {\hbox to 6.5 cm { Екзаменатор \dotfill Т.О.~Карнаух}} \end{scriptsize}}
%end
\newpage
%begin
{{\begin {scriptsize}\par
{ \center  Київський національний університет імені Тараса Шевченка \par}
{\parbox {5 cm}{Освітня програма \hfill}{\parbox {7 cm}{Інформатика\hfill}}\parbox {6 cm}{\hfill Семестр 2}}\par
{\parbox {5 cm} {Навчальний предмет \hfill}\parbox {7 cm}{Програмування \hfill}\parbox {6cm}{\hfill}\par}\vspace{-15pt} \end{scriptsize}}{\center Екзаменаційний білет \No \textbf { 27 }\par}}
{%beginitem
1. Що важливо забезпечити при роботі з динамічними змінними?%enditem
\par
\vspace{2mm}
%beginitem
2. Максимально змістовно повно сформулюйте назву оголошеного в наведеному фрагменті методу класу.
\vspace{-3mm}
\begin{verbatim}class A{
   ...
   A& operator =(A&&);
};\end{verbatim}
\vspace{-3mm}
%enditem
\par
\vspace{2mm}
%beginitem
3. Вкажіть, що не так з наведеним кодом.
\vspace{-3mm}
\begin{verbatim}int **ppi=new int*(new int(3)); delete ppi;\end{verbatim}
\vspace{-3mm}
%enditem
\par
\vspace{2mm}
%beginitem
4. За наявності означень
\vspace{-3mm}
\begin{verbatim}
class A{ };
class B{ };
class C:public A{B b[2];};\end{verbatim}
\vspace{-3mm}
вкажіть послідовність викликів конструкторів за виконання коду \code{\{C c;\}} 

(Якщо конструктор викликається кілька разів поспіль, то має зазначатися кожен його виклик, наприклад, A, A, A.)
%enditem

\vspace{2mm}
%beginitem
5. Яку часову оцінку складності має алгоритм сортування вибором?
%enditem
\par
\vspace{2mm}
%beginitem
6. Модифікувати клас точок простору 

\qquad\code{class Point\{ double x, y, z;\};}


додавши до нього оператор перетворення на рядок типу \code{std::string}, що містить рядкове зображення точки простору (у круглих дужках координати через кому), наприклад, (1.01, 2.35, 0.5), та його означення.%enditem
\par
\vspace{2mm}
%beginitem
7. 
Задано клас A (структуру, оператори < та == десь коректно означено) та клас вузла списку Node.

struct A\{

\quad  якісь коректні означення 

\quad  bool operator<(const A\&); 

\quad  bool operator==(const A\&); 

\quad A\& operator =(const A\&)=delete; 

\quad  A\& operator =(A\&\&)=delete;\};

struct Node\{Node *next; A a;\};

 Написати функцію, яка вилучає зі списку всі вузли, елемент в яких більший елемента в попередньому вузлі списку або дорівнює елементу в наступному вузлі. Список є циклічним та задається вказівником на перший вузол списку. Наприклад, зі списку, що містить послідовність 8, 10, 9, 4, мають бути вилучені вузли, що містять 8 (бо список циклічний і $8>4$) та 10. У цьому прикладі 9 залишається в списку.

Крім правильності та відповідності умові оцінюється якість коду, наявність витоків ресурсів та коректність роботи з пам'яттю. 

%enditem
}
{\smallskip\smallskip \begin {scriptsize} Затверджено на засіданні кафедри теоретичної кібернетики, протокол \No 4  від   25.01.2024 р.\par\smallskip\smallskip\smallskip\smallskip
{\parbox {6.5 cm} {Зав. кафедри \dotfill Ю.В.~Крак}}\hspace {0.7cm} {\hbox to 6.5 cm { Екзаменатор \dotfill Т.О.~Карнаух}} \end{scriptsize}}
%end
\newpage
%begin
{{\begin {scriptsize}\par
{ \center  Київський національний університет імені Тараса Шевченка \par}
{\parbox {5 cm}{Освітня програма \hfill}{\parbox {7 cm}{Інформатика\hfill}}\parbox {6 cm}{\hfill Семестр 2}}\par
{\parbox {5 cm} {Навчальний предмет \hfill}\parbox {7 cm}{Програмування \hfill}\parbox {6cm}{\hfill}\par}\vspace{-15pt} \end{scriptsize}}{\center Екзаменаційний білет \No \textbf { 28 }\par}}
{%beginitem
1. Вкажіть основну \textbf{синтаксичну} відмінність між конструктором копії та конструктором переміщення.%enditem
\par
\vspace{2mm}
%beginitem
2. Максимально змістовно повно сформулюйте назву оголошеного в наведеному фрагменті методу класу.
\vspace{-3mm}
\begin{verbatim}class A{
   ...
   A(A&&);
};\end{verbatim}
\vspace{-3mm}
%enditem
\par
\vspace{2mm}
%beginitem
3. Вкажіть, що не так з наведеним кодом.
\vspace{-3mm}
\begin{verbatim}int *pa=new int(2), *pb=*pa; delete pa;\end{verbatim}
\vspace{-3mm}
%enditem
\par
\vspace{2mm}
%beginitem
4. За наявності означень
\vspace{-3mm}
\begin{verbatim}
class A{ };
class B:public A{ };
class C:public B{ };\end{verbatim}
\vspace{-3mm}
вкажіть послідовність викликів деструкторів за виконання коду \code{\{C c;\}} 

(Якщо деструктор викликається кілька разів поспіль, то має зазначатися кожен його виклик, наприклад, $\sim$A, $\sim$A, $\sim$A.)
%enditem

\vspace{2mm}
%beginitem
5. Яку часову оцінку складності має алгоритм сортування вставками?
%enditem
\par
\vspace{2mm}
%beginitem
6. Модифікувати клас 

\qquad\code{class Session\{ int DM, MA, Alg, Prog;\};} 


додавши до нього оператор порівняння \code{<}: один об'єкт є меншим за інший тоді і тільки тоді, коли значення всіх атрибутів об'єктів не менші за 60 і сума атрибутів першого менша за суму атрибутів другого. Означення оператора теж має бути записане.%enditem
\par
\vspace{2mm}
%beginitem
7. 
Задано клас A (структуру, оператори < та == десь коректно означено) та клас вузла списку Node.

struct A\{

\quad  якісь коректні означення 

\quad  bool operator<(const A\&); 

\quad A\& operator =(const A\&)=delete; 

\quad  A\& operator =(A\&\&)=delete;\};

struct Node\{Node *prev; Node *next; A a;\};

Написати функцію, яка вилучає зі списку всі вузли, елемент в яких менший елемента в попередньому вузлі списку або дорівнює елементу в наступному вузлі. Список задається вказівником на перший вузол списку. Наприклад, зі списку, що містить послідовність 8, 5, 6, 9, 4, мають бути вилучені вузли, що містять 5 та 4. У цьому прикладі 6 залишається в списку.

Крім правильності та відповідності умові оцінюється якість коду, наявність витоків ресурсів та коректність роботи з пам'яттю. 

%enditem
}
{\smallskip\smallskip \begin {scriptsize} Затверджено на засіданні кафедри теоретичної кібернетики, протокол \No 4  від   25.01.2024 р.\par\smallskip\smallskip\smallskip\smallskip
{\parbox {6.5 cm} {Зав. кафедри \dotfill Ю.В.~Крак}}\hspace {0.7cm} {\hbox to 6.5 cm { Екзаменатор \dotfill Т.О.~Карнаух}} \end{scriptsize}}
%end
\newpage
%begin
{{\begin {scriptsize}\par
{ \center  Київський національний університет імені Тараса Шевченка \par}
{\parbox {5 cm}{Освітня програма \hfill}{\parbox {7 cm}{Інформатика\hfill}}\parbox {6 cm}{\hfill Семестр 2}}\par
{\parbox {5 cm} {Навчальний предмет \hfill}\parbox {7 cm}{Програмування \hfill}\parbox {6cm}{\hfill}\par}\vspace{-15pt} \end{scriptsize}}{\center Екзаменаційний білет \No \textbf { 29 }\par}}
{%beginitem
1. Вкажіть основну відмінність між віртуальними методами та тими, що не є віртуальними.%enditem
\par
\vspace{2mm}
%beginitem
2. Максимально змістовно повно сформулюйте назву оголошеного в наведеному фрагменті методу класу.
\vspace{-3mm}
\begin{verbatim}class A{
   ...
   A(const A&);
};\end{verbatim}
\vspace{-3mm}
%enditem
\par
\vspace{2mm}
%beginitem
3. Вкажіть, що не так з наведеним кодом.
\vspace{-3mm}
\begin{verbatim}int a, *p=&a; delete a;\end{verbatim}
\vspace{-3mm}
%enditem
\par
\vspace{2mm}
%beginitem
4. За наявності означень
\vspace{-3mm}
\begin{verbatim}
class A{ };
class B{ };
class C:public A{B b;};\end{verbatim}
\vspace{-3mm}
вкажіть послідовність викликів деструкторів за виконання коду \code{\{C c;\}} 

(Якщо деструктор викликається кілька разів поспіль, то має зазначатися кожен його виклик, наприклад, $\sim$A, $\sim$A, $\sim$A.)
%enditem

\vspace{2mm}
%beginitem
5. Яку часову оцінку складності має алгоритм сортування вибором?
%enditem
\par
\vspace{2mm}
%beginitem
6. Модифікувати клас 

\qquad\code{class Session\{ int DM, MA, Alg, Prog;\};} 


додавши до нього оператор перетворення на тип \code{bool}: об'єкт класу є істинним тоді і тільки тоді, коли значення всіх атрибутів не менші за 60. Означення оператора теж має бути записане.%enditem
\par
\vspace{2mm}
%beginitem
7. 
Реалізувати операцію перетину множин з повтореннями. Множина з повтореннями задається переліком своїх елементів, записаних за неспаданням у текстовому файлі через білі символи. Результат у такому ж самому форматі (і теж за неспаданням) має бути записаний у текстовий файл. Питома функція має отримувати три шляхи до файлів. Розмір вхідних множин такий, що вони не можуть бути завантажені повністю в пам’ять програми.

Крім правильності та відповідності умові оцінюється якість коду, наявність витоків ресурсів та коректність роботи з пам'яттю. 

%enditem
}
{\smallskip\smallskip \begin {scriptsize} Затверджено на засіданні кафедри теоретичної кібернетики, протокол \No 4  від   25.01.2024 р.\par\smallskip\smallskip\smallskip\smallskip
{\parbox {6.5 cm} {Зав. кафедри \dotfill Ю.В.~Крак}}\hspace {0.7cm} {\hbox to 6.5 cm { Екзаменатор \dotfill Т.О.~Карнаух}} \end{scriptsize}}
%end
\newpage
%begin
{{\begin {scriptsize}\par
{ \center  Київський національний університет імені Тараса Шевченка \par}
{\parbox {5 cm}{Освітня програма \hfill}{\parbox {7 cm}{Інформатика\hfill}}\parbox {6 cm}{\hfill Семестр 2}}\par
{\parbox {5 cm} {Навчальний предмет \hfill}\parbox {7 cm}{Програмування \hfill}\parbox {6cm}{\hfill}\par}\vspace{-15pt} \end{scriptsize}}{\center Екзаменаційний білет \No \textbf { 30 }\par}}
{%beginitem
1. Вкажіть основну відмінність між відкритими (public) та прихованими (private) методами класу.%enditem
\par
\vspace{2mm}
%beginitem
2. Максимально змістовно повно сформулюйте назву оголошеного в наведеному фрагменті методу класу.
\vspace{-3mm}
\begin{verbatim}class A{
   ...
   A();
};\end{verbatim}
\vspace{-3mm}
%enditem
\par
\vspace{2mm}
%beginitem
3. Вкажіть, що не так з наведеним кодом.
\vspace{-3mm}
\begin{verbatim}int **ppi=new int*(new int(3)); delete ppi;\end{verbatim}
\vspace{-3mm}
%enditem
\par
\vspace{2mm}
%beginitem
4. За наявності означень
\vspace{-3mm}
\begin{verbatim}
class A{ };
class B{ };
class C{B b;A a;};\end{verbatim}
\vspace{-3mm}
вкажіть послідовність викликів конструкторів за виконання коду \code{\{C c;\}} 

(Якщо конструктор викликається кілька разів поспіль, то має зазначатися кожен його виклик, наприклад, A, A, A.)
%enditem

\vspace{2mm}
%beginitem
5. Яку часову оцінку складності має алгоритм сортування вставками?
%enditem
\par
\vspace{2mm}
%beginitem
6. Модифікувати клас 

\qquad\code{class Session\{ int DM, MA, Alg, Prog;\};} 


додавши до нього оператор перетворення на тип \code{int}: якщо хоча б один атрибут менший за 60, то має повертатись 0, інакше -- сума атрибутів.%enditem
\par
\vspace{2mm}
%beginitem
7. 
Задано клас A (структуру, оператори > та == десь коректно означено) та клас вузла списку Node.

struct A\{

\quad  якісь коректні означення 

\quad  bool operator>(const A\&); 

\quad  bool operator==(const A\&); 

\quad A\& operator =(const A\&)=delete; 

\quad  A\& operator =(A\&\&)=delete;\};

struct Node\{Node *next; A a;\};

 Написати функцію, яка вилучає зі списку всі вузли, елемент в яких менший елемента в попередньому вузлі списку або дорівнює елементу в наступному вузлі. Список задається вказівником на перший вузол списку. Наприклад, зі списку, що містить послідовність 8, 5, 6, 9, 4, мають бути вилучені вузли, що містять 5 та 4. У цьому прикладі 6 залишається в списку.

Крім правильності та відповідності умові оцінюється якість коду, наявність витоків ресурсів та коректність роботи з пам'яттю. 

%enditem
}
{\smallskip\smallskip \begin {scriptsize} Затверджено на засіданні кафедри теоретичної кібернетики, протокол \No 4  від   25.01.2024 р.\par\smallskip\smallskip\smallskip\smallskip
{\parbox {6.5 cm} {Зав. кафедри \dotfill Ю.В.~Крак}}\hspace {0.7cm} {\hbox to 6.5 cm { Екзаменатор \dotfill Т.О.~Карнаух}} \end{scriptsize}}
%end
\newpage
%begin
{{\begin {scriptsize}\par
{ \center  Київський національний університет імені Тараса Шевченка \par}
{\parbox {5 cm}{Освітня програма \hfill}{\parbox {7 cm}{Інформатика\hfill}}\parbox {6 cm}{\hfill Семестр 2}}\par
{\parbox {5 cm} {Навчальний предмет \hfill}\parbox {7 cm}{Програмування \hfill}\parbox {6cm}{\hfill}\par}\vspace{-15pt} \end{scriptsize}}{\center Екзаменаційний білет \No \textbf { 31 }\par}}
{%beginitem
1. Вкажіть основну відмінність між віртуальними методами та тими, що не є віртуальними.%enditem
\par
\vspace{2mm}
%beginitem
2. Максимально змістовно повно сформулюйте назву оголошеного в наведеному фрагменті методу класу.
\vspace{-3mm}
\begin{verbatim}class A{
   ...
   A(const A&);
};\end{verbatim}
\vspace{-3mm}
%enditem
\par
\vspace{2mm}
%beginitem
3. Вкажіть, що не так з наведеним кодом.
\vspace{-3mm}
\begin{verbatim}int *pa=new int(2), *pb=&pa; delete pa;\end{verbatim}
\vspace{-3mm}
%enditem
\par
\vspace{2mm}
%beginitem
4. За наявності означень
\vspace{-3mm}
\begin{verbatim}
class A{ };
class B:public A{ };
class C:public B{ };\end{verbatim}
\vspace{-3mm}
вкажіть послідовність викликів конструкторів за виконання коду \code{\{C c;\}} 

(Якщо конструктор викликається кілька разів поспіль, то має зазначатися кожен його виклик, наприклад, A, A, A.)
%enditem

\vspace{2mm}
%beginitem
5. Яку часову оцінку складності має алгоритм сортування купою (він же деревом, пірамідою)?
%enditem
\par
\vspace{2mm}
%beginitem
6. Модифікувати клас точок простору 

\qquad\code{class Point\{ double x, y, z;\};}


додавши до нього оператор перетворення на тип \code{bool}: точка є істинною тоді і тільки тоді, коли містить хоча б одну ненульову координату. Означення оператора теж має бути записане.%enditem
\par
\vspace{2mm}
%beginitem
7. 
Реалізувати операцію різниці множин з повтореннями. Множина з повтореннями задається переліком своїх елементів, записаних за неспаданням у текстовому файлі через білі символи. Результат у такому ж самому форматі (і теж за неспаданням) має бути записаний у текстовий файл. Питома функція має отримувати три шляхи до файлів. Розмір вхідних множин такий, що вони не можуть бути завантажені повністю в пам’ять програми.

Крім правильності та відповідності умові оцінюється якість коду, наявність витоків ресурсів та коректність роботи з пам'яттю. 

%enditem
}
{\smallskip\smallskip \begin {scriptsize} Затверджено на засіданні кафедри теоретичної кібернетики, протокол \No 4  від   25.01.2024 р.\par\smallskip\smallskip\smallskip\smallskip
{\parbox {6.5 cm} {Зав. кафедри \dotfill Ю.В.~Крак}}\hspace {0.7cm} {\hbox to 6.5 cm { Екзаменатор \dotfill Т.О.~Карнаух}} \end{scriptsize}}
%end
\newpage
%begin
{{\begin {scriptsize}\par
{ \center  Київський національний університет імені Тараса Шевченка \par}
{\parbox {5 cm}{Освітня програма \hfill}{\parbox {7 cm}{Інформатика\hfill}}\parbox {6 cm}{\hfill Семестр 2}}\par
{\parbox {5 cm} {Навчальний предмет \hfill}\parbox {7 cm}{Програмування \hfill}\parbox {6cm}{\hfill}\par}\vspace{-15pt} \end{scriptsize}}{\center Екзаменаційний білет \No \textbf { 32 }\par}}
{%beginitem
1. Що важливо забезпечити при роботі з динамічними змінними?%enditem
\par
\vspace{2mm}
%beginitem
2. Максимально змістовно повно сформулюйте назву оголошеного в наведеному фрагменті методу класу.
\vspace{-3mm}
\begin{verbatim}class A{
   ...
   A();
};\end{verbatim}
\vspace{-3mm}
%enditem
\par
\vspace{2mm}
%beginitem
3. Вкажіть, що не так з наведеним кодом.
\vspace{-3mm}
\begin{verbatim}int *pa=new int(2), *pb=new int(3); pa=pb;\end{verbatim}
\vspace{-3mm}
%enditem
\par
\vspace{2mm}
%beginitem
4. За наявності означень
\vspace{-3mm}
\begin{verbatim}
class A{ };
class B:public A{ };
class C:public A{B b;};\end{verbatim}
\vspace{-3mm}
вкажіть послідовність викликів деструкторів за виконання коду \code{\{C c;\}} 

(Якщо деструктор викликається кілька разів поспіль, то має зазначатися кожен його виклик, наприклад, $\sim$A, $\sim$A, $\sim$A.)
%enditem

\vspace{2mm}
%beginitem
5. Яку часову оцінку складності має алгоритм швидкого сортування?
%enditem
\par
\vspace{2mm}
%beginitem
6. Модифікувати клас 

\qquad\code{class Session\{ int DM, MA, Alg, Prog;\};} 


додавши до нього оператор перетворення на тип \code{int}: якщо хоча б один атрибут менший за 60, то має повертатись 0, інакше -- сума атрибутів.%enditem
\par
\vspace{2mm}
%beginitem
7. 
Задано клас A (структуру, оператори < та == десь коректно означено) та клас вузла списку Node.

struct A\{

\quad  якісь коректні означення 

\quad  bool operator<(const A\&); 

\quad  bool operator==(const A\&); 

\quad A\& operator =(const A\&)=delete; 

\quad  A\& operator =(A\&\&)=delete;\};

struct Node\{Node *prev; Node *next; A a;\};

Написати функцію, яка вилучає зі списку всі вузли, елемент в яких більший елемента в попередньому вузлі списку або дорівнює елементу в наступному вузлі. Список є циклічним та задається вказівником на перший вузол списку. Наприклад, зі списку, що містить послідовність 8, 10, 9, 4, мають бути вилучені вузли, що містять 8 (бо список циклічний і $8>4$) та 10. У цьому прикладі 9 залишається в списку.

Крім правильності та відповідності умові оцінюється якість коду, наявність витоків ресурсів та коректність роботи з пам'яттю. 

%enditem
}
{\smallskip\smallskip \begin {scriptsize} Затверджено на засіданні кафедри теоретичної кібернетики, протокол \No 4  від   25.01.2024 р.\par\smallskip\smallskip\smallskip\smallskip
{\parbox {6.5 cm} {Зав. кафедри \dotfill Ю.В.~Крак}}\hspace {0.7cm} {\hbox to 6.5 cm { Екзаменатор \dotfill Т.О.~Карнаух}} \end{scriptsize}}
%end
\newpage
%begin
{{\begin {scriptsize}\par
{ \center  Київський національний університет імені Тараса Шевченка \par}
{\parbox {5 cm}{Освітня програма \hfill}{\parbox {7 cm}{Інформатика\hfill}}\parbox {6 cm}{\hfill Семестр 2}}\par
{\parbox {5 cm} {Навчальний предмет \hfill}\parbox {7 cm}{Програмування \hfill}\parbox {6cm}{\hfill}\par}\vspace{-15pt} \end{scriptsize}}{\center Екзаменаційний білет \No \textbf { 33 }\par}}
{%beginitem
1. Вкажіть основну \textbf{синтаксичну} відмінність між конструктором копії та конструктором переміщення.%enditem
\par
\vspace{2mm}
%beginitem
2. Максимально змістовно повно сформулюйте назву оголошеного в наведеному фрагменті методу класу.
\vspace{-3mm}
\begin{verbatim}class A{
   ...
   ~A();
};\end{verbatim}
\vspace{-3mm}
%enditem
\par
\vspace{2mm}
%beginitem
3. Вкажіть, що не так з наведеним кодом.
\vspace{-3mm}
\begin{verbatim}int *pa=new int(2), *pb=*pa; delete pa;\end{verbatim}
\vspace{-3mm}
%enditem
\par
\vspace{2mm}
%beginitem
4. За наявності означень
\vspace{-3mm}
\begin{verbatim}
class A{ };
class B{ };
class C{B b;A a;};\end{verbatim}
\vspace{-3mm}
вкажіть послідовність викликів деструкторів за виконання коду \code{\{C c;\}} 

(Якщо деструктор викликається кілька разів поспіль, то має зазначатися кожен його виклик, наприклад, $\sim$A, $\sim$A, $\sim$A.)
%enditem

\vspace{2mm}
%beginitem
5. Яку часову оцінку складності має алгоритм сортування бульбашкою?
%enditem
\par
\vspace{2mm}
%beginitem
6. Модифікувати клас 

\qquad\code{class Session\{ int DM, MA, Alg, Prog;\};} 


додавши до нього оператор порівняння \code{<}: один об'єкт є меншим за інший тоді і тільки тоді, коли значення всіх атрибутів об'єктів не менші за 60 і сума атрибутів першого менша за суму атрибутів другого. Означення оператора теж має бути записане.%enditem
\par
\vspace{2mm}
%beginitem
7. 
Задано клас A (структуру, оператори < та == десь коректно означено) та клас вузла списку Node.

struct A\{

\quad  якісь коректні означення 

\quad  bool operator<(const A\&); 

\quad  bool operator==(const A\&); 

\quad A\& operator =(const A\&)=delete; 

\quad  A\& operator =(A\&\&)=delete;\};

struct Node\{Node *next; A a;\};

 Написати функцію, яка вилучає зі списку всі вузли, елемент в яких більший елемента в попередньому вузлі списку або дорівнює елементу в наступному вузлі. Список є циклічним та задається вказівником на перший вузол списку. Наприклад, зі списку, що містить послідовність 8, 10, 9, 4, мають бути вилучені вузли, що містять 8 (бо список циклічний і $8>4$) та 10. У цьому прикладі 9 залишається в списку.

Крім правильності та відповідності умові оцінюється якість коду, наявність витоків ресурсів та коректність роботи з пам'яттю. 

%enditem
}
{\smallskip\smallskip \begin {scriptsize} Затверджено на засіданні кафедри теоретичної кібернетики, протокол \No 4  від   25.01.2024 р.\par\smallskip\smallskip\smallskip\smallskip
{\parbox {6.5 cm} {Зав. кафедри \dotfill Ю.В.~Крак}}\hspace {0.7cm} {\hbox to 6.5 cm { Екзаменатор \dotfill Т.О.~Карнаух}} \end{scriptsize}}
%end
\newpage
%begin
{{\begin {scriptsize}\par
{ \center  Київський національний університет імені Тараса Шевченка \par}
{\parbox {5 cm}{Освітня програма \hfill}{\parbox {7 cm}{Інформатика\hfill}}\parbox {6 cm}{\hfill Семестр 2}}\par
{\parbox {5 cm} {Навчальний предмет \hfill}\parbox {7 cm}{Програмування \hfill}\parbox {6cm}{\hfill}\par}\vspace{-15pt} \end{scriptsize}}{\center Екзаменаційний білет \No \textbf { 34 }\par}}
{%beginitem
1. Вкажіть основну відмінність між класами (class) та структурами (struct).%enditem
\par
\vspace{2mm}
%beginitem
2. Максимально змістовно повно сформулюйте назву оголошеного в наведеному фрагменті методу класу.
\vspace{-3mm}
\begin{verbatim}class A{
   ...
   A(A&&);
};\end{verbatim}
\vspace{-3mm}
%enditem
\par
\vspace{2mm}
%beginitem
3. Вкажіть, що не так з наведеним кодом.
\vspace{-3mm}
\begin{verbatim}int *pa=new int(2), *pb=pa; delete pb; *pa=3;\end{verbatim}
\vspace{-3mm}
%enditem
\par
\vspace{2mm}
%beginitem
4. За наявності означень
\vspace{-3mm}
\begin{verbatim}
class A{ };
class B{ };
class C:public A{B b;A a;};\end{verbatim}
\vspace{-3mm}
вкажіть послідовність викликів деструкторів за виконання коду \code{\{C c;\}} 

(Якщо деструктор викликається кілька разів поспіль, то має зазначатися кожен його виклик, наприклад, $\sim$A, $\sim$A, $\sim$A.)
%enditem

\vspace{2mm}
%beginitem
5. Яку часову оцінку складності має алгоритм швидкого сортування?
%enditem
\par
\vspace{2mm}
%beginitem
6. Модифікувати клас 

\qquad\code{class Session\{ int DM, MA, Alg, Prog;\};} 


додавши до нього оператор перетворення на тип \code{bool}: об'єкт класу є істинним тоді і тільки тоді, коли значення всіх атрибутів не менші за 60. Означення оператора теж має бути записане.%enditem
\par
\vspace{2mm}
%beginitem
7. 
Задано клас A (структуру, оператори < та == десь коректно означено) та клас вузла списку Node.

struct A\{

\quad  якісь коректні означення 

\quad  bool operator<(const A\&); 

\quad  bool operator==(const A\&); 

\quad A\& operator =(const A\&)=delete; 

\quad  A\& operator =(A\&\&)=delete;\};

struct Node\{Node *prev; Node *next; A a;\};

Написати функцію, яка вилучає зі списку всі вузли, елемент в яких більший елемента в попередньому вузлі списку або дорівнює елементу в наступному вузлі. Список є циклічним та задається вказівником на перший вузол списку. Наприклад, зі списку, що містить послідовність 8, 10, 9, 4, мають бути вилучені вузли, що містять 8 (бо список циклічний і $8>4$) та 10. У цьому прикладі 9 залишається в списку.

Крім правильності та відповідності умові оцінюється якість коду, наявність витоків ресурсів та коректність роботи з пам'яттю. 

%enditem
}
{\smallskip\smallskip \begin {scriptsize} Затверджено на засіданні кафедри теоретичної кібернетики, протокол \No 4  від   25.01.2024 р.\par\smallskip\smallskip\smallskip\smallskip
{\parbox {6.5 cm} {Зав. кафедри \dotfill Ю.В.~Крак}}\hspace {0.7cm} {\hbox to 6.5 cm { Екзаменатор \dotfill Т.О.~Карнаух}} \end{scriptsize}}
%end
\newpage
%begin
{{\begin {scriptsize}\par
{ \center  Київський національний університет імені Тараса Шевченка \par}
{\parbox {5 cm}{Освітня програма \hfill}{\parbox {7 cm}{Інформатика\hfill}}\parbox {6 cm}{\hfill Семестр 2}}\par
{\parbox {5 cm} {Навчальний предмет \hfill}\parbox {7 cm}{Програмування \hfill}\parbox {6cm}{\hfill}\par}\vspace{-15pt} \end{scriptsize}}{\center Екзаменаційний білет \No \textbf { 35 }\par}}
{%beginitem
1. Вкажіть основну відмінність між відкритими (public) та прихованими (private) методами класу.%enditem
\par
\vspace{2mm}
%beginitem
2. Максимально змістовно повно сформулюйте назву оголошеного в наведеному фрагменті методу класу.
\vspace{-3mm}
\begin{verbatim}class A{
   ...
   A& operator =(A&&);
};\end{verbatim}
\vspace{-3mm}
%enditem
\par
\vspace{2mm}
%beginitem
3. Вкажіть, що не так з наведеним кодом.
\vspace{-3mm}
\begin{verbatim}int *pa=new int(2), *pb=pa; delete pb; delete pa;\end{verbatim}
\vspace{-3mm}
%enditem
\par
\vspace{2mm}
%beginitem
4. За наявності означень
\vspace{-3mm}
\begin{verbatim}
class A{ };
class B:public A{ };
class C:public A{B b;};\end{verbatim}
\vspace{-3mm}
вкажіть послідовність викликів конструкторів за виконання коду \code{\{C c;\}} 

(Якщо конструктор викликається кілька разів поспіль, то має зазначатися кожен його виклик, наприклад, A, A, A.)
%enditem

\vspace{2mm}
%beginitem
5. Яку часову оцінку складності має алгоритм сортування злиттям?
%enditem
\par
\vspace{2mm}
%beginitem
6. Модифікувати клас точок простору 

\qquad\code{class Point\{ double x, y, z;\};}


додавши до нього оператор перетворення на рядок типу \code{std::string}, що містить рядкове зображення точки простору (у круглих дужках координати через кому), наприклад, (1.01, 2.35, 0.5), та його означення.%enditem
\par
\vspace{2mm}
%beginitem
7. 
Реалізувати операцію різниці множин з повтореннями. Множина з повтореннями задається переліком своїх елементів, записаних за неспаданням у текстовому файлі через білі символи. Результат у такому ж самому форматі (і теж за неспаданням) має бути записаний у текстовий файл. Питома функція має отримувати три шляхи до файлів. Розмір вхідних множин такий, що вони не можуть бути завантажені повністю в пам’ять програми.

Крім правильності та відповідності умові оцінюється якість коду, наявність витоків ресурсів та коректність роботи з пам'яттю. 

%enditem
}
{\smallskip\smallskip \begin {scriptsize} Затверджено на засіданні кафедри теоретичної кібернетики, протокол \No 4  від   25.01.2024 р.\par\smallskip\smallskip\smallskip\smallskip
{\parbox {6.5 cm} {Зав. кафедри \dotfill Ю.В.~Крак}}\hspace {0.7cm} {\hbox to 6.5 cm { Екзаменатор \dotfill Т.О.~Карнаух}} \end{scriptsize}}
%end
\newpage
%begin
{{\begin {scriptsize}\par
{ \center  Київський національний університет імені Тараса Шевченка \par}
{\parbox {5 cm}{Освітня програма \hfill}{\parbox {7 cm}{Інформатика\hfill}}\parbox {6 cm}{\hfill Семестр 2}}\par
{\parbox {5 cm} {Навчальний предмет \hfill}\parbox {7 cm}{Програмування \hfill}\parbox {6cm}{\hfill}\par}\vspace{-15pt} \end{scriptsize}}{\center Екзаменаційний білет \No \textbf { 36 }\par}}
{%beginitem
1. Вкажіть основну \textbf{семантичну} відмінність між конструктором копії та конструктором переміщення.%enditem
\par
\vspace{2mm}
%beginitem
2. Максимально змістовно повно сформулюйте назву оголошеного в наведеному фрагменті методу класу.
\vspace{-3mm}
\begin{verbatim}class A{
   ...
   A(const A&);
};\end{verbatim}
\vspace{-3mm}
%enditem
\par
\vspace{2mm}
%beginitem
3. Вкажіть, що не так з наведеним кодом.
\vspace{-3mm}
\begin{verbatim}int *pa=new int(2), *pb=pa; delete pb; delete pa;\end{verbatim}
\vspace{-3mm}
%enditem
\par
\vspace{2mm}
%beginitem
4. За наявності означень
\vspace{-3mm}
\begin{verbatim}
class A{ };
class B{ };
class C:public A{B b;A a;};\end{verbatim}
\vspace{-3mm}
вкажіть послідовність викликів конструкторів за виконання коду \code{\{C c;\}} 

(Якщо конструктор викликається кілька разів поспіль, то має зазначатися кожен його виклик, наприклад, A, A, A.)
%enditem

\vspace{2mm}
%beginitem
5. Яку часову оцінку складності має алгоритм сортування злиттям?
%enditem
\par
\vspace{2mm}
%beginitem
6. Модифікувати клас 

\qquad\code{class Session\{ int DM, MA, Alg, Prog;\};} 


додавши до нього оператор перетворення на тип \code{int}: якщо хоча б один атрибут менший за 60, то має повертатись 0, інакше -- сума атрибутів.%enditem
\par
\vspace{2mm}
%beginitem
7. 
Реалізувати операцію злиття множин з повтореннями. Множина з повтореннями задається переліком своїх елементів, записаних за неспаданням у текстовому файлі через білі символи. Результат у такому ж самому форматі (і теж за неспаданням) має бути записаний у текстовий файл. Питома функція має отримувати три шляхи до файлів. Розмір вхідних множин такий, що вони не можуть бути завантажені повністю в пам’ять програми.

Крім правильності та відповідності умові оцінюється якість коду, наявність витоків ресурсів та коректність роботи з пам'яттю. 

%enditem
}
{\smallskip\smallskip \begin {scriptsize} Затверджено на засіданні кафедри теоретичної кібернетики, протокол \No 4  від   25.01.2024 р.\par\smallskip\smallskip\smallskip\smallskip
{\parbox {6.5 cm} {Зав. кафедри \dotfill Ю.В.~Крак}}\hspace {0.7cm} {\hbox to 6.5 cm { Екзаменатор \dotfill Т.О.~Карнаух}} \end{scriptsize}}
%end
\newpage
%begin
{{\begin {scriptsize}\par
{ \center  Київський національний університет імені Тараса Шевченка \par}
{\parbox {5 cm}{Освітня програма \hfill}{\parbox {7 cm}{Інформатика\hfill}}\parbox {6 cm}{\hfill Семестр 2}}\par
{\parbox {5 cm} {Навчальний предмет \hfill}\parbox {7 cm}{Програмування \hfill}\parbox {6cm}{\hfill}\par}\vspace{-15pt} \end{scriptsize}}{\center Екзаменаційний білет \No \textbf { 37 }\par}}
{%beginitem
1. Вкажіть основну відмінність між класами (class) та структурами (struct).%enditem
\par
\vspace{2mm}
%beginitem
2. Максимально змістовно повно сформулюйте назву оголошеного в наведеному фрагменті методу класу.
\vspace{-3mm}
\begin{verbatim}class A{
   ...
   ~A();
};\end{verbatim}
\vspace{-3mm}
%enditem
\par
\vspace{2mm}
%beginitem
3. Вкажіть, що не так з наведеним кодом.
\vspace{-3mm}
\begin{verbatim}int a, *p=&a; delete a;\end{verbatim}
\vspace{-3mm}
%enditem
\par
\vspace{2mm}
%beginitem
4. За наявності означень
\vspace{-3mm}
\begin{verbatim}
class A{ };
class B{ };
class C{B b;A a;};\end{verbatim}
\vspace{-3mm}
вкажіть послідовність викликів деструкторів за виконання коду \code{\{C c;\}} 

(Якщо деструктор викликається кілька разів поспіль, то має зазначатися кожен його виклик, наприклад, $\sim$A, $\sim$A, $\sim$A.)
%enditem

\vspace{2mm}
%beginitem
5. Яку часову оцінку складності має алгоритм сортування купою (він же деревом, пірамідою)?
%enditem
\par
\vspace{2mm}
%beginitem
6. Модифікувати клас 

\qquad\code{class Session\{ int DM, MA, Alg, Prog;\};} 


додавши до нього оператор порівняння \code{<}: один об'єкт є меншим за інший тоді і тільки тоді, коли значення всіх атрибутів об'єктів не менші за 60 і сума атрибутів першого менша за суму атрибутів другого. Означення оператора теж має бути записане.%enditem
\par
\vspace{2mm}
%beginitem
7. 
Задано клас A (структуру, оператори > та == десь коректно означено) та клас вузла списку Node.

struct A\{

\quad  якісь коректні означення 

\quad  bool operator>(const A\&); 

\quad  bool operator==(const A\&); 

\quad A\& operator =(const A\&)=delete; 

\quad  A\& operator =(A\&\&)=delete;\};

struct Node\{Node *next; A a;\};

 Написати функцію, яка вилучає зі списку всі вузли, елемент в яких менший елемента в попередньому вузлі списку або дорівнює елементу в наступному вузлі. Список задається вказівником на перший вузол списку. Наприклад, зі списку, що містить послідовність 8, 5, 6, 9, 4, мають бути вилучені вузли, що містять 5 та 4. У цьому прикладі 6 залишається в списку.

Крім правильності та відповідності умові оцінюється якість коду, наявність витоків ресурсів та коректність роботи з пам'яттю. 

%enditem
}
{\smallskip\smallskip \begin {scriptsize} Затверджено на засіданні кафедри теоретичної кібернетики, протокол \No 4  від   25.01.2024 р.\par\smallskip\smallskip\smallskip\smallskip
{\parbox {6.5 cm} {Зав. кафедри \dotfill Ю.В.~Крак}}\hspace {0.7cm} {\hbox to 6.5 cm { Екзаменатор \dotfill Т.О.~Карнаух}} \end{scriptsize}}
%end
\newpage
%begin
{{\begin {scriptsize}\par
{ \center  Київський національний університет імені Тараса Шевченка \par}
{\parbox {5 cm}{Освітня програма \hfill}{\parbox {7 cm}{Інформатика\hfill}}\parbox {6 cm}{\hfill Семестр 2}}\par
{\parbox {5 cm} {Навчальний предмет \hfill}\parbox {7 cm}{Програмування \hfill}\parbox {6cm}{\hfill}\par}\vspace{-15pt} \end{scriptsize}}{\center Екзаменаційний білет \No \textbf { 38 }\par}}
{%beginitem
1. Вкажіть основну відмінність між відкритими (public) та прихованими (private) методами класу.%enditem
\par
\vspace{2mm}
%beginitem
2. Максимально змістовно повно сформулюйте назву оголошеного в наведеному фрагменті методу класу.
\vspace{-3mm}
\begin{verbatim}class A{
   ...
   A& operator =(A&&);
};\end{verbatim}
\vspace{-3mm}
%enditem
\par
\vspace{2mm}
%beginitem
3. Вкажіть, що не так з наведеним кодом.
\vspace{-3mm}
\begin{verbatim}int *pa=new int(2), *pb=pa; delete pb; *pa=3;\end{verbatim}
\vspace{-3mm}
%enditem
\par
\vspace{2mm}
%beginitem
4. За наявності означень
\vspace{-3mm}
\begin{verbatim}
class A{ };
class B:public A{ };
class C:public A{B b;};\end{verbatim}
\vspace{-3mm}
вкажіть послідовність викликів деструкторів за виконання коду \code{\{C c;\}} 

(Якщо деструктор викликається кілька разів поспіль, то має зазначатися кожен його виклик, наприклад, $\sim$A, $\sim$A, $\sim$A.)
%enditem

\vspace{2mm}
%beginitem
5. Яку часову оцінку складності має алгоритм швидкого сортування?
%enditem
\par
\vspace{2mm}
%beginitem
6. Модифікувати клас точок простору 

\qquad\code{class Point\{ double x, y, z;\};}


додавши до нього оператор перетворення на рядок типу \code{std::string}, що містить рядкове зображення точки простору (у круглих дужках координати через кому), наприклад, (1.01, 2.35, 0.5), та його означення.%enditem
\par
\vspace{2mm}
%beginitem
7. 
Реалізувати операцію перетину множин з повтореннями. Множина з повтореннями задається переліком своїх елементів, записаних за неспаданням у текстовому файлі через білі символи. Результат у такому ж самому форматі (і теж за неспаданням) має бути записаний у текстовий файл. Питома функція має отримувати три шляхи до файлів. Розмір вхідних множин такий, що вони не можуть бути завантажені повністю в пам’ять програми.

Крім правильності та відповідності умові оцінюється якість коду, наявність витоків ресурсів та коректність роботи з пам'яттю. 

%enditem
}
{\smallskip\smallskip \begin {scriptsize} Затверджено на засіданні кафедри теоретичної кібернетики, протокол \No 4  від   25.01.2024 р.\par\smallskip\smallskip\smallskip\smallskip
{\parbox {6.5 cm} {Зав. кафедри \dotfill Ю.В.~Крак}}\hspace {0.7cm} {\hbox to 6.5 cm { Екзаменатор \dotfill Т.О.~Карнаух}} \end{scriptsize}}
%end
\newpage
%begin
{{\begin {scriptsize}\par
{ \center  Київський національний університет імені Тараса Шевченка \par}
{\parbox {5 cm}{Освітня програма \hfill}{\parbox {7 cm}{Інформатика\hfill}}\parbox {6 cm}{\hfill Семестр 2}}\par
{\parbox {5 cm} {Навчальний предмет \hfill}\parbox {7 cm}{Програмування \hfill}\parbox {6cm}{\hfill}\par}\vspace{-15pt} \end{scriptsize}}{\center Екзаменаційний білет \No \textbf { 39 }\par}}
{%beginitem
1. Що важливо забезпечити при роботі з динамічними змінними?%enditem
\par
\vspace{2mm}
%beginitem
2. Максимально змістовно повно сформулюйте назву оголошеного в наведеному фрагменті методу класу.
\vspace{-3mm}
\begin{verbatim}class A{
   ...
   A(A&&);
};\end{verbatim}
\vspace{-3mm}
%enditem
\par
\vspace{2mm}
%beginitem
3. Вкажіть, що не так з наведеним кодом.
\vspace{-3mm}
\begin{verbatim}int *pa=new int(2), *pb=new int(3); pa=pb;\end{verbatim}
\vspace{-3mm}
%enditem
\par
\vspace{2mm}
%beginitem
4. За наявності означень
\vspace{-3mm}
\begin{verbatim}
class A{ };
class B{ };
class C:public A{B b[2];};\end{verbatim}
\vspace{-3mm}
вкажіть послідовність викликів конструкторів за виконання коду \code{\{C c;\}} 

(Якщо конструктор викликається кілька разів поспіль, то має зазначатися кожен його виклик, наприклад, A, A, A.)
%enditem

\vspace{2mm}
%beginitem
5. Яку часову оцінку складності має алгоритм сортування бульбашкою?
%enditem
\par
\vspace{2mm}
%beginitem
6. Модифікувати клас 

\qquad\code{class Session\{ int DM, MA, Alg, Prog;\};} 


додавши до нього оператор перетворення на тип \code{bool}: об'єкт класу є істинним тоді і тільки тоді, коли значення всіх атрибутів не менші за 60. Означення оператора теж має бути записане.%enditem
\par
\vspace{2mm}
%beginitem
7. 
Задано клас A (структуру, оператори < та == десь коректно означено) та клас вузла списку Node.

struct A\{

\quad  якісь коректні означення 

\quad  bool operator<(const A\&); 

\quad A\& operator =(const A\&)=delete; 

\quad  A\& operator =(A\&\&)=delete;\};

struct Node\{Node *prev; Node *next; A a;\};

Написати функцію, яка вилучає зі списку всі вузли, елемент в яких менший елемента в попередньому вузлі списку або дорівнює елементу в наступному вузлі. Список задається вказівником на перший вузол списку. Наприклад, зі списку, що містить послідовність 8, 5, 6, 9, 4, мають бути вилучені вузли, що містять 5 та 4. У цьому прикладі 6 залишається в списку.

Крім правильності та відповідності умові оцінюється якість коду, наявність витоків ресурсів та коректність роботи з пам'яттю. 

%enditem
}
{\smallskip\smallskip \begin {scriptsize} Затверджено на засіданні кафедри теоретичної кібернетики, протокол \No 4  від   25.01.2024 р.\par\smallskip\smallskip\smallskip\smallskip
{\parbox {6.5 cm} {Зав. кафедри \dotfill Ю.В.~Крак}}\hspace {0.7cm} {\hbox to 6.5 cm { Екзаменатор \dotfill Т.О.~Карнаух}} \end{scriptsize}}
%end
\newpage
%begin
{{\begin {scriptsize}\par
{ \center  Київський національний університет імені Тараса Шевченка \par}
{\parbox {5 cm}{Освітня програма \hfill}{\parbox {7 cm}{Інформатика\hfill}}\parbox {6 cm}{\hfill Семестр 2}}\par
{\parbox {5 cm} {Навчальний предмет \hfill}\parbox {7 cm}{Програмування \hfill}\parbox {6cm}{\hfill}\par}\vspace{-15pt} \end{scriptsize}}{\center Екзаменаційний білет \No \textbf { 40 }\par}}
{%beginitem
1. Вкажіть основну відмінність між віртуальними методами та тими, що не є віртуальними.%enditem
\par
\vspace{2mm}
%beginitem
2. Максимально змістовно повно сформулюйте назву оголошеного в наведеному фрагменті методу класу.
\vspace{-3mm}
\begin{verbatim}class A{
   ...
   A();
};\end{verbatim}
\vspace{-3mm}
%enditem
\par
\vspace{2mm}
%beginitem
3. Вкажіть, що не так з наведеним кодом.
\vspace{-3mm}
\begin{verbatim}int *pa=new int(2), *pb=&pa; delete pa;\end{verbatim}
\vspace{-3mm}
%enditem
\par
\vspace{2mm}
%beginitem
4. За наявності означень
\vspace{-3mm}
\begin{verbatim}
class A{ };
class B{ };
class C:public A{B b;};\end{verbatim}
\vspace{-3mm}
вкажіть послідовність викликів деструкторів за виконання коду \code{\{C c;\}} 

(Якщо деструктор викликається кілька разів поспіль, то має зазначатися кожен його виклик, наприклад, $\sim$A, $\sim$A, $\sim$A.)
%enditem

\vspace{2mm}
%beginitem
5. Яку часову оцінку складності має алгоритм швидкого сортування?
%enditem
\par
\vspace{2mm}
%beginitem
6. Модифікувати клас точок простору 

\qquad\code{class Point\{ double x, y, z;\};}


додавши до нього оператор перетворення на тип \code{bool}: точка є істинною тоді і тільки тоді, коли містить хоча б одну ненульову координату. Означення оператора теж має бути записане.%enditem
\par
\vspace{2mm}
%beginitem
7. 
Реалізувати операцію симетричної різниці множин з повтореннями. Множина з повтореннями задається переліком своїх елементів, записаних за неспаданням у текстовому файлі через білі символи. Результат у такому ж самому форматі (і теж за неспаданням) має бути записаний у текстовий файл. Питома функція має отримувати три шляхи до файлів. Розмір вхідних множин такий, що вони не можуть бути завантажені повністю в пам’ять програми.

Крім правильності та відповідності умові оцінюється якість коду, наявність витоків ресурсів та коректність роботи з пам'яттю. 

%enditem
}
{\smallskip\smallskip \begin {scriptsize} Затверджено на засіданні кафедри теоретичної кібернетики, протокол \No 4  від   25.01.2024 р.\par\smallskip\smallskip\smallskip\smallskip
{\parbox {6.5 cm} {Зав. кафедри \dotfill Ю.В.~Крак}}\hspace {0.7cm} {\hbox to 6.5 cm { Екзаменатор \dotfill Т.О.~Карнаух}} \end{scriptsize}}
%end
\newpage
%begin
{{\begin {scriptsize}\par
{ \center  Київський національний університет імені Тараса Шевченка \par}
{\parbox {5 cm}{Освітня програма \hfill}{\parbox {7 cm}{Інформатика\hfill}}\parbox {6 cm}{\hfill Семестр 2}}\par
{\parbox {5 cm} {Навчальний предмет \hfill}\parbox {7 cm}{Програмування \hfill}\parbox {6cm}{\hfill}\par}\vspace{-15pt} \end{scriptsize}}{\center Екзаменаційний білет \No \textbf { 41 }\par}}
{%beginitem
1. Вкажіть основну \textbf{синтаксичну} відмінність між конструктором копії та конструктором переміщення.%enditem
\par
\vspace{2mm}
%beginitem
2. Максимально змістовно повно сформулюйте назву оголошеного в наведеному фрагменті методу класу.
\vspace{-3mm}
\begin{verbatim}class A{
   ...
   A(const A&);
};\end{verbatim}
\vspace{-3mm}
%enditem
\par
\vspace{2mm}
%beginitem
3. Вкажіть, що не так з наведеним кодом.
\vspace{-3mm}
\begin{verbatim}int *pa=new int(2), *pb=*pa; delete pa;\end{verbatim}
\vspace{-3mm}
%enditem
\par
\vspace{2mm}
%beginitem
4. За наявності означень
\vspace{-3mm}
\begin{verbatim}
class A{ };
class B:public A{ };
class C:public A{B b;};\end{verbatim}
\vspace{-3mm}
вкажіть послідовність викликів конструкторів за виконання коду \code{\{C c;\}} 

(Якщо конструктор викликається кілька разів поспіль, то має зазначатися кожен його виклик, наприклад, A, A, A.)
%enditem

\vspace{2mm}
%beginitem
5. Яку часову оцінку складності має алгоритм сортування вибором?
%enditem
\par
\vspace{2mm}
%beginitem
6. Модифікувати клас точок простору 

\qquad\code{class Point\{ double x, y, z;\};}


додавши до нього оператор перетворення на рядок типу \code{std::string}, що містить рядкове зображення точки простору (у круглих дужках координати через кому), наприклад, (1.01, 2.35, 0.5), та його означення.%enditem
\par
\vspace{2mm}
%beginitem
7. 
Реалізувати операцію різниці множин з повтореннями. Множина з повтореннями задається переліком своїх елементів, записаних за неспаданням у текстовому файлі через білі символи. Результат у такому ж самому форматі (і теж за неспаданням) має бути записаний у текстовий файл. Питома функція має отримувати три шляхи до файлів. Розмір вхідних множин такий, що вони не можуть бути завантажені повністю в пам’ять програми.

Крім правильності та відповідності умові оцінюється якість коду, наявність витоків ресурсів та коректність роботи з пам'яттю. 

%enditem
}
{\smallskip\smallskip \begin {scriptsize} Затверджено на засіданні кафедри теоретичної кібернетики, протокол \No 4  від   25.01.2024 р.\par\smallskip\smallskip\smallskip\smallskip
{\parbox {6.5 cm} {Зав. кафедри \dotfill Ю.В.~Крак}}\hspace {0.7cm} {\hbox to 6.5 cm { Екзаменатор \dotfill Т.О.~Карнаух}} \end{scriptsize}}
%end
\newpage
%begin
{{\begin {scriptsize}\par
{ \center  Київський національний університет імені Тараса Шевченка \par}
{\parbox {5 cm}{Освітня програма \hfill}{\parbox {7 cm}{Інформатика\hfill}}\parbox {6 cm}{\hfill Семестр 2}}\par
{\parbox {5 cm} {Навчальний предмет \hfill}\parbox {7 cm}{Програмування \hfill}\parbox {6cm}{\hfill}\par}\vspace{-15pt} \end{scriptsize}}{\center Екзаменаційний білет \No \textbf { 42 }\par}}
{%beginitem
1. Вкажіть основну \textbf{семантичну} відмінність між конструктором копії та конструктором переміщення.%enditem
\par
\vspace{2mm}
%beginitem
2. Максимально змістовно повно сформулюйте назву оголошеного в наведеному фрагменті методу класу.
\vspace{-3mm}
\begin{verbatim}class A{
   ...
   A();
};\end{verbatim}
\vspace{-3mm}
%enditem
\par
\vspace{2mm}
%beginitem
3. Вкажіть, що не так з наведеним кодом.
\vspace{-3mm}
\begin{verbatim}int **ppi=new int*(new int(3)); delete ppi;\end{verbatim}
\vspace{-3mm}
%enditem
\par
\vspace{2mm}
%beginitem
4. За наявності означень
\vspace{-3mm}
\begin{verbatim}
class A{ };
class B{ };
class C:public A{B b[2];};\end{verbatim}
\vspace{-3mm}
вкажіть послідовність викликів деструкторів за виконання коду \code{\{C c;\}} 

(Якщо деструктор викликається кілька разів поспіль, то має зазначатися кожен його виклик, наприклад, $\sim$A, $\sim$A, $\sim$A.)
%enditem

\vspace{2mm}
%beginitem
5. Яку часову оцінку складності має алгоритм сортування вставками?
%enditem
\par
\vspace{2mm}
%beginitem
6. Модифікувати клас 

\qquad\code{class Session\{ int DM, MA, Alg, Prog;\};} 


додавши до нього оператор перетворення на тип \code{int}: якщо хоча б один атрибут менший за 60, то має повертатись 0, інакше -- сума атрибутів.%enditem
\par
\vspace{2mm}
%beginitem
7. 
Задано клас A (структуру, оператори < та == десь коректно означено) та клас вузла списку Node.

struct A\{

\quad  якісь коректні означення 

\quad  bool operator<(const A\&); 

\quad  bool operator==(const A\&); 

\quad A\& operator =(const A\&)=delete; 

\quad  A\& operator =(A\&\&)=delete;\};

struct Node\{Node *next; A a;\};

 Написати функцію, яка вилучає зі списку всі вузли, елемент в яких більший елемента в попередньому вузлі списку або дорівнює елементу в наступному вузлі. Список є циклічним та задається вказівником на перший вузол списку. Наприклад, зі списку, що містить послідовність 8, 10, 9, 4, мають бути вилучені вузли, що містять 8 (бо список циклічний і $8>4$) та 10. У цьому прикладі 9 залишається в списку.

Крім правильності та відповідності умові оцінюється якість коду, наявність витоків ресурсів та коректність роботи з пам'яттю. 

%enditem
}
{\smallskip\smallskip \begin {scriptsize} Затверджено на засіданні кафедри теоретичної кібернетики, протокол \No 4  від   25.01.2024 р.\par\smallskip\smallskip\smallskip\smallskip
{\parbox {6.5 cm} {Зав. кафедри \dotfill Ю.В.~Крак}}\hspace {0.7cm} {\hbox to 6.5 cm { Екзаменатор \dotfill Т.О.~Карнаух}} \end{scriptsize}}
%end
\newpage
%begin
{{\begin {scriptsize}\par
{ \center  Київський національний університет імені Тараса Шевченка \par}
{\parbox {5 cm}{Освітня програма \hfill}{\parbox {7 cm}{Інформатика\hfill}}\parbox {6 cm}{\hfill Семестр 2}}\par
{\parbox {5 cm} {Навчальний предмет \hfill}\parbox {7 cm}{Програмування \hfill}\parbox {6cm}{\hfill}\par}\vspace{-15pt} \end{scriptsize}}{\center Екзаменаційний білет \No \textbf { 43 }\par}}
{%beginitem
1. Вкажіть основну відмінність між відкритими (public) та прихованими (private) методами класу.%enditem
\par
\vspace{2mm}
%beginitem
2. Максимально змістовно повно сформулюйте назву оголошеного в наведеному фрагменті методу класу.
\vspace{-3mm}
\begin{verbatim}class A{
   ...
   A& operator =(A&&);
};\end{verbatim}
\vspace{-3mm}
%enditem
\par
\vspace{2mm}
%beginitem
3. Вкажіть, що не так з наведеним кодом.
\vspace{-3mm}
\begin{verbatim}int a, *p=&a; delete a;\end{verbatim}
\vspace{-3mm}
%enditem
\par
\vspace{2mm}
%beginitem
4. За наявності означень
\vspace{-3mm}
\begin{verbatim}
class A{ };
class B:public A{ };
class C:public B{ };\end{verbatim}
\vspace{-3mm}
вкажіть послідовність викликів деструкторів за виконання коду \code{\{C c;\}} 

(Якщо деструктор викликається кілька разів поспіль, то має зазначатися кожен його виклик, наприклад, $\sim$A, $\sim$A, $\sim$A.)
%enditem

\vspace{2mm}
%beginitem
5. Яку часову оцінку складності має алгоритм швидкого сортування?
%enditem
\par
\vspace{2mm}
%beginitem
6. Модифікувати клас точок простору 

\qquad\code{class Point\{ double x, y, z;\};}


додавши до нього оператор перетворення на тип \code{bool}: точка є істинною тоді і тільки тоді, коли містить хоча б одну ненульову координату. Означення оператора теж має бути записане.%enditem
\par
\vspace{2mm}
%beginitem
7. 
Задано клас A (структуру, оператори < та == десь коректно означено) та клас вузла списку Node.

struct A\{

\quad  якісь коректні означення 

\quad  bool operator<(const A\&); 

\quad  bool operator==(const A\&); 

\quad A\& operator =(const A\&)=delete; 

\quad  A\& operator =(A\&\&)=delete;\};

struct Node\{Node *prev; Node *next; A a;\};

Написати функцію, яка вилучає зі списку всі вузли, елемент в яких більший елемента в попередньому вузлі списку або дорівнює елементу в наступному вузлі. Список є циклічним та задається вказівником на перший вузол списку. Наприклад, зі списку, що містить послідовність 8, 10, 9, 4, мають бути вилучені вузли, що містять 8 (бо список циклічний і $8>4$) та 10. У цьому прикладі 9 залишається в списку.

Крім правильності та відповідності умові оцінюється якість коду, наявність витоків ресурсів та коректність роботи з пам'яттю. 

%enditem
}
{\smallskip\smallskip \begin {scriptsize} Затверджено на засіданні кафедри теоретичної кібернетики, протокол \No 4  від   25.01.2024 р.\par\smallskip\smallskip\smallskip\smallskip
{\parbox {6.5 cm} {Зав. кафедри \dotfill Ю.В.~Крак}}\hspace {0.7cm} {\hbox to 6.5 cm { Екзаменатор \dotfill Т.О.~Карнаух}} \end{scriptsize}}
%end
\newpage
%begin
{{\begin {scriptsize}\par
{ \center  Київський національний університет імені Тараса Шевченка \par}
{\parbox {5 cm}{Освітня програма \hfill}{\parbox {7 cm}{Інформатика\hfill}}\parbox {6 cm}{\hfill Семестр 2}}\par
{\parbox {5 cm} {Навчальний предмет \hfill}\parbox {7 cm}{Програмування \hfill}\parbox {6cm}{\hfill}\par}\vspace{-15pt} \end{scriptsize}}{\center Екзаменаційний білет \No \textbf { 44 }\par}}
{%beginitem
1. Вкажіть основну відмінність між класами (class) та структурами (struct).%enditem
\par
\vspace{2mm}
%beginitem
2. Максимально змістовно повно сформулюйте назву оголошеного в наведеному фрагменті методу класу.
\vspace{-3mm}
\begin{verbatim}class A{
   ...
   A(A&&);
};\end{verbatim}
\vspace{-3mm}
%enditem
\par
\vspace{2mm}
%beginitem
3. Вкажіть, що не так з наведеним кодом.
\vspace{-3mm}
\begin{verbatim}int *pa=new int(2), *pb=&pa; delete pa;\end{verbatim}
\vspace{-3mm}
%enditem
\par
\vspace{2mm}
%beginitem
4. За наявності означень
\vspace{-3mm}
\begin{verbatim}
class A{ };
class B{ };
class C:public A{B b;A a;};\end{verbatim}
\vspace{-3mm}
вкажіть послідовність викликів деструкторів за виконання коду \code{\{C c;\}} 

(Якщо деструктор викликається кілька разів поспіль, то має зазначатися кожен його виклик, наприклад, $\sim$A, $\sim$A, $\sim$A.)
%enditem

\vspace{2mm}
%beginitem
5. Яку часову оцінку складності має алгоритм сортування злиттям?
%enditem
\par
\vspace{2mm}
%beginitem
6. Модифікувати клас 

\qquad\code{class Session\{ int DM, MA, Alg, Prog;\};} 


додавши до нього оператор перетворення на тип \code{bool}: об'єкт класу є істинним тоді і тільки тоді, коли значення всіх атрибутів не менші за 60. Означення оператора теж має бути записане.%enditem
\par
\vspace{2mm}
%beginitem
7. 
Реалізувати операцію перетину множин з повтореннями. Множина з повтореннями задається переліком своїх елементів, записаних за неспаданням у текстовому файлі через білі символи. Результат у такому ж самому форматі (і теж за неспаданням) має бути записаний у текстовий файл. Питома функція має отримувати три шляхи до файлів. Розмір вхідних множин такий, що вони не можуть бути завантажені повністю в пам’ять програми.

Крім правильності та відповідності умові оцінюється якість коду, наявність витоків ресурсів та коректність роботи з пам'яттю. 

%enditem
}
{\smallskip\smallskip \begin {scriptsize} Затверджено на засіданні кафедри теоретичної кібернетики, протокол \No 4  від   25.01.2024 р.\par\smallskip\smallskip\smallskip\smallskip
{\parbox {6.5 cm} {Зав. кафедри \dotfill Ю.В.~Крак}}\hspace {0.7cm} {\hbox to 6.5 cm { Екзаменатор \dotfill Т.О.~Карнаух}} \end{scriptsize}}
%end
\newpage
%begin
{{\begin {scriptsize}\par
{ \center  Київський національний університет імені Тараса Шевченка \par}
{\parbox {5 cm}{Освітня програма \hfill}{\parbox {7 cm}{Інформатика\hfill}}\parbox {6 cm}{\hfill Семестр 2}}\par
{\parbox {5 cm} {Навчальний предмет \hfill}\parbox {7 cm}{Програмування \hfill}\parbox {6cm}{\hfill}\par}\vspace{-15pt} \end{scriptsize}}{\center Екзаменаційний білет \No \textbf { 45 }\par}}
{%beginitem
1. Що важливо забезпечити при роботі з динамічними змінними?%enditem
\par
\vspace{2mm}
%beginitem
2. Максимально змістовно повно сформулюйте назву оголошеного в наведеному фрагменті методу класу.
\vspace{-3mm}
\begin{verbatim}class A{
   ...
   ~A();
};\end{verbatim}
\vspace{-3mm}
%enditem
\par
\vspace{2mm}
%beginitem
3. Вкажіть, що не так з наведеним кодом.
\vspace{-3mm}
\begin{verbatim}int *pa=new int(2), *pb=new int(3); pa=pb;\end{verbatim}
\vspace{-3mm}
%enditem
\par
\vspace{2mm}
%beginitem
4. За наявності означень
\vspace{-3mm}
\begin{verbatim}
class A{ };
class B{ };
class C:public A{B b;};\end{verbatim}
\vspace{-3mm}
вкажіть послідовність викликів конструкторів за виконання коду \code{\{C c;\}} 

(Якщо конструктор викликається кілька разів поспіль, то має зазначатися кожен його виклик, наприклад, A, A, A.)
%enditem

\vspace{2mm}
%beginitem
5. Яку часову оцінку складності має алгоритм сортування бульбашкою?
%enditem
\par
\vspace{2mm}
%beginitem
6. Модифікувати клас 

\qquad\code{class Session\{ int DM, MA, Alg, Prog;\};} 


додавши до нього оператор порівняння \code{<}: один об'єкт є меншим за інший тоді і тільки тоді, коли значення всіх атрибутів об'єктів не менші за 60 і сума атрибутів першого менша за суму атрибутів другого. Означення оператора теж має бути записане.%enditem
\par
\vspace{2mm}
%beginitem
7. 
Задано клас A (структуру, оператори > та == десь коректно означено) та клас вузла списку Node.

struct A\{

\quad  якісь коректні означення 

\quad  bool operator>(const A\&); 

\quad  bool operator==(const A\&); 

\quad A\& operator =(const A\&)=delete; 

\quad  A\& operator =(A\&\&)=delete;\};

struct Node\{Node *next; A a;\};

 Написати функцію, яка вилучає зі списку всі вузли, елемент в яких менший елемента в попередньому вузлі списку або дорівнює елементу в наступному вузлі. Список задається вказівником на перший вузол списку. Наприклад, зі списку, що містить послідовність 8, 5, 6, 9, 4, мають бути вилучені вузли, що містять 5 та 4. У цьому прикладі 6 залишається в списку.

Крім правильності та відповідності умові оцінюється якість коду, наявність витоків ресурсів та коректність роботи з пам'яттю. 

%enditem
}
{\smallskip\smallskip \begin {scriptsize} Затверджено на засіданні кафедри теоретичної кібернетики, протокол \No 4  від   25.01.2024 р.\par\smallskip\smallskip\smallskip\smallskip
{\parbox {6.5 cm} {Зав. кафедри \dotfill Ю.В.~Крак}}\hspace {0.7cm} {\hbox to 6.5 cm { Екзаменатор \dotfill Т.О.~Карнаух}} \end{scriptsize}}
%end
\newpage
%begin
{{\begin {scriptsize}\par
{ \center  Київський національний університет імені Тараса Шевченка \par}
{\parbox {5 cm}{Освітня програма \hfill}{\parbox {7 cm}{Інформатика\hfill}}\parbox {6 cm}{\hfill Семестр 2}}\par
{\parbox {5 cm} {Навчальний предмет \hfill}\parbox {7 cm}{Програмування \hfill}\parbox {6cm}{\hfill}\par}\vspace{-15pt} \end{scriptsize}}{\center Екзаменаційний білет \No \textbf { 46 }\par}}
{%beginitem
1. Вкажіть основну \textbf{семантичну} відмінність між конструктором копії та конструктором переміщення.%enditem
\par
\vspace{2mm}
%beginitem
2. Максимально змістовно повно сформулюйте назву оголошеного в наведеному фрагменті методу класу.
\vspace{-3mm}
\begin{verbatim}class A{
   ...
   A& operator =(A&&);
};\end{verbatim}
\vspace{-3mm}
%enditem
\par
\vspace{2mm}
%beginitem
3. Вкажіть, що не так з наведеним кодом.
\vspace{-3mm}
\begin{verbatim}int *pa=new int(2), *pb=pa; delete pb; delete pa;\end{verbatim}
\vspace{-3mm}
%enditem
\par
\vspace{2mm}
%beginitem
4. За наявності означень
\vspace{-3mm}
\begin{verbatim}
class A{ };
class B{ };
class C{B b;A a;};\end{verbatim}
\vspace{-3mm}
вкажіть послідовність викликів конструкторів за виконання коду \code{\{C c;\}} 

(Якщо конструктор викликається кілька разів поспіль, то має зазначатися кожен його виклик, наприклад, A, A, A.)
%enditem

\vspace{2mm}
%beginitem
5. Яку часову оцінку складності має алгоритм швидкого сортування?
%enditem
\par
\vspace{2mm}
%beginitem
6. Модифікувати клас 

\qquad\code{class Session\{ int DM, MA, Alg, Prog;\};} 


додавши до нього оператор перетворення на тип \code{int}: якщо хоча б один атрибут менший за 60, то має повертатись 0, інакше -- сума атрибутів.%enditem
\par
\vspace{2mm}
%beginitem
7. 
Задано клас A (структуру, оператори < та == десь коректно означено) та клас вузла списку Node.

struct A\{

\quad  якісь коректні означення 

\quad  bool operator<(const A\&); 

\quad A\& operator =(const A\&)=delete; 

\quad  A\& operator =(A\&\&)=delete;\};

struct Node\{Node *prev; Node *next; A a;\};

Написати функцію, яка вилучає зі списку всі вузли, елемент в яких менший елемента в попередньому вузлі списку або дорівнює елементу в наступному вузлі. Список задається вказівником на перший вузол списку. Наприклад, зі списку, що містить послідовність 8, 5, 6, 9, 4, мають бути вилучені вузли, що містять 5 та 4. У цьому прикладі 6 залишається в списку.

Крім правильності та відповідності умові оцінюється якість коду, наявність витоків ресурсів та коректність роботи з пам'яттю. 

%enditem
}
{\smallskip\smallskip \begin {scriptsize} Затверджено на засіданні кафедри теоретичної кібернетики, протокол \No 4  від   25.01.2024 р.\par\smallskip\smallskip\smallskip\smallskip
{\parbox {6.5 cm} {Зав. кафедри \dotfill Ю.В.~Крак}}\hspace {0.7cm} {\hbox to 6.5 cm { Екзаменатор \dotfill Т.О.~Карнаух}} \end{scriptsize}}
%end
\newpage
%begin
{{\begin {scriptsize}\par
{ \center  Київський національний університет імені Тараса Шевченка \par}
{\parbox {5 cm}{Освітня програма \hfill}{\parbox {7 cm}{Інформатика\hfill}}\parbox {6 cm}{\hfill Семестр 2}}\par
{\parbox {5 cm} {Навчальний предмет \hfill}\parbox {7 cm}{Програмування \hfill}\parbox {6cm}{\hfill}\par}\vspace{-15pt} \end{scriptsize}}{\center Екзаменаційний білет \No \textbf { 47 }\par}}
{%beginitem
1. Вкажіть основну відмінність між віртуальними методами та тими, що не є віртуальними.%enditem
\par
\vspace{2mm}
%beginitem
2. Максимально змістовно повно сформулюйте назву оголошеного в наведеному фрагменті методу класу.
\vspace{-3mm}
\begin{verbatim}class A{
   ...
   A();
};\end{verbatim}
\vspace{-3mm}
%enditem
\par
\vspace{2mm}
%beginitem
3. Вкажіть, що не так з наведеним кодом.
\vspace{-3mm}
\begin{verbatim}int **ppi=new int*(new int(3)); delete ppi;\end{verbatim}
\vspace{-3mm}
%enditem
\par
\vspace{2mm}
%beginitem
4. За наявності означень
\vspace{-3mm}
\begin{verbatim}
class A{ };
class B:public A{ };
class C:public B{ };\end{verbatim}
\vspace{-3mm}
вкажіть послідовність викликів конструкторів за виконання коду \code{\{C c;\}} 

(Якщо конструктор викликається кілька разів поспіль, то має зазначатися кожен його виклик, наприклад, A, A, A.)
%enditem

\vspace{2mm}
%beginitem
5. Яку часову оцінку складності має алгоритм сортування вибором?
%enditem
\par
\vspace{2mm}
%beginitem
6. Модифікувати клас точок простору 

\qquad\code{class Point\{ double x, y, z;\};}


додавши до нього оператор перетворення на тип \code{bool}: точка є істинною тоді і тільки тоді, коли містить хоча б одну ненульову координату. Означення оператора теж має бути записане.%enditem
\par
\vspace{2mm}
%beginitem
7. 
Реалізувати операцію симетричної різниці множин з повтореннями. Множина з повтореннями задається переліком своїх елементів, записаних за неспаданням у текстовому файлі через білі символи. Результат у такому ж самому форматі (і теж за неспаданням) має бути записаний у текстовий файл. Питома функція має отримувати три шляхи до файлів. Розмір вхідних множин такий, що вони не можуть бути завантажені повністю в пам’ять програми.

Крім правильності та відповідності умові оцінюється якість коду, наявність витоків ресурсів та коректність роботи з пам'яттю. 

%enditem
}
{\smallskip\smallskip \begin {scriptsize} Затверджено на засіданні кафедри теоретичної кібернетики, протокол \No 4  від   25.01.2024 р.\par\smallskip\smallskip\smallskip\smallskip
{\parbox {6.5 cm} {Зав. кафедри \dotfill Ю.В.~Крак}}\hspace {0.7cm} {\hbox to 6.5 cm { Екзаменатор \dotfill Т.О.~Карнаух}} \end{scriptsize}}
%end
\newpage
%begin
{{\begin {scriptsize}\par
{ \center  Київський національний університет імені Тараса Шевченка \par}
{\parbox {5 cm}{Освітня програма \hfill}{\parbox {7 cm}{Інформатика\hfill}}\parbox {6 cm}{\hfill Семестр 2}}\par
{\parbox {5 cm} {Навчальний предмет \hfill}\parbox {7 cm}{Програмування \hfill}\parbox {6cm}{\hfill}\par}\vspace{-15pt} \end{scriptsize}}{\center Екзаменаційний білет \No \textbf { 48 }\par}}
{%beginitem
1. Вкажіть основну \textbf{синтаксичну} відмінність між конструктором копії та конструктором переміщення.%enditem
\par
\vspace{2mm}
%beginitem
2. Максимально змістовно повно сформулюйте назву оголошеного в наведеному фрагменті методу класу.
\vspace{-3mm}
\begin{verbatim}class A{
   ...
   A(const A&);
};\end{verbatim}
\vspace{-3mm}
%enditem
\par
\vspace{2mm}
%beginitem
3. Вкажіть, що не так з наведеним кодом.
\vspace{-3mm}
\begin{verbatim}int *pa=new int(2), *pb=pa; delete pb; *pa=3;\end{verbatim}
\vspace{-3mm}
%enditem
\par
\vspace{2mm}
%beginitem
4. За наявності означень
\vspace{-3mm}
\begin{verbatim}
class A{ };
class B{ };
class C:public A{B b;A a;};\end{verbatim}
\vspace{-3mm}
вкажіть послідовність викликів конструкторів за виконання коду \code{\{C c;\}} 

(Якщо конструктор викликається кілька разів поспіль, то має зазначатися кожен його виклик, наприклад, A, A, A.)
%enditem

\vspace{2mm}
%beginitem
5. Яку часову оцінку складності має алгоритм сортування вставками?
%enditem
\par
\vspace{2mm}
%beginitem
6. Модифікувати клас 

\qquad\code{class Session\{ int DM, MA, Alg, Prog;\};} 


додавши до нього оператор порівняння \code{<}: один об'єкт є меншим за інший тоді і тільки тоді, коли значення всіх атрибутів об'єктів не менші за 60 і сума атрибутів першого менша за суму атрибутів другого. Означення оператора теж має бути записане.%enditem
\par
\vspace{2mm}
%beginitem
7. 
Реалізувати операцію злиття множин з повтореннями. Множина з повтореннями задається переліком своїх елементів, записаних за неспаданням у текстовому файлі через білі символи. Результат у такому ж самому форматі (і теж за неспаданням) має бути записаний у текстовий файл. Питома функція має отримувати три шляхи до файлів. Розмір вхідних множин такий, що вони не можуть бути завантажені повністю в пам’ять програми.

Крім правильності та відповідності умові оцінюється якість коду, наявність витоків ресурсів та коректність роботи з пам'яттю. 

%enditem
}
{\smallskip\smallskip \begin {scriptsize} Затверджено на засіданні кафедри теоретичної кібернетики, протокол \No 4  від   25.01.2024 р.\par\smallskip\smallskip\smallskip\smallskip
{\parbox {6.5 cm} {Зав. кафедри \dotfill Ю.В.~Крак}}\hspace {0.7cm} {\hbox to 6.5 cm { Екзаменатор \dotfill Т.О.~Карнаух}} \end{scriptsize}}
%end
\newpage
%begin
{{\begin {scriptsize}\par
{ \center  Київський національний університет імені Тараса Шевченка \par}
{\parbox {5 cm}{Освітня програма \hfill}{\parbox {7 cm}{Інформатика\hfill}}\parbox {6 cm}{\hfill Семестр 2}}\par
{\parbox {5 cm} {Навчальний предмет \hfill}\parbox {7 cm}{Програмування \hfill}\parbox {6cm}{\hfill}\par}\vspace{-15pt} \end{scriptsize}}{\center Екзаменаційний білет \No \textbf { 49 }\par}}
{%beginitem
1. Вкажіть основну відмінність між класами (class) та структурами (struct).%enditem
\par
\vspace{2mm}
%beginitem
2. Максимально змістовно повно сформулюйте назву оголошеного в наведеному фрагменті методу класу.
\vspace{-3mm}
\begin{verbatim}class A{
   ...
   A(A&&);
};\end{verbatim}
\vspace{-3mm}
%enditem
\par
\vspace{2mm}
%beginitem
3. Вкажіть, що не так з наведеним кодом.
\vspace{-3mm}
\begin{verbatim}int *pa=new int(2), *pb=*pa; delete pa;\end{verbatim}
\vspace{-3mm}
%enditem
\par
\vspace{2mm}
%beginitem
4. За наявності означень
\vspace{-3mm}
\begin{verbatim}
class A{ };
class B{ };
class C:public A{B b;};\end{verbatim}
\vspace{-3mm}
вкажіть послідовність викликів конструкторів за виконання коду \code{\{C c;\}} 

(Якщо конструктор викликається кілька разів поспіль, то має зазначатися кожен його виклик, наприклад, A, A, A.)
%enditem

\vspace{2mm}
%beginitem
5. Яку часову оцінку складності має алгоритм сортування купою (він же деревом, пірамідою)?
%enditem
\par
\vspace{2mm}
%beginitem
6. Модифікувати клас точок простору 

\qquad\code{class Point\{ double x, y, z;\};}


додавши до нього оператор перетворення на рядок типу \code{std::string}, що містить рядкове зображення точки простору (у круглих дужках координати через кому), наприклад, (1.01, 2.35, 0.5), та його означення.%enditem
\par
\vspace{2mm}
%beginitem
7. 
Реалізувати операцію різниці множин з повтореннями. Множина з повтореннями задається переліком своїх елементів, записаних за неспаданням у текстовому файлі через білі символи. Результат у такому ж самому форматі (і теж за неспаданням) має бути записаний у текстовий файл. Питома функція має отримувати три шляхи до файлів. Розмір вхідних множин такий, що вони не можуть бути завантажені повністю в пам’ять програми.

Крім правильності та відповідності умові оцінюється якість коду, наявність витоків ресурсів та коректність роботи з пам'яттю. 

%enditem
}
{\smallskip\smallskip \begin {scriptsize} Затверджено на засіданні кафедри теоретичної кібернетики, протокол \No 4  від   25.01.2024 р.\par\smallskip\smallskip\smallskip\smallskip
{\parbox {6.5 cm} {Зав. кафедри \dotfill Ю.В.~Крак}}\hspace {0.7cm} {\hbox to 6.5 cm { Екзаменатор \dotfill Т.О.~Карнаух}} \end{scriptsize}}
%end
\newpage
%begin
{{\begin {scriptsize}\par
{ \center  Київський національний університет імені Тараса Шевченка \par}
{\parbox {5 cm}{Освітня програма \hfill}{\parbox {7 cm}{Інформатика\hfill}}\parbox {6 cm}{\hfill Семестр 2}}\par
{\parbox {5 cm} {Навчальний предмет \hfill}\parbox {7 cm}{Програмування \hfill}\parbox {6cm}{\hfill}\par}\vspace{-15pt} \end{scriptsize}}{\center Екзаменаційний білет \No \textbf { 50 }\par}}
{%beginitem
1. Вкажіть основну відмінність між відкритими (public) та прихованими (private) методами класу.%enditem
\par
\vspace{2mm}
%beginitem
2. Максимально змістовно повно сформулюйте назву оголошеного в наведеному фрагменті методу класу.
\vspace{-3mm}
\begin{verbatim}class A{
   ...
   ~A();
};\end{verbatim}
\vspace{-3mm}
%enditem
\par
\vspace{2mm}
%beginitem
3. Вкажіть, що не так з наведеним кодом.
\vspace{-3mm}
\begin{verbatim}int *pa=new int(2), *pb=pa; delete pb; *pa=3;\end{verbatim}
\vspace{-3mm}
%enditem
\par
\vspace{2mm}
%beginitem
4. За наявності означень
\vspace{-3mm}
\begin{verbatim}
class A{ };
class B{ };
class C:public A{B b;A a;};\end{verbatim}
\vspace{-3mm}
вкажіть послідовність викликів деструкторів за виконання коду \code{\{C c;\}} 

(Якщо деструктор викликається кілька разів поспіль, то має зазначатися кожен його виклик, наприклад, $\sim$A, $\sim$A, $\sim$A.)
%enditem

\vspace{2mm}
%beginitem
5. Яку часову оцінку складності має алгоритм сортування вставками?
%enditem
\par
\vspace{2mm}
%beginitem
6. Модифікувати клас 

\qquad\code{class Session\{ int DM, MA, Alg, Prog;\};} 


додавши до нього оператор перетворення на тип \code{bool}: об'єкт класу є істинним тоді і тільки тоді, коли значення всіх атрибутів не менші за 60. Означення оператора теж має бути записане.%enditem
\par
\vspace{2mm}
%beginitem
7. 
Задано клас A (структуру, оператори > та == десь коректно означено) та клас вузла списку Node.

struct A\{

\quad  якісь коректні означення 

\quad  bool operator>(const A\&); 

\quad  bool operator==(const A\&); 

\quad A\& operator =(const A\&)=delete; 

\quad  A\& operator =(A\&\&)=delete;\};

struct Node\{Node *next; A a;\};

 Написати функцію, яка вилучає зі списку всі вузли, елемент в яких менший елемента в попередньому вузлі списку або дорівнює елементу в наступному вузлі. Список задається вказівником на перший вузол списку. Наприклад, зі списку, що містить послідовність 8, 5, 6, 9, 4, мають бути вилучені вузли, що містять 5 та 4. У цьому прикладі 6 залишається в списку.

Крім правильності та відповідності умові оцінюється якість коду, наявність витоків ресурсів та коректність роботи з пам'яттю. 

%enditem
}
{\smallskip\smallskip \begin {scriptsize} Затверджено на засіданні кафедри теоретичної кібернетики, протокол \No 4  від   25.01.2024 р.\par\smallskip\smallskip\smallskip\smallskip
{\parbox {6.5 cm} {Зав. кафедри \dotfill Ю.В.~Крак}}\hspace {0.7cm} {\hbox to 6.5 cm { Екзаменатор \dotfill Т.О.~Карнаух}} \end{scriptsize}}
%end
\newpage
%begin
{{\begin {scriptsize}\par
{ \center  Київський національний університет імені Тараса Шевченка \par}
{\parbox {5 cm}{Освітня програма \hfill}{\parbox {7 cm}{Інформатика\hfill}}\parbox {6 cm}{\hfill Семестр 2}}\par
{\parbox {5 cm} {Навчальний предмет \hfill}\parbox {7 cm}{Програмування \hfill}\parbox {6cm}{\hfill}\par}\vspace{-15pt} \end{scriptsize}}{\center Екзаменаційний білет \No \textbf { 51 }\par}}
{%beginitem
1. Вкажіть основну \textbf{семантичну} відмінність між конструктором копії та конструктором переміщення.%enditem
\par
\vspace{2mm}
%beginitem
2. Максимально змістовно повно сформулюйте назву оголошеного в наведеному фрагменті методу класу.
\vspace{-3mm}
\begin{verbatim}class A{
   ...
   A();
};\end{verbatim}
\vspace{-3mm}
%enditem
\par
\vspace{2mm}
%beginitem
3. Вкажіть, що не так з наведеним кодом.
\vspace{-3mm}
\begin{verbatim}int *pa=new int(2), *pb=new int(3); pa=pb;\end{verbatim}
\vspace{-3mm}
%enditem
\par
\vspace{2mm}
%beginitem
4. За наявності означень
\vspace{-3mm}
\begin{verbatim}
class A{ };
class B:public A{ };
class C:public B{ };\end{verbatim}
\vspace{-3mm}
вкажіть послідовність викликів деструкторів за виконання коду \code{\{C c;\}} 

(Якщо деструктор викликається кілька разів поспіль, то має зазначатися кожен його виклик, наприклад, $\sim$A, $\sim$A, $\sim$A.)
%enditem

\vspace{2mm}
%beginitem
5. Яку часову оцінку складності має алгоритм сортування бульбашкою?
%enditem
\par
\vspace{2mm}
%beginitem
6. Модифікувати клас 

\qquad\code{class Session\{ int DM, MA, Alg, Prog;\};} 


додавши до нього оператор перетворення на тип \code{bool}: об'єкт класу є істинним тоді і тільки тоді, коли значення всіх атрибутів не менші за 60. Означення оператора теж має бути записане.%enditem
\par
\vspace{2mm}
%beginitem
7. 
Задано клас A (структуру, оператори < та == десь коректно означено) та клас вузла списку Node.

struct A\{

\quad  якісь коректні означення 

\quad  bool operator<(const A\&); 

\quad A\& operator =(const A\&)=delete; 

\quad  A\& operator =(A\&\&)=delete;\};

struct Node\{Node *prev; Node *next; A a;\};

Написати функцію, яка вилучає зі списку всі вузли, елемент в яких менший елемента в попередньому вузлі списку або дорівнює елементу в наступному вузлі. Список задається вказівником на перший вузол списку. Наприклад, зі списку, що містить послідовність 8, 5, 6, 9, 4, мають бути вилучені вузли, що містять 5 та 4. У цьому прикладі 6 залишається в списку.

Крім правильності та відповідності умові оцінюється якість коду, наявність витоків ресурсів та коректність роботи з пам'яттю. 

%enditem
}
{\smallskip\smallskip \begin {scriptsize} Затверджено на засіданні кафедри теоретичної кібернетики, протокол \No 4  від   25.01.2024 р.\par\smallskip\smallskip\smallskip\smallskip
{\parbox {6.5 cm} {Зав. кафедри \dotfill Ю.В.~Крак}}\hspace {0.7cm} {\hbox to 6.5 cm { Екзаменатор \dotfill Т.О.~Карнаух}} \end{scriptsize}}
%end
\newpage
%begin
{{\begin {scriptsize}\par
{ \center  Київський національний університет імені Тараса Шевченка \par}
{\parbox {5 cm}{Освітня програма \hfill}{\parbox {7 cm}{Інформатика\hfill}}\parbox {6 cm}{\hfill Семестр 2}}\par
{\parbox {5 cm} {Навчальний предмет \hfill}\parbox {7 cm}{Програмування \hfill}\parbox {6cm}{\hfill}\par}\vspace{-15pt} \end{scriptsize}}{\center Екзаменаційний білет \No \textbf { 52 }\par}}
{%beginitem
1. Вкажіть основну відмінність між віртуальними методами та тими, що не є віртуальними.%enditem
\par
\vspace{2mm}
%beginitem
2. Максимально змістовно повно сформулюйте назву оголошеного в наведеному фрагменті методу класу.
\vspace{-3mm}
\begin{verbatim}class A{
   ...
   A(const A&);
};\end{verbatim}
\vspace{-3mm}
%enditem
\par
\vspace{2mm}
%beginitem
3. Вкажіть, що не так з наведеним кодом.
\vspace{-3mm}
\begin{verbatim}int *pa=new int(2), *pb=pa; delete pb; delete pa;\end{verbatim}
\vspace{-3mm}
%enditem
\par
\vspace{2mm}
%beginitem
4. За наявності означень
\vspace{-3mm}
\begin{verbatim}
class A{ };
class B:public A{ };
class C:public A{B b;};\end{verbatim}
\vspace{-3mm}
вкажіть послідовність викликів конструкторів за виконання коду \code{\{C c;\}} 

(Якщо конструктор викликається кілька разів поспіль, то має зазначатися кожен його виклик, наприклад, A, A, A.)
%enditem

\vspace{2mm}
%beginitem
5. Яку часову оцінку складності має алгоритм швидкого сортування?
%enditem
\par
\vspace{2mm}
%beginitem
6. Модифікувати клас точок простору 

\qquad\code{class Point\{ double x, y, z;\};}


додавши до нього оператор перетворення на тип \code{bool}: точка є істинною тоді і тільки тоді, коли містить хоча б одну ненульову координату. Означення оператора теж має бути записане.%enditem
\par
\vspace{2mm}
%beginitem
7. 
Реалізувати операцію перетину множин з повтореннями. Множина з повтореннями задається переліком своїх елементів, записаних за неспаданням у текстовому файлі через білі символи. Результат у такому ж самому форматі (і теж за неспаданням) має бути записаний у текстовий файл. Питома функція має отримувати три шляхи до файлів. Розмір вхідних множин такий, що вони не можуть бути завантажені повністю в пам’ять програми.

Крім правильності та відповідності умові оцінюється якість коду, наявність витоків ресурсів та коректність роботи з пам'яттю. 

%enditem
}
{\smallskip\smallskip \begin {scriptsize} Затверджено на засіданні кафедри теоретичної кібернетики, протокол \No 4  від   25.01.2024 р.\par\smallskip\smallskip\smallskip\smallskip
{\parbox {6.5 cm} {Зав. кафедри \dotfill Ю.В.~Крак}}\hspace {0.7cm} {\hbox to 6.5 cm { Екзаменатор \dotfill Т.О.~Карнаух}} \end{scriptsize}}
%end
\newpage
%begin
{{\begin {scriptsize}\par
{ \center  Київський національний університет імені Тараса Шевченка \par}
{\parbox {5 cm}{Освітня програма \hfill}{\parbox {7 cm}{Інформатика\hfill}}\parbox {6 cm}{\hfill Семестр 2}}\par
{\parbox {5 cm} {Навчальний предмет \hfill}\parbox {7 cm}{Програмування \hfill}\parbox {6cm}{\hfill}\par}\vspace{-15pt} \end{scriptsize}}{\center Екзаменаційний білет \No \textbf { 53 }\par}}
{%beginitem
1. Вкажіть основну \textbf{синтаксичну} відмінність між конструктором копії та конструктором переміщення.%enditem
\par
\vspace{2mm}
%beginitem
2. Максимально змістовно повно сформулюйте назву оголошеного в наведеному фрагменті методу класу.
\vspace{-3mm}
\begin{verbatim}class A{
   ...
   A& operator =(A&&);
};\end{verbatim}
\vspace{-3mm}
%enditem
\par
\vspace{2mm}
%beginitem
3. Вкажіть, що не так з наведеним кодом.
\vspace{-3mm}
\begin{verbatim}int *pa=new int(2), *pb=*pa; delete pa;\end{verbatim}
\vspace{-3mm}
%enditem
\par
\vspace{2mm}
%beginitem
4. За наявності означень
\vspace{-3mm}
\begin{verbatim}
class A{ };
class B{ };
class C:public A{B b[2];};\end{verbatim}
\vspace{-3mm}
вкажіть послідовність викликів конструкторів за виконання коду \code{\{C c;\}} 

(Якщо конструктор викликається кілька разів поспіль, то має зазначатися кожен його виклик, наприклад, A, A, A.)
%enditem

\vspace{2mm}
%beginitem
5. Яку часову оцінку складності має алгоритм сортування купою (він же деревом, пірамідою)?
%enditem
\par
\vspace{2mm}
%beginitem
6. Модифікувати клас точок простору 

\qquad\code{class Point\{ double x, y, z;\};}


додавши до нього оператор перетворення на рядок типу \code{std::string}, що містить рядкове зображення точки простору (у круглих дужках координати через кому), наприклад, (1.01, 2.35, 0.5), та його означення.%enditem
\par
\vspace{2mm}
%beginitem
7. 
Реалізувати операцію симетричної різниці множин з повтореннями. Множина з повтореннями задається переліком своїх елементів, записаних за неспаданням у текстовому файлі через білі символи. Результат у такому ж самому форматі (і теж за неспаданням) має бути записаний у текстовий файл. Питома функція має отримувати три шляхи до файлів. Розмір вхідних множин такий, що вони не можуть бути завантажені повністю в пам’ять програми.

Крім правильності та відповідності умові оцінюється якість коду, наявність витоків ресурсів та коректність роботи з пам'яттю. 

%enditem
}
{\smallskip\smallskip \begin {scriptsize} Затверджено на засіданні кафедри теоретичної кібернетики, протокол \No 4  від   25.01.2024 р.\par\smallskip\smallskip\smallskip\smallskip
{\parbox {6.5 cm} {Зав. кафедри \dotfill Ю.В.~Крак}}\hspace {0.7cm} {\hbox to 6.5 cm { Екзаменатор \dotfill Т.О.~Карнаух}} \end{scriptsize}}
%end
\newpage
%begin
{{\begin {scriptsize}\par
{ \center  Київський національний університет імені Тараса Шевченка \par}
{\parbox {5 cm}{Освітня програма \hfill}{\parbox {7 cm}{Інформатика\hfill}}\parbox {6 cm}{\hfill Семестр 2}}\par
{\parbox {5 cm} {Навчальний предмет \hfill}\parbox {7 cm}{Програмування \hfill}\parbox {6cm}{\hfill}\par}\vspace{-15pt} \end{scriptsize}}{\center Екзаменаційний білет \No \textbf { 54 }\par}}
{%beginitem
1. Що важливо забезпечити при роботі з динамічними змінними?%enditem
\par
\vspace{2mm}
%beginitem
2. Максимально змістовно повно сформулюйте назву оголошеного в наведеному фрагменті методу класу.
\vspace{-3mm}
\begin{verbatim}class A{
   ...
   A(A&&);
};\end{verbatim}
\vspace{-3mm}
%enditem
\par
\vspace{2mm}
%beginitem
3. Вкажіть, що не так з наведеним кодом.
\vspace{-3mm}
\begin{verbatim}int **ppi=new int*(new int(3)); delete ppi;\end{verbatim}
\vspace{-3mm}
%enditem
\par
\vspace{2mm}
%beginitem
4. За наявності означень
\vspace{-3mm}
\begin{verbatim}
class A{ };
class B:public A{ };
class C:public B{ };\end{verbatim}
\vspace{-3mm}
вкажіть послідовність викликів конструкторів за виконання коду \code{\{C c;\}} 

(Якщо конструктор викликається кілька разів поспіль, то має зазначатися кожен його виклик, наприклад, A, A, A.)
%enditem

\vspace{2mm}
%beginitem
5. Яку часову оцінку складності має алгоритм сортування злиттям?
%enditem
\par
\vspace{2mm}
%beginitem
6. Модифікувати клас 

\qquad\code{class Session\{ int DM, MA, Alg, Prog;\};} 


додавши до нього оператор порівняння \code{<}: один об'єкт є меншим за інший тоді і тільки тоді, коли значення всіх атрибутів об'єктів не менші за 60 і сума атрибутів першого менша за суму атрибутів другого. Означення оператора теж має бути записане.%enditem
\par
\vspace{2mm}
%beginitem
7. 
Задано клас A (структуру, оператори < та == десь коректно означено) та клас вузла списку Node.

struct A\{

\quad  якісь коректні означення 

\quad  bool operator<(const A\&); 

\quad  bool operator==(const A\&); 

\quad A\& operator =(const A\&)=delete; 

\quad  A\& operator =(A\&\&)=delete;\};

struct Node\{Node *next; A a;\};

 Написати функцію, яка вилучає зі списку всі вузли, елемент в яких більший елемента в попередньому вузлі списку або дорівнює елементу в наступному вузлі. Список є циклічним та задається вказівником на перший вузол списку. Наприклад, зі списку, що містить послідовність 8, 10, 9, 4, мають бути вилучені вузли, що містять 8 (бо список циклічний і $8>4$) та 10. У цьому прикладі 9 залишається в списку.

Крім правильності та відповідності умові оцінюється якість коду, наявність витоків ресурсів та коректність роботи з пам'яттю. 

%enditem
}
{\smallskip\smallskip \begin {scriptsize} Затверджено на засіданні кафедри теоретичної кібернетики, протокол \No 4  від   25.01.2024 р.\par\smallskip\smallskip\smallskip\smallskip
{\parbox {6.5 cm} {Зав. кафедри \dotfill Ю.В.~Крак}}\hspace {0.7cm} {\hbox to 6.5 cm { Екзаменатор \dotfill Т.О.~Карнаух}} \end{scriptsize}}
%end
\newpage
%begin
{{\begin {scriptsize}\par
{ \center  Київський національний університет імені Тараса Шевченка \par}
{\parbox {5 cm}{Освітня програма \hfill}{\parbox {7 cm}{Інформатика\hfill}}\parbox {6 cm}{\hfill Семестр 2}}\par
{\parbox {5 cm} {Навчальний предмет \hfill}\parbox {7 cm}{Програмування \hfill}\parbox {6cm}{\hfill}\par}\vspace{-15pt} \end{scriptsize}}{\center Екзаменаційний білет \No \textbf { 55 }\par}}
{%beginitem
1. Вкажіть основну відмінність між віртуальними методами та тими, що не є віртуальними.%enditem
\par
\vspace{2mm}
%beginitem
2. Максимально змістовно повно сформулюйте назву оголошеного в наведеному фрагменті методу класу.
\vspace{-3mm}
\begin{verbatim}class A{
   ...
   ~A();
};\end{verbatim}
\vspace{-3mm}
%enditem
\par
\vspace{2mm}
%beginitem
3. Вкажіть, що не так з наведеним кодом.
\vspace{-3mm}
\begin{verbatim}int *pa=new int(2), *pb=&pa; delete pa;\end{verbatim}
\vspace{-3mm}
%enditem
\par
\vspace{2mm}
%beginitem
4. За наявності означень
\vspace{-3mm}
\begin{verbatim}
class A{ };
class B{ };
class C{B b;A a;};\end{verbatim}
\vspace{-3mm}
вкажіть послідовність викликів деструкторів за виконання коду \code{\{C c;\}} 

(Якщо деструктор викликається кілька разів поспіль, то має зазначатися кожен його виклик, наприклад, $\sim$A, $\sim$A, $\sim$A.)
%enditem

\vspace{2mm}
%beginitem
5. Яку часову оцінку складності має алгоритм швидкого сортування?
%enditem
\par
\vspace{2mm}
%beginitem
6. Модифікувати клас 

\qquad\code{class Session\{ int DM, MA, Alg, Prog;\};} 


додавши до нього оператор перетворення на тип \code{int}: якщо хоча б один атрибут менший за 60, то має повертатись 0, інакше -- сума атрибутів.%enditem
\par
\vspace{2mm}
%beginitem
7. 
Реалізувати операцію злиття множин з повтореннями. Множина з повтореннями задається переліком своїх елементів, записаних за неспаданням у текстовому файлі через білі символи. Результат у такому ж самому форматі (і теж за неспаданням) має бути записаний у текстовий файл. Питома функція має отримувати три шляхи до файлів. Розмір вхідних множин такий, що вони не можуть бути завантажені повністю в пам’ять програми.

Крім правильності та відповідності умові оцінюється якість коду, наявність витоків ресурсів та коректність роботи з пам'яттю. 

%enditem
}
{\smallskip\smallskip \begin {scriptsize} Затверджено на засіданні кафедри теоретичної кібернетики, протокол \No 4  від   25.01.2024 р.\par\smallskip\smallskip\smallskip\smallskip
{\parbox {6.5 cm} {Зав. кафедри \dotfill Ю.В.~Крак}}\hspace {0.7cm} {\hbox to 6.5 cm { Екзаменатор \dotfill Т.О.~Карнаух}} \end{scriptsize}}
%end
\newpage
%begin
{{\begin {scriptsize}\par
{ \center  Київський національний університет імені Тараса Шевченка \par}
{\parbox {5 cm}{Освітня програма \hfill}{\parbox {7 cm}{Інформатика\hfill}}\parbox {6 cm}{\hfill Семестр 2}}\par
{\parbox {5 cm} {Навчальний предмет \hfill}\parbox {7 cm}{Програмування \hfill}\parbox {6cm}{\hfill}\par}\vspace{-15pt} \end{scriptsize}}{\center Екзаменаційний білет \No \textbf { 56 }\par}}
{%beginitem
1. Що важливо забезпечити при роботі з динамічними змінними?%enditem
\par
\vspace{2mm}
%beginitem
2. Максимально змістовно повно сформулюйте назву оголошеного в наведеному фрагменті методу класу.
\vspace{-3mm}
\begin{verbatim}class A{
   ...
   A();
};\end{verbatim}
\vspace{-3mm}
%enditem
\par
\vspace{2mm}
%beginitem
3. Вкажіть, що не так з наведеним кодом.
\vspace{-3mm}
\begin{verbatim}int a, *p=&a; delete a;\end{verbatim}
\vspace{-3mm}
%enditem
\par
\vspace{2mm}
%beginitem
4. За наявності означень
\vspace{-3mm}
\begin{verbatim}
class A{ };
class B{ };
class C:public A{B b;A a;};\end{verbatim}
\vspace{-3mm}
вкажіть послідовність викликів конструкторів за виконання коду \code{\{C c;\}} 

(Якщо конструктор викликається кілька разів поспіль, то має зазначатися кожен його виклик, наприклад, A, A, A.)
%enditem

\vspace{2mm}
%beginitem
5. Яку часову оцінку складності має алгоритм сортування вибором?
%enditem
\par
\vspace{2mm}
%beginitem
6. Модифікувати клас точок простору 

\qquad\code{class Point\{ double x, y, z;\};}


додавши до нього оператор перетворення на рядок типу \code{std::string}, що містить рядкове зображення точки простору (у круглих дужках координати через кому), наприклад, (1.01, 2.35, 0.5), та його означення.%enditem
\par
\vspace{2mm}
%beginitem
7. 
Задано клас A (структуру, оператори < та == десь коректно означено) та клас вузла списку Node.

struct A\{

\quad  якісь коректні означення 

\quad  bool operator<(const A\&); 

\quad  bool operator==(const A\&); 

\quad A\& operator =(const A\&)=delete; 

\quad  A\& operator =(A\&\&)=delete;\};

struct Node\{Node *prev; Node *next; A a;\};

Написати функцію, яка вилучає зі списку всі вузли, елемент в яких більший елемента в попередньому вузлі списку або дорівнює елементу в наступному вузлі. Список є циклічним та задається вказівником на перший вузол списку. Наприклад, зі списку, що містить послідовність 8, 10, 9, 4, мають бути вилучені вузли, що містять 8 (бо список циклічний і $8>4$) та 10. У цьому прикладі 9 залишається в списку.

Крім правильності та відповідності умові оцінюється якість коду, наявність витоків ресурсів та коректність роботи з пам'яттю. 

%enditem
}
{\smallskip\smallskip \begin {scriptsize} Затверджено на засіданні кафедри теоретичної кібернетики, протокол \No 4  від   25.01.2024 р.\par\smallskip\smallskip\smallskip\smallskip
{\parbox {6.5 cm} {Зав. кафедри \dotfill Ю.В.~Крак}}\hspace {0.7cm} {\hbox to 6.5 cm { Екзаменатор \dotfill Т.О.~Карнаух}} \end{scriptsize}}
%end
\newpage
%begin
{{\begin {scriptsize}\par
{ \center  Київський національний університет імені Тараса Шевченка \par}
{\parbox {5 cm}{Освітня програма \hfill}{\parbox {7 cm}{Інформатика\hfill}}\parbox {6 cm}{\hfill Семестр 2}}\par
{\parbox {5 cm} {Навчальний предмет \hfill}\parbox {7 cm}{Програмування \hfill}\parbox {6cm}{\hfill}\par}\vspace{-15pt} \end{scriptsize}}{\center Екзаменаційний білет \No \textbf { 57 }\par}}
{%beginitem
1. Вкажіть основну \textbf{семантичну} відмінність між конструктором копії та конструктором переміщення.%enditem
\par
\vspace{2mm}
%beginitem
2. Максимально змістовно повно сформулюйте назву оголошеного в наведеному фрагменті методу класу.
\vspace{-3mm}
\begin{verbatim}class A{
   ...
   ~A();
};\end{verbatim}
\vspace{-3mm}
%enditem
\par
\vspace{2mm}
%beginitem
3. Вкажіть, що не так з наведеним кодом.
\vspace{-3mm}
\begin{verbatim}int *pa=new int(2), *pb=*pa; delete pa;\end{verbatim}
\vspace{-3mm}
%enditem
\par
\vspace{2mm}
%beginitem
4. За наявності означень
\vspace{-3mm}
\begin{verbatim}
class A{ };
class B:public A{ };
class C:public A{B b;};\end{verbatim}
\vspace{-3mm}
вкажіть послідовність викликів деструкторів за виконання коду \code{\{C c;\}} 

(Якщо деструктор викликається кілька разів поспіль, то має зазначатися кожен його виклик, наприклад, $\sim$A, $\sim$A, $\sim$A.)
%enditem

\vspace{2mm}
%beginitem
5. Яку часову оцінку складності має алгоритм сортування вставками?
%enditem
\par
\vspace{2mm}
%beginitem
6. Модифікувати клас 

\qquad\code{class Session\{ int DM, MA, Alg, Prog;\};} 


додавши до нього оператор перетворення на тип \code{bool}: об'єкт класу є істинним тоді і тільки тоді, коли значення всіх атрибутів не менші за 60. Означення оператора теж має бути записане.%enditem
\par
\vspace{2mm}
%beginitem
7. 
Задано клас A (структуру, оператори < та == десь коректно означено) та клас вузла списку Node.

struct A\{

\quad  якісь коректні означення 

\quad  bool operator<(const A\&); 

\quad  bool operator==(const A\&); 

\quad A\& operator =(const A\&)=delete; 

\quad  A\& operator =(A\&\&)=delete;\};

struct Node\{Node *next; A a;\};

 Написати функцію, яка вилучає зі списку всі вузли, елемент в яких більший елемента в попередньому вузлі списку або дорівнює елементу в наступному вузлі. Список є циклічним та задається вказівником на перший вузол списку. Наприклад, зі списку, що містить послідовність 8, 10, 9, 4, мають бути вилучені вузли, що містять 8 (бо список циклічний і $8>4$) та 10. У цьому прикладі 9 залишається в списку.

Крім правильності та відповідності умові оцінюється якість коду, наявність витоків ресурсів та коректність роботи з пам'яттю. 

%enditem
}
{\smallskip\smallskip \begin {scriptsize} Затверджено на засіданні кафедри теоретичної кібернетики, протокол \No 4  від   25.01.2024 р.\par\smallskip\smallskip\smallskip\smallskip
{\parbox {6.5 cm} {Зав. кафедри \dotfill Ю.В.~Крак}}\hspace {0.7cm} {\hbox to 6.5 cm { Екзаменатор \dotfill Т.О.~Карнаух}} \end{scriptsize}}
%end
\newpage
%begin
{{\begin {scriptsize}\par
{ \center  Київський національний університет імені Тараса Шевченка \par}
{\parbox {5 cm}{Освітня програма \hfill}{\parbox {7 cm}{Інформатика\hfill}}\parbox {6 cm}{\hfill Семестр 2}}\par
{\parbox {5 cm} {Навчальний предмет \hfill}\parbox {7 cm}{Програмування \hfill}\parbox {6cm}{\hfill}\par}\vspace{-15pt} \end{scriptsize}}{\center Екзаменаційний білет \No \textbf { 58 }\par}}
{%beginitem
1. Вкажіть основну відмінність між класами (class) та структурами (struct).%enditem
\par
\vspace{2mm}
%beginitem
2. Максимально змістовно повно сформулюйте назву оголошеного в наведеному фрагменті методу класу.
\vspace{-3mm}
\begin{verbatim}class A{
   ...
   A(A&&);
};\end{verbatim}
\vspace{-3mm}
%enditem
\par
\vspace{2mm}
%beginitem
3. Вкажіть, що не так з наведеним кодом.
\vspace{-3mm}
\begin{verbatim}int **ppi=new int*(new int(3)); delete ppi;\end{verbatim}
\vspace{-3mm}
%enditem
\par
\vspace{2mm}
%beginitem
4. За наявності означень
\vspace{-3mm}
\begin{verbatim}
class A{ };
class B{ };
class C{B b;A a;};\end{verbatim}
\vspace{-3mm}
вкажіть послідовність викликів конструкторів за виконання коду \code{\{C c;\}} 

(Якщо конструктор викликається кілька разів поспіль, то має зазначатися кожен його виклик, наприклад, A, A, A.)
%enditem

\vspace{2mm}
%beginitem
5. Яку часову оцінку складності має алгоритм сортування купою (він же деревом, пірамідою)?
%enditem
\par
\vspace{2mm}
%beginitem
6. Модифікувати клас точок простору 

\qquad\code{class Point\{ double x, y, z;\};}


додавши до нього оператор перетворення на тип \code{bool}: точка є істинною тоді і тільки тоді, коли містить хоча б одну ненульову координату. Означення оператора теж має бути записане.%enditem
\par
\vspace{2mm}
%beginitem
7. 
Задано клас A (структуру, оператори > та == десь коректно означено) та клас вузла списку Node.

struct A\{

\quad  якісь коректні означення 

\quad  bool operator>(const A\&); 

\quad  bool operator==(const A\&); 

\quad A\& operator =(const A\&)=delete; 

\quad  A\& operator =(A\&\&)=delete;\};

struct Node\{Node *next; A a;\};

 Написати функцію, яка вилучає зі списку всі вузли, елемент в яких менший елемента в попередньому вузлі списку або дорівнює елементу в наступному вузлі. Список задається вказівником на перший вузол списку. Наприклад, зі списку, що містить послідовність 8, 5, 6, 9, 4, мають бути вилучені вузли, що містять 5 та 4. У цьому прикладі 6 залишається в списку.

Крім правильності та відповідності умові оцінюється якість коду, наявність витоків ресурсів та коректність роботи з пам'яттю. 

%enditem
}
{\smallskip\smallskip \begin {scriptsize} Затверджено на засіданні кафедри теоретичної кібернетики, протокол \No 4  від   25.01.2024 р.\par\smallskip\smallskip\smallskip\smallskip
{\parbox {6.5 cm} {Зав. кафедри \dotfill Ю.В.~Крак}}\hspace {0.7cm} {\hbox to 6.5 cm { Екзаменатор \dotfill Т.О.~Карнаух}} \end{scriptsize}}
%end
\newpage
%begin
{{\begin {scriptsize}\par
{ \center  Київський національний університет імені Тараса Шевченка \par}
{\parbox {5 cm}{Освітня програма \hfill}{\parbox {7 cm}{Інформатика\hfill}}\parbox {6 cm}{\hfill Семестр 2}}\par
{\parbox {5 cm} {Навчальний предмет \hfill}\parbox {7 cm}{Програмування \hfill}\parbox {6cm}{\hfill}\par}\vspace{-15pt} \end{scriptsize}}{\center Екзаменаційний білет \No \textbf { 59 }\par}}
{%beginitem
1. Вкажіть основну відмінність між відкритими (public) та прихованими (private) методами класу.%enditem
\par
\vspace{2mm}
%beginitem
2. Максимально змістовно повно сформулюйте назву оголошеного в наведеному фрагменті методу класу.
\vspace{-3mm}
\begin{verbatim}class A{
   ...
   A& operator =(A&&);
};\end{verbatim}
\vspace{-3mm}
%enditem
\par
\vspace{2mm}
%beginitem
3. Вкажіть, що не так з наведеним кодом.
\vspace{-3mm}
\begin{verbatim}int *pa=new int(2), *pb=pa; delete pb; delete pa;\end{verbatim}
\vspace{-3mm}
%enditem
\par
\vspace{2mm}
%beginitem
4. За наявності означень
\vspace{-3mm}
\begin{verbatim}
class A{ };
class B{ };
class C:public A{B b;};\end{verbatim}
\vspace{-3mm}
вкажіть послідовність викликів деструкторів за виконання коду \code{\{C c;\}} 

(Якщо деструктор викликається кілька разів поспіль, то має зазначатися кожен його виклик, наприклад, $\sim$A, $\sim$A, $\sim$A.)
%enditem

\vspace{2mm}
%beginitem
5. Яку часову оцінку складності має алгоритм сортування бульбашкою?
%enditem
\par
\vspace{2mm}
%beginitem
6. Модифікувати клас 

\qquad\code{class Session\{ int DM, MA, Alg, Prog;\};} 


додавши до нього оператор порівняння \code{<}: один об'єкт є меншим за інший тоді і тільки тоді, коли значення всіх атрибутів об'єктів не менші за 60 і сума атрибутів першого менша за суму атрибутів другого. Означення оператора теж має бути записане.%enditem
\par
\vspace{2mm}
%beginitem
7. 
Реалізувати операцію різниці множин з повтореннями. Множина з повтореннями задається переліком своїх елементів, записаних за неспаданням у текстовому файлі через білі символи. Результат у такому ж самому форматі (і теж за неспаданням) має бути записаний у текстовий файл. Питома функція має отримувати три шляхи до файлів. Розмір вхідних множин такий, що вони не можуть бути завантажені повністю в пам’ять програми.

Крім правильності та відповідності умові оцінюється якість коду, наявність витоків ресурсів та коректність роботи з пам'яттю. 

%enditem
}
{\smallskip\smallskip \begin {scriptsize} Затверджено на засіданні кафедри теоретичної кібернетики, протокол \No 4  від   25.01.2024 р.\par\smallskip\smallskip\smallskip\smallskip
{\parbox {6.5 cm} {Зав. кафедри \dotfill Ю.В.~Крак}}\hspace {0.7cm} {\hbox to 6.5 cm { Екзаменатор \dotfill Т.О.~Карнаух}} \end{scriptsize}}
%end
\newpage
%begin
{{\begin {scriptsize}\par
{ \center  Київський національний університет імені Тараса Шевченка \par}
{\parbox {5 cm}{Освітня програма \hfill}{\parbox {7 cm}{Інформатика\hfill}}\parbox {6 cm}{\hfill Семестр 2}}\par
{\parbox {5 cm} {Навчальний предмет \hfill}\parbox {7 cm}{Програмування \hfill}\parbox {6cm}{\hfill}\par}\vspace{-15pt} \end{scriptsize}}{\center Екзаменаційний білет \No \textbf { 60 }\par}}
{%beginitem
1. Вкажіть основну \textbf{синтаксичну} відмінність між конструктором копії та конструктором переміщення.%enditem
\par
\vspace{2mm}
%beginitem
2. Максимально змістовно повно сформулюйте назву оголошеного в наведеному фрагменті методу класу.
\vspace{-3mm}
\begin{verbatim}class A{
   ...
   A(const A&);
};\end{verbatim}
\vspace{-3mm}
%enditem
\par
\vspace{2mm}
%beginitem
3. Вкажіть, що не так з наведеним кодом.
\vspace{-3mm}
\begin{verbatim}int *pa=new int(2), *pb=new int(3); pa=pb;\end{verbatim}
\vspace{-3mm}
%enditem
\par
\vspace{2mm}
%beginitem
4. За наявності означень
\vspace{-3mm}
\begin{verbatim}
class A{ };
class B{ };
class C:public A{B b[2];};\end{verbatim}
\vspace{-3mm}
вкажіть послідовність викликів деструкторів за виконання коду \code{\{C c;\}} 

(Якщо деструктор викликається кілька разів поспіль, то має зазначатися кожен його виклик, наприклад, $\sim$A, $\sim$A, $\sim$A.)
%enditem

\vspace{2mm}
%beginitem
5. Яку часову оцінку складності має алгоритм швидкого сортування?
%enditem
\par
\vspace{2mm}
%beginitem
6. Модифікувати клас 

\qquad\code{class Session\{ int DM, MA, Alg, Prog;\};} 


додавши до нього оператор перетворення на тип \code{int}: якщо хоча б один атрибут менший за 60, то має повертатись 0, інакше -- сума атрибутів.%enditem
\par
\vspace{2mm}
%beginitem
7. 
Задано клас A (структуру, оператори < та == десь коректно означено) та клас вузла списку Node.

struct A\{

\quad  якісь коректні означення 

\quad  bool operator<(const A\&); 

\quad A\& operator =(const A\&)=delete; 

\quad  A\& operator =(A\&\&)=delete;\};

struct Node\{Node *prev; Node *next; A a;\};

Написати функцію, яка вилучає зі списку всі вузли, елемент в яких менший елемента в попередньому вузлі списку або дорівнює елементу в наступному вузлі. Список задається вказівником на перший вузол списку. Наприклад, зі списку, що містить послідовність 8, 5, 6, 9, 4, мають бути вилучені вузли, що містять 5 та 4. У цьому прикладі 6 залишається в списку.

Крім правильності та відповідності умові оцінюється якість коду, наявність витоків ресурсів та коректність роботи з пам'яттю. 

%enditem
}
{\smallskip\smallskip \begin {scriptsize} Затверджено на засіданні кафедри теоретичної кібернетики, протокол \No 4  від   25.01.2024 р.\par\smallskip\smallskip\smallskip\smallskip
{\parbox {6.5 cm} {Зав. кафедри \dotfill Ю.В.~Крак}}\hspace {0.7cm} {\hbox to 6.5 cm { Екзаменатор \dotfill Т.О.~Карнаух}} \end{scriptsize}}
%end
\newpage
%begin
{{\begin {scriptsize}\par
{ \center  Київський національний університет імені Тараса Шевченка \par}
{\parbox {5 cm}{Освітня програма \hfill}{\parbox {7 cm}{Інформатика\hfill}}\parbox {6 cm}{\hfill Семестр 2}}\par
{\parbox {5 cm} {Навчальний предмет \hfill}\parbox {7 cm}{Програмування \hfill}\parbox {6cm}{\hfill}\par}\vspace{-15pt} \end{scriptsize}}{\center Екзаменаційний білет \No \textbf { 61 }\par}}
{%beginitem
1. Вкажіть основну відмінність між відкритими (public) та прихованими (private) методами класу.%enditem
\par
\vspace{2mm}
%beginitem
2. Максимально змістовно повно сформулюйте назву оголошеного в наведеному фрагменті методу класу.
\vspace{-3mm}
\begin{verbatim}class A{
   ...
   A(const A&);
};\end{verbatim}
\vspace{-3mm}
%enditem
\par
\vspace{2mm}
%beginitem
3. Вкажіть, що не так з наведеним кодом.
\vspace{-3mm}
\begin{verbatim}int *pa=new int(2), *pb=pa; delete pb; *pa=3;\end{verbatim}
\vspace{-3mm}
%enditem
\par
\vspace{2mm}
%beginitem
4. За наявності означень
\vspace{-3mm}
\begin{verbatim}
class A{ };
class B{ };
class C:public A{B b;};\end{verbatim}
\vspace{-3mm}
вкажіть послідовність викликів конструкторів за виконання коду \code{\{C c;\}} 

(Якщо конструктор викликається кілька разів поспіль, то має зазначатися кожен його виклик, наприклад, A, A, A.)
%enditem

\vspace{2mm}
%beginitem
5. Яку часову оцінку складності має алгоритм сортування злиттям?
%enditem
\par
\vspace{2mm}
%beginitem
6. Модифікувати клас 

\qquad\code{class Session\{ int DM, MA, Alg, Prog;\};} 


додавши до нього оператор порівняння \code{<}: один об'єкт є меншим за інший тоді і тільки тоді, коли значення всіх атрибутів об'єктів не менші за 60 і сума атрибутів першого менша за суму атрибутів другого. Означення оператора теж має бути записане.%enditem
\par
\vspace{2mm}
%beginitem
7. 
Реалізувати операцію симетричної різниці множин з повтореннями. Множина з повтореннями задається переліком своїх елементів, записаних за неспаданням у текстовому файлі через білі символи. Результат у такому ж самому форматі (і теж за неспаданням) має бути записаний у текстовий файл. Питома функція має отримувати три шляхи до файлів. Розмір вхідних множин такий, що вони не можуть бути завантажені повністю в пам’ять програми.

Крім правильності та відповідності умові оцінюється якість коду, наявність витоків ресурсів та коректність роботи з пам'яттю. 

%enditem
}
{\smallskip\smallskip \begin {scriptsize} Затверджено на засіданні кафедри теоретичної кібернетики, протокол \No 4  від   25.01.2024 р.\par\smallskip\smallskip\smallskip\smallskip
{\parbox {6.5 cm} {Зав. кафедри \dotfill Ю.В.~Крак}}\hspace {0.7cm} {\hbox to 6.5 cm { Екзаменатор \dotfill Т.О.~Карнаух}} \end{scriptsize}}
%end
\newpage
%begin
{{\begin {scriptsize}\par
{ \center  Київський національний університет імені Тараса Шевченка \par}
{\parbox {5 cm}{Освітня програма \hfill}{\parbox {7 cm}{Інформатика\hfill}}\parbox {6 cm}{\hfill Семестр 2}}\par
{\parbox {5 cm} {Навчальний предмет \hfill}\parbox {7 cm}{Програмування \hfill}\parbox {6cm}{\hfill}\par}\vspace{-15pt} \end{scriptsize}}{\center Екзаменаційний білет \No \textbf { 62 }\par}}
{%beginitem
1. Що важливо забезпечити при роботі з динамічними змінними?%enditem
\par
\vspace{2mm}
%beginitem
2. Максимально змістовно повно сформулюйте назву оголошеного в наведеному фрагменті методу класу.
\vspace{-3mm}
\begin{verbatim}class A{
   ...
   A& operator =(A&&);
};\end{verbatim}
\vspace{-3mm}
%enditem
\par
\vspace{2mm}
%beginitem
3. Вкажіть, що не так з наведеним кодом.
\vspace{-3mm}
\begin{verbatim}int a, *p=&a; delete a;\end{verbatim}
\vspace{-3mm}
%enditem
\par
\vspace{2mm}
%beginitem
4. За наявності означень
\vspace{-3mm}
\begin{verbatim}
class A{ };
class B:public A{ };
class C:public A{B b;};\end{verbatim}
\vspace{-3mm}
вкажіть послідовність викликів конструкторів за виконання коду \code{\{C c;\}} 

(Якщо конструктор викликається кілька разів поспіль, то має зазначатися кожен його виклик, наприклад, A, A, A.)
%enditem

\vspace{2mm}
%beginitem
5. Яку часову оцінку складності має алгоритм сортування вибором?
%enditem
\par
\vspace{2mm}
%beginitem
6. Модифікувати клас точок простору 

\qquad\code{class Point\{ double x, y, z;\};}


додавши до нього оператор перетворення на рядок типу \code{std::string}, що містить рядкове зображення точки простору (у круглих дужках координати через кому), наприклад, (1.01, 2.35, 0.5), та його означення.%enditem
\par
\vspace{2mm}
%beginitem
7. 
Задано клас A (структуру, оператори < та == десь коректно означено) та клас вузла списку Node.

struct A\{

\quad  якісь коректні означення 

\quad  bool operator<(const A\&); 

\quad  bool operator==(const A\&); 

\quad A\& operator =(const A\&)=delete; 

\quad  A\& operator =(A\&\&)=delete;\};

struct Node\{Node *prev; Node *next; A a;\};

Написати функцію, яка вилучає зі списку всі вузли, елемент в яких більший елемента в попередньому вузлі списку або дорівнює елементу в наступному вузлі. Список є циклічним та задається вказівником на перший вузол списку. Наприклад, зі списку, що містить послідовність 8, 10, 9, 4, мають бути вилучені вузли, що містять 8 (бо список циклічний і $8>4$) та 10. У цьому прикладі 9 залишається в списку.

Крім правильності та відповідності умові оцінюється якість коду, наявність витоків ресурсів та коректність роботи з пам'яттю. 

%enditem
}
{\smallskip\smallskip \begin {scriptsize} Затверджено на засіданні кафедри теоретичної кібернетики, протокол \No 4  від   25.01.2024 р.\par\smallskip\smallskip\smallskip\smallskip
{\parbox {6.5 cm} {Зав. кафедри \dotfill Ю.В.~Крак}}\hspace {0.7cm} {\hbox to 6.5 cm { Екзаменатор \dotfill Т.О.~Карнаух}} \end{scriptsize}}
%end
\newpage
%begin
{{\begin {scriptsize}\par
{ \center  Київський національний університет імені Тараса Шевченка \par}
{\parbox {5 cm}{Освітня програма \hfill}{\parbox {7 cm}{Інформатика\hfill}}\parbox {6 cm}{\hfill Семестр 2}}\par
{\parbox {5 cm} {Навчальний предмет \hfill}\parbox {7 cm}{Програмування \hfill}\parbox {6cm}{\hfill}\par}\vspace{-15pt} \end{scriptsize}}{\center Екзаменаційний білет \No \textbf { 63 }\par}}
{%beginitem
1. Вкажіть основну відмінність між класами (class) та структурами (struct).%enditem
\par
\vspace{2mm}
%beginitem
2. Максимально змістовно повно сформулюйте назву оголошеного в наведеному фрагменті методу класу.
\vspace{-3mm}
\begin{verbatim}class A{
   ...
   A();
};\end{verbatim}
\vspace{-3mm}
%enditem
\par
\vspace{2mm}
%beginitem
3. Вкажіть, що не так з наведеним кодом.
\vspace{-3mm}
\begin{verbatim}int *pa=new int(2), *pb=&pa; delete pa;\end{verbatim}
\vspace{-3mm}
%enditem
\par
\vspace{2mm}
%beginitem
4. За наявності означень
\vspace{-3mm}
\begin{verbatim}
class A{ };
class B{ };
class C:public A{B b;};\end{verbatim}
\vspace{-3mm}
вкажіть послідовність викликів деструкторів за виконання коду \code{\{C c;\}} 

(Якщо деструктор викликається кілька разів поспіль, то має зазначатися кожен його виклик, наприклад, $\sim$A, $\sim$A, $\sim$A.)
%enditem

\vspace{2mm}
%beginitem
5. Яку часову оцінку складності має алгоритм швидкого сортування?
%enditem
\par
\vspace{2mm}
%beginitem
6. Модифікувати клас 

\qquad\code{class Session\{ int DM, MA, Alg, Prog;\};} 


додавши до нього оператор перетворення на тип \code{int}: якщо хоча б один атрибут менший за 60, то має повертатись 0, інакше -- сума атрибутів.%enditem
\par
\vspace{2mm}
%beginitem
7. 
Реалізувати операцію злиття множин з повтореннями. Множина з повтореннями задається переліком своїх елементів, записаних за неспаданням у текстовому файлі через білі символи. Результат у такому ж самому форматі (і теж за неспаданням) має бути записаний у текстовий файл. Питома функція має отримувати три шляхи до файлів. Розмір вхідних множин такий, що вони не можуть бути завантажені повністю в пам’ять програми.

Крім правильності та відповідності умові оцінюється якість коду, наявність витоків ресурсів та коректність роботи з пам'яттю. 

%enditem
}
{\smallskip\smallskip \begin {scriptsize} Затверджено на засіданні кафедри теоретичної кібернетики, протокол \No 4  від   25.01.2024 р.\par\smallskip\smallskip\smallskip\smallskip
{\parbox {6.5 cm} {Зав. кафедри \dotfill Ю.В.~Крак}}\hspace {0.7cm} {\hbox to 6.5 cm { Екзаменатор \dotfill Т.О.~Карнаух}} \end{scriptsize}}
%end
\newpage
%begin
{{\begin {scriptsize}\par
{ \center  Київський національний університет імені Тараса Шевченка \par}
{\parbox {5 cm}{Освітня програма \hfill}{\parbox {7 cm}{Інформатика\hfill}}\parbox {6 cm}{\hfill Семестр 2}}\par
{\parbox {5 cm} {Навчальний предмет \hfill}\parbox {7 cm}{Програмування \hfill}\parbox {6cm}{\hfill}\par}\vspace{-15pt} \end{scriptsize}}{\center Екзаменаційний білет \No \textbf { 64 }\par}}
{%beginitem
1. Вкажіть основну відмінність між віртуальними методами та тими, що не є віртуальними.%enditem
\par
\vspace{2mm}
%beginitem
2. Максимально змістовно повно сформулюйте назву оголошеного в наведеному фрагменті методу класу.
\vspace{-3mm}
\begin{verbatim}class A{
   ...
   A(A&&);
};\end{verbatim}
\vspace{-3mm}
%enditem
\par
\vspace{2mm}
%beginitem
3. Вкажіть, що не так з наведеним кодом.
\vspace{-3mm}
\begin{verbatim}int a, *p=&a; delete a;\end{verbatim}
\vspace{-3mm}
%enditem
\par
\vspace{2mm}
%beginitem
4. За наявності означень
\vspace{-3mm}
\begin{verbatim}
class A{ };
class B:public A{ };
class C:public B{ };\end{verbatim}
\vspace{-3mm}
вкажіть послідовність викликів конструкторів за виконання коду \code{\{C c;\}} 

(Якщо конструктор викликається кілька разів поспіль, то має зазначатися кожен його виклик, наприклад, A, A, A.)
%enditem

\vspace{2mm}
%beginitem
5. Яку часову оцінку складності має алгоритм сортування злиттям?
%enditem
\par
\vspace{2mm}
%beginitem
6. Модифікувати клас точок простору 

\qquad\code{class Point\{ double x, y, z;\};}


додавши до нього оператор перетворення на тип \code{bool}: точка є істинною тоді і тільки тоді, коли містить хоча б одну ненульову координату. Означення оператора теж має бути записане.%enditem
\par
\vspace{2mm}
%beginitem
7. 
Реалізувати операцію перетину множин з повтореннями. Множина з повтореннями задається переліком своїх елементів, записаних за неспаданням у текстовому файлі через білі символи. Результат у такому ж самому форматі (і теж за неспаданням) має бути записаний у текстовий файл. Питома функція має отримувати три шляхи до файлів. Розмір вхідних множин такий, що вони не можуть бути завантажені повністю в пам’ять програми.

Крім правильності та відповідності умові оцінюється якість коду, наявність витоків ресурсів та коректність роботи з пам'яттю. 

%enditem
}
{\smallskip\smallskip \begin {scriptsize} Затверджено на засіданні кафедри теоретичної кібернетики, протокол \No 4  від   25.01.2024 р.\par\smallskip\smallskip\smallskip\smallskip
{\parbox {6.5 cm} {Зав. кафедри \dotfill Ю.В.~Крак}}\hspace {0.7cm} {\hbox to 6.5 cm { Екзаменатор \dotfill Т.О.~Карнаух}} \end{scriptsize}}
%end
\newpage
%begin
{{\begin {scriptsize}\par
{ \center  Київський національний університет імені Тараса Шевченка \par}
{\parbox {5 cm}{Освітня програма \hfill}{\parbox {7 cm}{Інформатика\hfill}}\parbox {6 cm}{\hfill Семестр 2}}\par
{\parbox {5 cm} {Навчальний предмет \hfill}\parbox {7 cm}{Програмування \hfill}\parbox {6cm}{\hfill}\par}\vspace{-15pt} \end{scriptsize}}{\center Екзаменаційний білет \No \textbf { 65 }\par}}
{%beginitem
1. Вкажіть основну \textbf{семантичну} відмінність між конструктором копії та конструктором переміщення.%enditem
\par
\vspace{2mm}
%beginitem
2. Максимально змістовно повно сформулюйте назву оголошеного в наведеному фрагменті методу класу.
\vspace{-3mm}
\begin{verbatim}class A{
   ...
   ~A();
};\end{verbatim}
\vspace{-3mm}
%enditem
\par
\vspace{2mm}
%beginitem
3. Вкажіть, що не так з наведеним кодом.
\vspace{-3mm}
\begin{verbatim}int *pa=new int(2), *pb=&pa; delete pa;\end{verbatim}
\vspace{-3mm}
%enditem
\par
\vspace{2mm}
%beginitem
4. За наявності означень
\vspace{-3mm}
\begin{verbatim}
class A{ };
class B:public A{ };
class C:public B{ };\end{verbatim}
\vspace{-3mm}
вкажіть послідовність викликів деструкторів за виконання коду \code{\{C c;\}} 

(Якщо деструктор викликається кілька разів поспіль, то має зазначатися кожен його виклик, наприклад, $\sim$A, $\sim$A, $\sim$A.)
%enditem

\vspace{2mm}
%beginitem
5. Яку часову оцінку складності має алгоритм сортування купою (він же деревом, пірамідою)?
%enditem
\par
\vspace{2mm}
%beginitem
6. Модифікувати клас 

\qquad\code{class Session\{ int DM, MA, Alg, Prog;\};} 


додавши до нього оператор перетворення на тип \code{bool}: об'єкт класу є істинним тоді і тільки тоді, коли значення всіх атрибутів не менші за 60. Означення оператора теж має бути записане.%enditem
\par
\vspace{2mm}
%beginitem
7. 
Реалізувати операцію злиття множин з повтореннями. Множина з повтореннями задається переліком своїх елементів, записаних за неспаданням у текстовому файлі через білі символи. Результат у такому ж самому форматі (і теж за неспаданням) має бути записаний у текстовий файл. Питома функція має отримувати три шляхи до файлів. Розмір вхідних множин такий, що вони не можуть бути завантажені повністю в пам’ять програми.

Крім правильності та відповідності умові оцінюється якість коду, наявність витоків ресурсів та коректність роботи з пам'яттю. 

%enditem
}
{\smallskip\smallskip \begin {scriptsize} Затверджено на засіданні кафедри теоретичної кібернетики, протокол \No 4  від   25.01.2024 р.\par\smallskip\smallskip\smallskip\smallskip
{\parbox {6.5 cm} {Зав. кафедри \dotfill Ю.В.~Крак}}\hspace {0.7cm} {\hbox to 6.5 cm { Екзаменатор \dotfill Т.О.~Карнаух}} \end{scriptsize}}
%end
\newpage
%begin
{{\begin {scriptsize}\par
{ \center  Київський національний університет імені Тараса Шевченка \par}
{\parbox {5 cm}{Освітня програма \hfill}{\parbox {7 cm}{Інформатика\hfill}}\parbox {6 cm}{\hfill Семестр 2}}\par
{\parbox {5 cm} {Навчальний предмет \hfill}\parbox {7 cm}{Програмування \hfill}\parbox {6cm}{\hfill}\par}\vspace{-15pt} \end{scriptsize}}{\center Екзаменаційний білет \No \textbf { 66 }\par}}
{%beginitem
1. Вкажіть основну \textbf{синтаксичну} відмінність між конструктором копії та конструктором переміщення.%enditem
\par
\vspace{2mm}
%beginitem
2. Максимально змістовно повно сформулюйте назву оголошеного в наведеному фрагменті методу класу.
\vspace{-3mm}
\begin{verbatim}class A{
   ...
   A& operator =(A&&);
};\end{verbatim}
\vspace{-3mm}
%enditem
\par
\vspace{2mm}
%beginitem
3. Вкажіть, що не так з наведеним кодом.
\vspace{-3mm}
\begin{verbatim}int *pa=new int(2), *pb=pa; delete pb; *pa=3;\end{verbatim}
\vspace{-3mm}
%enditem
\par
\vspace{2mm}
%beginitem
4. За наявності означень
\vspace{-3mm}
\begin{verbatim}
class A{ };
class B{ };
class C:public A{B b[2];};\end{verbatim}
\vspace{-3mm}
вкажіть послідовність викликів деструкторів за виконання коду \code{\{C c;\}} 

(Якщо деструктор викликається кілька разів поспіль, то має зазначатися кожен його виклик, наприклад, $\sim$A, $\sim$A, $\sim$A.)
%enditem

\vspace{2mm}
%beginitem
5. Яку часову оцінку складності має алгоритм сортування вибором?
%enditem
\par
\vspace{2mm}
%beginitem
6. Модифікувати клас 

\qquad\code{class Session\{ int DM, MA, Alg, Prog;\};} 


додавши до нього оператор перетворення на тип \code{bool}: об'єкт класу є істинним тоді і тільки тоді, коли значення всіх атрибутів не менші за 60. Означення оператора теж має бути записане.%enditem
\par
\vspace{2mm}
%beginitem
7. 
Реалізувати операцію перетину множин з повтореннями. Множина з повтореннями задається переліком своїх елементів, записаних за неспаданням у текстовому файлі через білі символи. Результат у такому ж самому форматі (і теж за неспаданням) має бути записаний у текстовий файл. Питома функція має отримувати три шляхи до файлів. Розмір вхідних множин такий, що вони не можуть бути завантажені повністю в пам’ять програми.

Крім правильності та відповідності умові оцінюється якість коду, наявність витоків ресурсів та коректність роботи з пам'яттю. 

%enditem
}
{\smallskip\smallskip \begin {scriptsize} Затверджено на засіданні кафедри теоретичної кібернетики, протокол \No 4  від   25.01.2024 р.\par\smallskip\smallskip\smallskip\smallskip
{\parbox {6.5 cm} {Зав. кафедри \dotfill Ю.В.~Крак}}\hspace {0.7cm} {\hbox to 6.5 cm { Екзаменатор \dotfill Т.О.~Карнаух}} \end{scriptsize}}
%end
\newpage
%begin
{{\begin {scriptsize}\par
{ \center  Київський національний університет імені Тараса Шевченка \par}
{\parbox {5 cm}{Освітня програма \hfill}{\parbox {7 cm}{Інформатика\hfill}}\parbox {6 cm}{\hfill Семестр 2}}\par
{\parbox {5 cm} {Навчальний предмет \hfill}\parbox {7 cm}{Програмування \hfill}\parbox {6cm}{\hfill}\par}\vspace{-15pt} \end{scriptsize}}{\center Екзаменаційний білет \No \textbf { 67 }\par}}
{%beginitem
1. Вкажіть основну відмінність між віртуальними методами та тими, що не є віртуальними.%enditem
\par
\vspace{2mm}
%beginitem
2. Максимально змістовно повно сформулюйте назву оголошеного в наведеному фрагменті методу класу.
\vspace{-3mm}
\begin{verbatim}class A{
   ...
   ~A();
};\end{verbatim}
\vspace{-3mm}
%enditem
\par
\vspace{2mm}
%beginitem
3. Вкажіть, що не так з наведеним кодом.
\vspace{-3mm}
\begin{verbatim}int *pa=new int(2), *pb=*pa; delete pa;\end{verbatim}
\vspace{-3mm}
%enditem
\par
\vspace{2mm}
%beginitem
4. За наявності означень
\vspace{-3mm}
\begin{verbatim}
class A{ };
class B{ };
class C{B b;A a;};\end{verbatim}
\vspace{-3mm}
вкажіть послідовність викликів деструкторів за виконання коду \code{\{C c;\}} 

(Якщо деструктор викликається кілька разів поспіль, то має зазначатися кожен його виклик, наприклад, $\sim$A, $\sim$A, $\sim$A.)
%enditem

\vspace{2mm}
%beginitem
5. Яку часову оцінку складності має алгоритм сортування вставками?
%enditem
\par
\vspace{2mm}
%beginitem
6. Модифікувати клас точок простору 

\qquad\code{class Point\{ double x, y, z;\};}


додавши до нього оператор перетворення на тип \code{bool}: точка є істинною тоді і тільки тоді, коли містить хоча б одну ненульову координату. Означення оператора теж має бути записане.%enditem
\par
\vspace{2mm}
%beginitem
7. 
Задано клас A (структуру, оператори < та == десь коректно означено) та клас вузла списку Node.

struct A\{

\quad  якісь коректні означення 

\quad  bool operator<(const A\&); 

\quad  bool operator==(const A\&); 

\quad A\& operator =(const A\&)=delete; 

\quad  A\& operator =(A\&\&)=delete;\};

struct Node\{Node *prev; Node *next; A a;\};

Написати функцію, яка вилучає зі списку всі вузли, елемент в яких більший елемента в попередньому вузлі списку або дорівнює елементу в наступному вузлі. Список є циклічним та задається вказівником на перший вузол списку. Наприклад, зі списку, що містить послідовність 8, 10, 9, 4, мають бути вилучені вузли, що містять 8 (бо список циклічний і $8>4$) та 10. У цьому прикладі 9 залишається в списку.

Крім правильності та відповідності умові оцінюється якість коду, наявність витоків ресурсів та коректність роботи з пам'яттю. 

%enditem
}
{\smallskip\smallskip \begin {scriptsize} Затверджено на засіданні кафедри теоретичної кібернетики, протокол \No 4  від   25.01.2024 р.\par\smallskip\smallskip\smallskip\smallskip
{\parbox {6.5 cm} {Зав. кафедри \dotfill Ю.В.~Крак}}\hspace {0.7cm} {\hbox to 6.5 cm { Екзаменатор \dotfill Т.О.~Карнаух}} \end{scriptsize}}
%end
\newpage
%begin
{{\begin {scriptsize}\par
{ \center  Київський національний університет імені Тараса Шевченка \par}
{\parbox {5 cm}{Освітня програма \hfill}{\parbox {7 cm}{Інформатика\hfill}}\parbox {6 cm}{\hfill Семестр 2}}\par
{\parbox {5 cm} {Навчальний предмет \hfill}\parbox {7 cm}{Програмування \hfill}\parbox {6cm}{\hfill}\par}\vspace{-15pt} \end{scriptsize}}{\center Екзаменаційний білет \No \textbf { 68 }\par}}
{%beginitem
1. Що важливо забезпечити при роботі з динамічними змінними?%enditem
\par
\vspace{2mm}
%beginitem
2. Максимально змістовно повно сформулюйте назву оголошеного в наведеному фрагменті методу класу.
\vspace{-3mm}
\begin{verbatim}class A{
   ...
   A(const A&);
};\end{verbatim}
\vspace{-3mm}
%enditem
\par
\vspace{2mm}
%beginitem
3. Вкажіть, що не так з наведеним кодом.
\vspace{-3mm}
\begin{verbatim}int **ppi=new int*(new int(3)); delete ppi;\end{verbatim}
\vspace{-3mm}
%enditem
\par
\vspace{2mm}
%beginitem
4. За наявності означень
\vspace{-3mm}
\begin{verbatim}
class A{ };
class B{ };
class C:public A{B b;A a;};\end{verbatim}
\vspace{-3mm}
вкажіть послідовність викликів деструкторів за виконання коду \code{\{C c;\}} 

(Якщо деструктор викликається кілька разів поспіль, то має зазначатися кожен його виклик, наприклад, $\sim$A, $\sim$A, $\sim$A.)
%enditem

\vspace{2mm}
%beginitem
5. Яку часову оцінку складності має алгоритм сортування бульбашкою?
%enditem
\par
\vspace{2mm}
%beginitem
6. Модифікувати клас 

\qquad\code{class Session\{ int DM, MA, Alg, Prog;\};} 


додавши до нього оператор порівняння \code{<}: один об'єкт є меншим за інший тоді і тільки тоді, коли значення всіх атрибутів об'єктів не менші за 60 і сума атрибутів першого менша за суму атрибутів другого. Означення оператора теж має бути записане.%enditem
\par
\vspace{2mm}
%beginitem
7. 
Реалізувати операцію різниці множин з повтореннями. Множина з повтореннями задається переліком своїх елементів, записаних за неспаданням у текстовому файлі через білі символи. Результат у такому ж самому форматі (і теж за неспаданням) має бути записаний у текстовий файл. Питома функція має отримувати три шляхи до файлів. Розмір вхідних множин такий, що вони не можуть бути завантажені повністю в пам’ять програми.

Крім правильності та відповідності умові оцінюється якість коду, наявність витоків ресурсів та коректність роботи з пам'яттю. 

%enditem
}
{\smallskip\smallskip \begin {scriptsize} Затверджено на засіданні кафедри теоретичної кібернетики, протокол \No 4  від   25.01.2024 р.\par\smallskip\smallskip\smallskip\smallskip
{\parbox {6.5 cm} {Зав. кафедри \dotfill Ю.В.~Крак}}\hspace {0.7cm} {\hbox to 6.5 cm { Екзаменатор \dotfill Т.О.~Карнаух}} \end{scriptsize}}
%end
\newpage
%begin
{{\begin {scriptsize}\par
{ \center  Київський національний університет імені Тараса Шевченка \par}
{\parbox {5 cm}{Освітня програма \hfill}{\parbox {7 cm}{Інформатика\hfill}}\parbox {6 cm}{\hfill Семестр 2}}\par
{\parbox {5 cm} {Навчальний предмет \hfill}\parbox {7 cm}{Програмування \hfill}\parbox {6cm}{\hfill}\par}\vspace{-15pt} \end{scriptsize}}{\center Екзаменаційний білет \No \textbf { 69 }\par}}
{%beginitem
1. Вкажіть основну \textbf{семантичну} відмінність між конструктором копії та конструктором переміщення.%enditem
\par
\vspace{2mm}
%beginitem
2. Максимально змістовно повно сформулюйте назву оголошеного в наведеному фрагменті методу класу.
\vspace{-3mm}
\begin{verbatim}class A{
   ...
   A(A&&);
};\end{verbatim}
\vspace{-3mm}
%enditem
\par
\vspace{2mm}
%beginitem
3. Вкажіть, що не так з наведеним кодом.
\vspace{-3mm}
\begin{verbatim}int *pa=new int(2), *pb=new int(3); pa=pb;\end{verbatim}
\vspace{-3mm}
%enditem
\par
\vspace{2mm}
%beginitem
4. За наявності означень
\vspace{-3mm}
\begin{verbatim}
class A{ };
class B{ };
class C{B b;A a;};\end{verbatim}
\vspace{-3mm}
вкажіть послідовність викликів конструкторів за виконання коду \code{\{C c;\}} 

(Якщо конструктор викликається кілька разів поспіль, то має зазначатися кожен його виклик, наприклад, A, A, A.)
%enditem

\vspace{2mm}
%beginitem
5. Яку часову оцінку складності має алгоритм швидкого сортування?
%enditem
\par
\vspace{2mm}
%beginitem
6. Модифікувати клас 

\qquad\code{class Session\{ int DM, MA, Alg, Prog;\};} 


додавши до нього оператор перетворення на тип \code{int}: якщо хоча б один атрибут менший за 60, то має повертатись 0, інакше -- сума атрибутів.%enditem
\par
\vspace{2mm}
%beginitem
7. 
Задано клас A (структуру, оператори > та == десь коректно означено) та клас вузла списку Node.

struct A\{

\quad  якісь коректні означення 

\quad  bool operator>(const A\&); 

\quad  bool operator==(const A\&); 

\quad A\& operator =(const A\&)=delete; 

\quad  A\& operator =(A\&\&)=delete;\};

struct Node\{Node *next; A a;\};

 Написати функцію, яка вилучає зі списку всі вузли, елемент в яких менший елемента в попередньому вузлі списку або дорівнює елементу в наступному вузлі. Список задається вказівником на перший вузол списку. Наприклад, зі списку, що містить послідовність 8, 5, 6, 9, 4, мають бути вилучені вузли, що містять 5 та 4. У цьому прикладі 6 залишається в списку.

Крім правильності та відповідності умові оцінюється якість коду, наявність витоків ресурсів та коректність роботи з пам'яттю. 

%enditem
}
{\smallskip\smallskip \begin {scriptsize} Затверджено на засіданні кафедри теоретичної кібернетики, протокол \No 4  від   25.01.2024 р.\par\smallskip\smallskip\smallskip\smallskip
{\parbox {6.5 cm} {Зав. кафедри \dotfill Ю.В.~Крак}}\hspace {0.7cm} {\hbox to 6.5 cm { Екзаменатор \dotfill Т.О.~Карнаух}} \end{scriptsize}}
%end
\newpage
%begin
{{\begin {scriptsize}\par
{ \center  Київський національний університет імені Тараса Шевченка \par}
{\parbox {5 cm}{Освітня програма \hfill}{\parbox {7 cm}{Інформатика\hfill}}\parbox {6 cm}{\hfill Семестр 2}}\par
{\parbox {5 cm} {Навчальний предмет \hfill}\parbox {7 cm}{Програмування \hfill}\parbox {6cm}{\hfill}\par}\vspace{-15pt} \end{scriptsize}}{\center Екзаменаційний білет \No \textbf { 70 }\par}}
{%beginitem
1. Вкажіть основну відмінність між класами (class) та структурами (struct).%enditem
\par
\vspace{2mm}
%beginitem
2. Максимально змістовно повно сформулюйте назву оголошеного в наведеному фрагменті методу класу.
\vspace{-3mm}
\begin{verbatim}class A{
   ...
   A();
};\end{verbatim}
\vspace{-3mm}
%enditem
\par
\vspace{2mm}
%beginitem
3. Вкажіть, що не так з наведеним кодом.
\vspace{-3mm}
\begin{verbatim}int *pa=new int(2), *pb=pa; delete pb; delete pa;\end{verbatim}
\vspace{-3mm}
%enditem
\par
\vspace{2mm}
%beginitem
4. За наявності означень
\vspace{-3mm}
\begin{verbatim}
class A{ };
class B{ };
class C:public A{B b;A a;};\end{verbatim}
\vspace{-3mm}
вкажіть послідовність викликів конструкторів за виконання коду \code{\{C c;\}} 

(Якщо конструктор викликається кілька разів поспіль, то має зазначатися кожен його виклик, наприклад, A, A, A.)
%enditem

\vspace{2mm}
%beginitem
5. Яку часову оцінку складності має алгоритм швидкого сортування?
%enditem
\par
\vspace{2mm}
%beginitem
6. Модифікувати клас точок простору 

\qquad\code{class Point\{ double x, y, z;\};}


додавши до нього оператор перетворення на рядок типу \code{std::string}, що містить рядкове зображення точки простору (у круглих дужках координати через кому), наприклад, (1.01, 2.35, 0.5), та його означення.%enditem
\par
\vspace{2mm}
%beginitem
7. 
Реалізувати операцію симетричної різниці множин з повтореннями. Множина з повтореннями задається переліком своїх елементів, записаних за неспаданням у текстовому файлі через білі символи. Результат у такому ж самому форматі (і теж за неспаданням) має бути записаний у текстовий файл. Питома функція має отримувати три шляхи до файлів. Розмір вхідних множин такий, що вони не можуть бути завантажені повністю в пам’ять програми.

Крім правильності та відповідності умові оцінюється якість коду, наявність витоків ресурсів та коректність роботи з пам'яттю. 

%enditem
}
{\smallskip\smallskip \begin {scriptsize} Затверджено на засіданні кафедри теоретичної кібернетики, протокол \No 4  від   25.01.2024 р.\par\smallskip\smallskip\smallskip\smallskip
{\parbox {6.5 cm} {Зав. кафедри \dotfill Ю.В.~Крак}}\hspace {0.7cm} {\hbox to 6.5 cm { Екзаменатор \dotfill Т.О.~Карнаух}} \end{scriptsize}}
%end
\newpage
%begin
{{\begin {scriptsize}\par
{ \center  Київський національний університет імені Тараса Шевченка \par}
{\parbox {5 cm}{Освітня програма \hfill}{\parbox {7 cm}{Інформатика\hfill}}\parbox {6 cm}{\hfill Семестр 2}}\par
{\parbox {5 cm} {Навчальний предмет \hfill}\parbox {7 cm}{Програмування \hfill}\parbox {6cm}{\hfill}\par}\vspace{-15pt} \end{scriptsize}}{\center Екзаменаційний білет \No \textbf { 71 }\par}}
{%beginitem
1. Вкажіть основну відмінність між відкритими (public) та прихованими (private) методами класу.%enditem
\par
\vspace{2mm}
%beginitem
2. Максимально змістовно повно сформулюйте назву оголошеного в наведеному фрагменті методу класу.
\vspace{-3mm}
\begin{verbatim}class A{
   ...
   A& operator =(A&&);
};\end{verbatim}
\vspace{-3mm}
%enditem
\par
\vspace{2mm}
%beginitem
3. Вкажіть, що не так з наведеним кодом.
\vspace{-3mm}
\begin{verbatim}int *pa=new int(2), *pb=new int(3); pa=pb;\end{verbatim}
\vspace{-3mm}
%enditem
\par
\vspace{2mm}
%beginitem
4. За наявності означень
\vspace{-3mm}
\begin{verbatim}
class A{ };
class B:public A{ };
class C:public A{B b;};\end{verbatim}
\vspace{-3mm}
вкажіть послідовність викликів деструкторів за виконання коду \code{\{C c;\}} 

(Якщо деструктор викликається кілька разів поспіль, то має зазначатися кожен його виклик, наприклад, $\sim$A, $\sim$A, $\sim$A.)
%enditem

\vspace{2mm}
%beginitem
5. Яку часову оцінку складності має алгоритм сортування вибором?
%enditem
\par
\vspace{2mm}
%beginitem
6. Модифікувати клас 

\qquad\code{class Session\{ int DM, MA, Alg, Prog;\};} 


додавши до нього оператор порівняння \code{<}: один об'єкт є меншим за інший тоді і тільки тоді, коли значення всіх атрибутів об'єктів не менші за 60 і сума атрибутів першого менша за суму атрибутів другого. Означення оператора теж має бути записане.%enditem
\par
\vspace{2mm}
%beginitem
7. 
Задано клас A (структуру, оператори < та == десь коректно означено) та клас вузла списку Node.

struct A\{

\quad  якісь коректні означення 

\quad  bool operator<(const A\&); 

\quad A\& operator =(const A\&)=delete; 

\quad  A\& operator =(A\&\&)=delete;\};

struct Node\{Node *prev; Node *next; A a;\};

Написати функцію, яка вилучає зі списку всі вузли, елемент в яких менший елемента в попередньому вузлі списку або дорівнює елементу в наступному вузлі. Список задається вказівником на перший вузол списку. Наприклад, зі списку, що містить послідовність 8, 5, 6, 9, 4, мають бути вилучені вузли, що містять 5 та 4. У цьому прикладі 6 залишається в списку.

Крім правильності та відповідності умові оцінюється якість коду, наявність витоків ресурсів та коректність роботи з пам'яттю. 

%enditem
}
{\smallskip\smallskip \begin {scriptsize} Затверджено на засіданні кафедри теоретичної кібернетики, протокол \No 4  від   25.01.2024 р.\par\smallskip\smallskip\smallskip\smallskip
{\parbox {6.5 cm} {Зав. кафедри \dotfill Ю.В.~Крак}}\hspace {0.7cm} {\hbox to 6.5 cm { Екзаменатор \dotfill Т.О.~Карнаух}} \end{scriptsize}}
%end
\newpage
%begin
{{\begin {scriptsize}\par
{ \center  Київський національний університет імені Тараса Шевченка \par}
{\parbox {5 cm}{Освітня програма \hfill}{\parbox {7 cm}{Інформатика\hfill}}\parbox {6 cm}{\hfill Семестр 2}}\par
{\parbox {5 cm} {Навчальний предмет \hfill}\parbox {7 cm}{Програмування \hfill}\parbox {6cm}{\hfill}\par}\vspace{-15pt} \end{scriptsize}}{\center Екзаменаційний білет \No \textbf { 72 }\par}}
{%beginitem
1. Вкажіть основну \textbf{синтаксичну} відмінність між конструктором копії та конструктором переміщення.%enditem
\par
\vspace{2mm}
%beginitem
2. Максимально змістовно повно сформулюйте назву оголошеного в наведеному фрагменті методу класу.
\vspace{-3mm}
\begin{verbatim}class A{
   ...
   ~A();
};\end{verbatim}
\vspace{-3mm}
%enditem
\par
\vspace{2mm}
%beginitem
3. Вкажіть, що не так з наведеним кодом.
\vspace{-3mm}
\begin{verbatim}int a, *p=&a; delete a;\end{verbatim}
\vspace{-3mm}
%enditem
\par
\vspace{2mm}
%beginitem
4. За наявності означень
\vspace{-3mm}
\begin{verbatim}
class A{ };
class B{ };
class C:public A{B b[2];};\end{verbatim}
\vspace{-3mm}
вкажіть послідовність викликів конструкторів за виконання коду \code{\{C c;\}} 

(Якщо конструктор викликається кілька разів поспіль, то має зазначатися кожен його виклик, наприклад, A, A, A.)
%enditem

\vspace{2mm}
%beginitem
5. Яку часову оцінку складності має алгоритм сортування купою (він же деревом, пірамідою)?
%enditem
\par
\vspace{2mm}
%beginitem
6. Модифікувати клас 

\qquad\code{class Session\{ int DM, MA, Alg, Prog;\};} 


додавши до нього оператор перетворення на тип \code{int}: якщо хоча б один атрибут менший за 60, то має повертатись 0, інакше -- сума атрибутів.%enditem
\par
\vspace{2mm}
%beginitem
7. 
Задано клас A (структуру, оператори < та == десь коректно означено) та клас вузла списку Node.

struct A\{

\quad  якісь коректні означення 

\quad  bool operator<(const A\&); 

\quad  bool operator==(const A\&); 

\quad A\& operator =(const A\&)=delete; 

\quad  A\& operator =(A\&\&)=delete;\};

struct Node\{Node *next; A a;\};

 Написати функцію, яка вилучає зі списку всі вузли, елемент в яких більший елемента в попередньому вузлі списку або дорівнює елементу в наступному вузлі. Список є циклічним та задається вказівником на перший вузол списку. Наприклад, зі списку, що містить послідовність 8, 10, 9, 4, мають бути вилучені вузли, що містять 8 (бо список циклічний і $8>4$) та 10. У цьому прикладі 9 залишається в списку.

Крім правильності та відповідності умові оцінюється якість коду, наявність витоків ресурсів та коректність роботи з пам'яттю. 

%enditem
}
{\smallskip\smallskip \begin {scriptsize} Затверджено на засіданні кафедри теоретичної кібернетики, протокол \No 4  від   25.01.2024 р.\par\smallskip\smallskip\smallskip\smallskip
{\parbox {6.5 cm} {Зав. кафедри \dotfill Ю.В.~Крак}}\hspace {0.7cm} {\hbox to 6.5 cm { Екзаменатор \dotfill Т.О.~Карнаух}} \end{scriptsize}}
%end
\newpage
%begin
{{\begin {scriptsize}\par
{ \center  Київський національний університет імені Тараса Шевченка \par}
{\parbox {5 cm}{Освітня програма \hfill}{\parbox {7 cm}{Інформатика\hfill}}\parbox {6 cm}{\hfill Семестр 2}}\par
{\parbox {5 cm} {Навчальний предмет \hfill}\parbox {7 cm}{Програмування \hfill}\parbox {6cm}{\hfill}\par}\vspace{-15pt} \end{scriptsize}}{\center Екзаменаційний білет \No \textbf { 73 }\par}}
{%beginitem
1. Вкажіть основну відмінність між віртуальними методами та тими, що не є віртуальними.%enditem
\par
\vspace{2mm}
%beginitem
2. Максимально змістовно повно сформулюйте назву оголошеного в наведеному фрагменті методу класу.
\vspace{-3mm}
\begin{verbatim}class A{
   ...
   A(const A&);
};\end{verbatim}
\vspace{-3mm}
%enditem
\par
\vspace{2mm}
%beginitem
3. Вкажіть, що не так з наведеним кодом.
\vspace{-3mm}
\begin{verbatim}int *pa=new int(2), *pb=pa; delete pb; *pa=3;\end{verbatim}
\vspace{-3mm}
%enditem
\par
\vspace{2mm}
%beginitem
4. За наявності означень
\vspace{-3mm}
\begin{verbatim}
class A{ };
class B{ };
class C:public A{B b[2];};\end{verbatim}
\vspace{-3mm}
вкажіть послідовність викликів конструкторів за виконання коду \code{\{C c;\}} 

(Якщо конструктор викликається кілька разів поспіль, то має зазначатися кожен його виклик, наприклад, A, A, A.)
%enditem

\vspace{2mm}
%beginitem
5. Яку часову оцінку складності має алгоритм сортування вставками?
%enditem
\par
\vspace{2mm}
%beginitem
6. Модифікувати клас точок простору 

\qquad\code{class Point\{ double x, y, z;\};}


додавши до нього оператор перетворення на тип \code{bool}: точка є істинною тоді і тільки тоді, коли містить хоча б одну ненульову координату. Означення оператора теж має бути записане.%enditem
\par
\vspace{2mm}
%beginitem
7. 
Реалізувати операцію перетину множин з повтореннями. Множина з повтореннями задається переліком своїх елементів, записаних за неспаданням у текстовому файлі через білі символи. Результат у такому ж самому форматі (і теж за неспаданням) має бути записаний у текстовий файл. Питома функція має отримувати три шляхи до файлів. Розмір вхідних множин такий, що вони не можуть бути завантажені повністю в пам’ять програми.

Крім правильності та відповідності умові оцінюється якість коду, наявність витоків ресурсів та коректність роботи з пам'яттю. 

%enditem
}
{\smallskip\smallskip \begin {scriptsize} Затверджено на засіданні кафедри теоретичної кібернетики, протокол \No 4  від   25.01.2024 р.\par\smallskip\smallskip\smallskip\smallskip
{\parbox {6.5 cm} {Зав. кафедри \dotfill Ю.В.~Крак}}\hspace {0.7cm} {\hbox to 6.5 cm { Екзаменатор \dotfill Т.О.~Карнаух}} \end{scriptsize}}
%end
\newpage
%begin
{{\begin {scriptsize}\par
{ \center  Київський національний університет імені Тараса Шевченка \par}
{\parbox {5 cm}{Освітня програма \hfill}{\parbox {7 cm}{Інформатика\hfill}}\parbox {6 cm}{\hfill Семестр 2}}\par
{\parbox {5 cm} {Навчальний предмет \hfill}\parbox {7 cm}{Програмування \hfill}\parbox {6cm}{\hfill}\par}\vspace{-15pt} \end{scriptsize}}{\center Екзаменаційний білет \No \textbf { 74 }\par}}
{%beginitem
1. Що важливо забезпечити при роботі з динамічними змінними?%enditem
\par
\vspace{2mm}
%beginitem
2. Максимально змістовно повно сформулюйте назву оголошеного в наведеному фрагменті методу класу.
\vspace{-3mm}
\begin{verbatim}class A{
   ...
   A();
};\end{verbatim}
\vspace{-3mm}
%enditem
\par
\vspace{2mm}
%beginitem
3. Вкажіть, що не так з наведеним кодом.
\vspace{-3mm}
\begin{verbatim}int *pa=new int(2), *pb=pa; delete pb; delete pa;\end{verbatim}
\vspace{-3mm}
%enditem
\par
\vspace{2mm}
%beginitem
4. За наявності означень
\vspace{-3mm}
\begin{verbatim}
class A{ };
class B{ };
class C:public A{B b;};\end{verbatim}
\vspace{-3mm}
вкажіть послідовність викликів деструкторів за виконання коду \code{\{C c;\}} 

(Якщо деструктор викликається кілька разів поспіль, то має зазначатися кожен його виклик, наприклад, $\sim$A, $\sim$A, $\sim$A.)
%enditem

\vspace{2mm}
%beginitem
5. Яку часову оцінку складності має алгоритм швидкого сортування?
%enditem
\par
\vspace{2mm}
%beginitem
6. Модифікувати клас точок простору 

\qquad\code{class Point\{ double x, y, z;\};}


додавши до нього оператор перетворення на рядок типу \code{std::string}, що містить рядкове зображення точки простору (у круглих дужках координати через кому), наприклад, (1.01, 2.35, 0.5), та його означення.%enditem
\par
\vspace{2mm}
%beginitem
7. 
Задано клас A (структуру, оператори < та == десь коректно означено) та клас вузла списку Node.

struct A\{

\quad  якісь коректні означення 

\quad  bool operator<(const A\&); 

\quad A\& operator =(const A\&)=delete; 

\quad  A\& operator =(A\&\&)=delete;\};

struct Node\{Node *prev; Node *next; A a;\};

Написати функцію, яка вилучає зі списку всі вузли, елемент в яких менший елемента в попередньому вузлі списку або дорівнює елементу в наступному вузлі. Список задається вказівником на перший вузол списку. Наприклад, зі списку, що містить послідовність 8, 5, 6, 9, 4, мають бути вилучені вузли, що містять 5 та 4. У цьому прикладі 6 залишається в списку.

Крім правильності та відповідності умові оцінюється якість коду, наявність витоків ресурсів та коректність роботи з пам'яттю. 

%enditem
}
{\smallskip\smallskip \begin {scriptsize} Затверджено на засіданні кафедри теоретичної кібернетики, протокол \No 4  від   25.01.2024 р.\par\smallskip\smallskip\smallskip\smallskip
{\parbox {6.5 cm} {Зав. кафедри \dotfill Ю.В.~Крак}}\hspace {0.7cm} {\hbox to 6.5 cm { Екзаменатор \dotfill Т.О.~Карнаух}} \end{scriptsize}}
%end
\newpage
%begin
{{\begin {scriptsize}\par
{ \center  Київський національний університет імені Тараса Шевченка \par}
{\parbox {5 cm}{Освітня програма \hfill}{\parbox {7 cm}{Інформатика\hfill}}\parbox {6 cm}{\hfill Семестр 2}}\par
{\parbox {5 cm} {Навчальний предмет \hfill}\parbox {7 cm}{Програмування \hfill}\parbox {6cm}{\hfill}\par}\vspace{-15pt} \end{scriptsize}}{\center Екзаменаційний білет \No \textbf { 75 }\par}}
{%beginitem
1. Вкажіть основну \textbf{семантичну} відмінність між конструктором копії та конструктором переміщення.%enditem
\par
\vspace{2mm}
%beginitem
2. Максимально змістовно повно сформулюйте назву оголошеного в наведеному фрагменті методу класу.
\vspace{-3mm}
\begin{verbatim}class A{
   ...
   A(A&&);
};\end{verbatim}
\vspace{-3mm}
%enditem
\par
\vspace{2mm}
%beginitem
3. Вкажіть, що не так з наведеним кодом.
\vspace{-3mm}
\begin{verbatim}int *pa=new int(2), *pb=&pa; delete pa;\end{verbatim}
\vspace{-3mm}
%enditem
\par
\vspace{2mm}
%beginitem
4. За наявності означень
\vspace{-3mm}
\begin{verbatim}
class A{ };
class B:public A{ };
class C:public B{ };\end{verbatim}
\vspace{-3mm}
вкажіть послідовність викликів деструкторів за виконання коду \code{\{C c;\}} 

(Якщо деструктор викликається кілька разів поспіль, то має зазначатися кожен його виклик, наприклад, $\sim$A, $\sim$A, $\sim$A.)
%enditem

\vspace{2mm}
%beginitem
5. Яку часову оцінку складності має алгоритм швидкого сортування?
%enditem
\par
\vspace{2mm}
%beginitem
6. Модифікувати клас 

\qquad\code{class Session\{ int DM, MA, Alg, Prog;\};} 


додавши до нього оператор перетворення на тип \code{bool}: об'єкт класу є істинним тоді і тільки тоді, коли значення всіх атрибутів не менші за 60. Означення оператора теж має бути записане.%enditem
\par
\vspace{2mm}
%beginitem
7. 
Реалізувати операцію симетричної різниці множин з повтореннями. Множина з повтореннями задається переліком своїх елементів, записаних за неспаданням у текстовому файлі через білі символи. Результат у такому ж самому форматі (і теж за неспаданням) має бути записаний у текстовий файл. Питома функція має отримувати три шляхи до файлів. Розмір вхідних множин такий, що вони не можуть бути завантажені повністю в пам’ять програми.

Крім правильності та відповідності умові оцінюється якість коду, наявність витоків ресурсів та коректність роботи з пам'яттю. 

%enditem
}
{\smallskip\smallskip \begin {scriptsize} Затверджено на засіданні кафедри теоретичної кібернетики, протокол \No 4  від   25.01.2024 р.\par\smallskip\smallskip\smallskip\smallskip
{\parbox {6.5 cm} {Зав. кафедри \dotfill Ю.В.~Крак}}\hspace {0.7cm} {\hbox to 6.5 cm { Екзаменатор \dotfill Т.О.~Карнаух}} \end{scriptsize}}
%end
\newpage
%begin
{{\begin {scriptsize}\par
{ \center  Київський національний університет імені Тараса Шевченка \par}
{\parbox {5 cm}{Освітня програма \hfill}{\parbox {7 cm}{Інформатика\hfill}}\parbox {6 cm}{\hfill Семестр 2}}\par
{\parbox {5 cm} {Навчальний предмет \hfill}\parbox {7 cm}{Програмування \hfill}\parbox {6cm}{\hfill}\par}\vspace{-15pt} \end{scriptsize}}{\center Екзаменаційний білет \No \textbf { 76 }\par}}
{%beginitem
1. Вкажіть основну відмінність між відкритими (public) та прихованими (private) методами класу.%enditem
\par
\vspace{2mm}
%beginitem
2. Максимально змістовно повно сформулюйте назву оголошеного в наведеному фрагменті методу класу.
\vspace{-3mm}
\begin{verbatim}class A{
   ...
   A(A&&);
};\end{verbatim}
\vspace{-3mm}
%enditem
\par
\vspace{2mm}
%beginitem
3. Вкажіть, що не так з наведеним кодом.
\vspace{-3mm}
\begin{verbatim}int *pa=new int(2), *pb=*pa; delete pa;\end{verbatim}
\vspace{-3mm}
%enditem
\par
\vspace{2mm}
%beginitem
4. За наявності означень
\vspace{-3mm}
\begin{verbatim}
class A{ };
class B:public A{ };
class C:public A{B b;};\end{verbatim}
\vspace{-3mm}
вкажіть послідовність викликів конструкторів за виконання коду \code{\{C c;\}} 

(Якщо конструктор викликається кілька разів поспіль, то має зазначатися кожен його виклик, наприклад, A, A, A.)
%enditem

\vspace{2mm}
%beginitem
5. Яку часову оцінку складності має алгоритм сортування бульбашкою?
%enditem
\par
\vspace{2mm}
%beginitem
6. Модифікувати клас 

\qquad\code{class Session\{ int DM, MA, Alg, Prog;\};} 


додавши до нього оператор перетворення на тип \code{int}: якщо хоча б один атрибут менший за 60, то має повертатись 0, інакше -- сума атрибутів.%enditem
\par
\vspace{2mm}
%beginitem
7. 
Задано клас A (структуру, оператори > та == десь коректно означено) та клас вузла списку Node.

struct A\{

\quad  якісь коректні означення 

\quad  bool operator>(const A\&); 

\quad  bool operator==(const A\&); 

\quad A\& operator =(const A\&)=delete; 

\quad  A\& operator =(A\&\&)=delete;\};

struct Node\{Node *next; A a;\};

 Написати функцію, яка вилучає зі списку всі вузли, елемент в яких менший елемента в попередньому вузлі списку або дорівнює елементу в наступному вузлі. Список задається вказівником на перший вузол списку. Наприклад, зі списку, що містить послідовність 8, 5, 6, 9, 4, мають бути вилучені вузли, що містять 5 та 4. У цьому прикладі 6 залишається в списку.

Крім правильності та відповідності умові оцінюється якість коду, наявність витоків ресурсів та коректність роботи з пам'яттю. 

%enditem
}
{\smallskip\smallskip \begin {scriptsize} Затверджено на засіданні кафедри теоретичної кібернетики, протокол \No 4  від   25.01.2024 р.\par\smallskip\smallskip\smallskip\smallskip
{\parbox {6.5 cm} {Зав. кафедри \dotfill Ю.В.~Крак}}\hspace {0.7cm} {\hbox to 6.5 cm { Екзаменатор \dotfill Т.О.~Карнаух}} \end{scriptsize}}
%end
\newpage
%begin
{{\begin {scriptsize}\par
{ \center  Київський національний університет імені Тараса Шевченка \par}
{\parbox {5 cm}{Освітня програма \hfill}{\parbox {7 cm}{Інформатика\hfill}}\parbox {6 cm}{\hfill Семестр 2}}\par
{\parbox {5 cm} {Навчальний предмет \hfill}\parbox {7 cm}{Програмування \hfill}\parbox {6cm}{\hfill}\par}\vspace{-15pt} \end{scriptsize}}{\center Екзаменаційний білет \No \textbf { 77 }\par}}
{%beginitem
1. Вкажіть основну відмінність між класами (class) та структурами (struct).%enditem
\par
\vspace{2mm}
%beginitem
2. Максимально змістовно повно сформулюйте назву оголошеного в наведеному фрагменті методу класу.
\vspace{-3mm}
\begin{verbatim}class A{
   ...
   A(const A&);
};\end{verbatim}
\vspace{-3mm}
%enditem
\par
\vspace{2mm}
%beginitem
3. Вкажіть, що не так з наведеним кодом.
\vspace{-3mm}
\begin{verbatim}int **ppi=new int*(new int(3)); delete ppi;\end{verbatim}
\vspace{-3mm}
%enditem
\par
\vspace{2mm}
%beginitem
4. За наявності означень
\vspace{-3mm}
\begin{verbatim}
class A{ };
class B:public A{ };
class C:public B{ };\end{verbatim}
\vspace{-3mm}
вкажіть послідовність викликів конструкторів за виконання коду \code{\{C c;\}} 

(Якщо конструктор викликається кілька разів поспіль, то має зазначатися кожен його виклик, наприклад, A, A, A.)
%enditem

\vspace{2mm}
%beginitem
5. Яку часову оцінку складності має алгоритм сортування злиттям?
%enditem
\par
\vspace{2mm}
%beginitem
6. Модифікувати клас точок простору 

\qquad\code{class Point\{ double x, y, z;\};}


додавши до нього оператор перетворення на рядок типу \code{std::string}, що містить рядкове зображення точки простору (у круглих дужках координати через кому), наприклад, (1.01, 2.35, 0.5), та його означення.%enditem
\par
\vspace{2mm}
%beginitem
7. 
Задано клас A (структуру, оператори < та == десь коректно означено) та клас вузла списку Node.

struct A\{

\quad  якісь коректні означення 

\quad  bool operator<(const A\&); 

\quad  bool operator==(const A\&); 

\quad A\& operator =(const A\&)=delete; 

\quad  A\& operator =(A\&\&)=delete;\};

struct Node\{Node *next; A a;\};

 Написати функцію, яка вилучає зі списку всі вузли, елемент в яких більший елемента в попередньому вузлі списку або дорівнює елементу в наступному вузлі. Список є циклічним та задається вказівником на перший вузол списку. Наприклад, зі списку, що містить послідовність 8, 10, 9, 4, мають бути вилучені вузли, що містять 8 (бо список циклічний і $8>4$) та 10. У цьому прикладі 9 залишається в списку.

Крім правильності та відповідності умові оцінюється якість коду, наявність витоків ресурсів та коректність роботи з пам'яттю. 

%enditem
}
{\smallskip\smallskip \begin {scriptsize} Затверджено на засіданні кафедри теоретичної кібернетики, протокол \No 4  від   25.01.2024 р.\par\smallskip\smallskip\smallskip\smallskip
{\parbox {6.5 cm} {Зав. кафедри \dotfill Ю.В.~Крак}}\hspace {0.7cm} {\hbox to 6.5 cm { Екзаменатор \dotfill Т.О.~Карнаух}} \end{scriptsize}}
%end
\newpage
%begin
{{\begin {scriptsize}\par
{ \center  Київський національний університет імені Тараса Шевченка \par}
{\parbox {5 cm}{Освітня програма \hfill}{\parbox {7 cm}{Інформатика\hfill}}\parbox {6 cm}{\hfill Семестр 2}}\par
{\parbox {5 cm} {Навчальний предмет \hfill}\parbox {7 cm}{Програмування \hfill}\parbox {6cm}{\hfill}\par}\vspace{-15pt} \end{scriptsize}}{\center Екзаменаційний білет \No \textbf { 78 }\par}}
{%beginitem
1. Вкажіть основну \textbf{синтаксичну} відмінність між конструктором копії та конструктором переміщення.%enditem
\par
\vspace{2mm}
%beginitem
2. Максимально змістовно повно сформулюйте назву оголошеного в наведеному фрагменті методу класу.
\vspace{-3mm}
\begin{verbatim}class A{
   ...
   ~A();
};\end{verbatim}
\vspace{-3mm}
%enditem
\par
\vspace{2mm}
%beginitem
3. Вкажіть, що не так з наведеним кодом.
\vspace{-3mm}
\begin{verbatim}int *pa=new int(2), *pb=&pa; delete pa;\end{verbatim}
\vspace{-3mm}
%enditem
\par
\vspace{2mm}
%beginitem
4. За наявності означень
\vspace{-3mm}
\begin{verbatim}
class A{ };
class B{ };
class C:public A{B b[2];};\end{verbatim}
\vspace{-3mm}
вкажіть послідовність викликів деструкторів за виконання коду \code{\{C c;\}} 

(Якщо деструктор викликається кілька разів поспіль, то має зазначатися кожен його виклик, наприклад, $\sim$A, $\sim$A, $\sim$A.)
%enditem

\vspace{2mm}
%beginitem
5. Яку часову оцінку складності має алгоритм швидкого сортування?
%enditem
\par
\vspace{2mm}
%beginitem
6. Модифікувати клас 

\qquad\code{class Session\{ int DM, MA, Alg, Prog;\};} 


додавши до нього оператор порівняння \code{<}: один об'єкт є меншим за інший тоді і тільки тоді, коли значення всіх атрибутів об'єктів не менші за 60 і сума атрибутів першого менша за суму атрибутів другого. Означення оператора теж має бути записане.%enditem
\par
\vspace{2mm}
%beginitem
7. 
Реалізувати операцію злиття множин з повтореннями. Множина з повтореннями задається переліком своїх елементів, записаних за неспаданням у текстовому файлі через білі символи. Результат у такому ж самому форматі (і теж за неспаданням) має бути записаний у текстовий файл. Питома функція має отримувати три шляхи до файлів. Розмір вхідних множин такий, що вони не можуть бути завантажені повністю в пам’ять програми.

Крім правильності та відповідності умові оцінюється якість коду, наявність витоків ресурсів та коректність роботи з пам'яттю. 

%enditem
}
{\smallskip\smallskip \begin {scriptsize} Затверджено на засіданні кафедри теоретичної кібернетики, протокол \No 4  від   25.01.2024 р.\par\smallskip\smallskip\smallskip\smallskip
{\parbox {6.5 cm} {Зав. кафедри \dotfill Ю.В.~Крак}}\hspace {0.7cm} {\hbox to 6.5 cm { Екзаменатор \dotfill Т.О.~Карнаух}} \end{scriptsize}}
%end
\newpage
%begin
{{\begin {scriptsize}\par
{ \center  Київський національний університет імені Тараса Шевченка \par}
{\parbox {5 cm}{Освітня програма \hfill}{\parbox {7 cm}{Інформатика\hfill}}\parbox {6 cm}{\hfill Семестр 2}}\par
{\parbox {5 cm} {Навчальний предмет \hfill}\parbox {7 cm}{Програмування \hfill}\parbox {6cm}{\hfill}\par}\vspace{-15pt} \end{scriptsize}}{\center Екзаменаційний білет \No \textbf { 79 }\par}}
{%beginitem
1. Вкажіть основну відмінність між віртуальними методами та тими, що не є віртуальними.%enditem
\par
\vspace{2mm}
%beginitem
2. Максимально змістовно повно сформулюйте назву оголошеного в наведеному фрагменті методу класу.
\vspace{-3mm}
\begin{verbatim}class A{
   ...
   A();
};\end{verbatim}
\vspace{-3mm}
%enditem
\par
\vspace{2mm}
%beginitem
3. Вкажіть, що не так з наведеним кодом.
\vspace{-3mm}
\begin{verbatim}int *pa=new int(2), *pb=pa; delete pb; delete pa;\end{verbatim}
\vspace{-3mm}
%enditem
\par
\vspace{2mm}
%beginitem
4. За наявності означень
\vspace{-3mm}
\begin{verbatim}
class A{ };
class B{ };
class C{B b;A a;};\end{verbatim}
\vspace{-3mm}
вкажіть послідовність викликів конструкторів за виконання коду \code{\{C c;\}} 

(Якщо конструктор викликається кілька разів поспіль, то має зазначатися кожен його виклик, наприклад, A, A, A.)
%enditem

\vspace{2mm}
%beginitem
5. Яку часову оцінку складності має алгоритм сортування злиттям?
%enditem
\par
\vspace{2mm}
%beginitem
6. Модифікувати клас 

\qquad\code{class Session\{ int DM, MA, Alg, Prog;\};} 


додавши до нього оператор перетворення на тип \code{bool}: об'єкт класу є істинним тоді і тільки тоді, коли значення всіх атрибутів не менші за 60. Означення оператора теж має бути записане.%enditem
\par
\vspace{2mm}
%beginitem
7. 
Задано клас A (структуру, оператори < та == десь коректно означено) та клас вузла списку Node.

struct A\{

\quad  якісь коректні означення 

\quad  bool operator<(const A\&); 

\quad  bool operator==(const A\&); 

\quad A\& operator =(const A\&)=delete; 

\quad  A\& operator =(A\&\&)=delete;\};

struct Node\{Node *prev; Node *next; A a;\};

Написати функцію, яка вилучає зі списку всі вузли, елемент в яких більший елемента в попередньому вузлі списку або дорівнює елементу в наступному вузлі. Список є циклічним та задається вказівником на перший вузол списку. Наприклад, зі списку, що містить послідовність 8, 10, 9, 4, мають бути вилучені вузли, що містять 8 (бо список циклічний і $8>4$) та 10. У цьому прикладі 9 залишається в списку.

Крім правильності та відповідності умові оцінюється якість коду, наявність витоків ресурсів та коректність роботи з пам'яттю. 

%enditem
}
{\smallskip\smallskip \begin {scriptsize} Затверджено на засіданні кафедри теоретичної кібернетики, протокол \No 4  від   25.01.2024 р.\par\smallskip\smallskip\smallskip\smallskip
{\parbox {6.5 cm} {Зав. кафедри \dotfill Ю.В.~Крак}}\hspace {0.7cm} {\hbox to 6.5 cm { Екзаменатор \dotfill Т.О.~Карнаух}} \end{scriptsize}}
%end
\newpage
%begin
{{\begin {scriptsize}\par
{ \center  Київський національний університет імені Тараса Шевченка \par}
{\parbox {5 cm}{Освітня програма \hfill}{\parbox {7 cm}{Інформатика\hfill}}\parbox {6 cm}{\hfill Семестр 2}}\par
{\parbox {5 cm} {Навчальний предмет \hfill}\parbox {7 cm}{Програмування \hfill}\parbox {6cm}{\hfill}\par}\vspace{-15pt} \end{scriptsize}}{\center Екзаменаційний білет \No \textbf { 80 }\par}}
{%beginitem
1. Вкажіть основну відмінність між відкритими (public) та прихованими (private) методами класу.%enditem
\par
\vspace{2mm}
%beginitem
2. Максимально змістовно повно сформулюйте назву оголошеного в наведеному фрагменті методу класу.
\vspace{-3mm}
\begin{verbatim}class A{
   ...
   A& operator =(A&&);
};\end{verbatim}
\vspace{-3mm}
%enditem
\par
\vspace{2mm}
%beginitem
3. Вкажіть, що не так з наведеним кодом.
\vspace{-3mm}
\begin{verbatim}int *pa=new int(2), *pb=pa; delete pb; *pa=3;\end{verbatim}
\vspace{-3mm}
%enditem
\par
\vspace{2mm}
%beginitem
4. За наявності означень
\vspace{-3mm}
\begin{verbatim}
class A{ };
class B:public A{ };
class C:public A{B b;};\end{verbatim}
\vspace{-3mm}
вкажіть послідовність викликів деструкторів за виконання коду \code{\{C c;\}} 

(Якщо деструктор викликається кілька разів поспіль, то має зазначатися кожен його виклик, наприклад, $\sim$A, $\sim$A, $\sim$A.)
%enditem

\vspace{2mm}
%beginitem
5. Яку часову оцінку складності має алгоритм сортування вибором?
%enditem
\par
\vspace{2mm}
%beginitem
6. Модифікувати клас точок простору 

\qquad\code{class Point\{ double x, y, z;\};}


додавши до нього оператор перетворення на тип \code{bool}: точка є істинною тоді і тільки тоді, коли містить хоча б одну ненульову координату. Означення оператора теж має бути записане.%enditem
\par
\vspace{2mm}
%beginitem
7. 
Реалізувати операцію різниці множин з повтореннями. Множина з повтореннями задається переліком своїх елементів, записаних за неспаданням у текстовому файлі через білі символи. Результат у такому ж самому форматі (і теж за неспаданням) має бути записаний у текстовий файл. Питома функція має отримувати три шляхи до файлів. Розмір вхідних множин такий, що вони не можуть бути завантажені повністю в пам’ять програми.

Крім правильності та відповідності умові оцінюється якість коду, наявність витоків ресурсів та коректність роботи з пам'яттю. 

%enditem
}
{\smallskip\smallskip \begin {scriptsize} Затверджено на засіданні кафедри теоретичної кібернетики, протокол \No 4  від   25.01.2024 р.\par\smallskip\smallskip\smallskip\smallskip
{\parbox {6.5 cm} {Зав. кафедри \dotfill Ю.В.~Крак}}\hspace {0.7cm} {\hbox to 6.5 cm { Екзаменатор \dotfill Т.О.~Карнаух}} \end{scriptsize}}
%end
\newpage
%begin
{{\begin {scriptsize}\par
{ \center  Київський національний університет імені Тараса Шевченка \par}
{\parbox {5 cm}{Освітня програма \hfill}{\parbox {7 cm}{Інформатика\hfill}}\parbox {6 cm}{\hfill Семестр 2}}\par
{\parbox {5 cm} {Навчальний предмет \hfill}\parbox {7 cm}{Програмування \hfill}\parbox {6cm}{\hfill}\par}\vspace{-15pt} \end{scriptsize}}{\center Екзаменаційний білет \No \textbf { 81 }\par}}
{%beginitem
1. Вкажіть основну \textbf{синтаксичну} відмінність між конструктором копії та конструктором переміщення.%enditem
\par
\vspace{2mm}
%beginitem
2. Максимально змістовно повно сформулюйте назву оголошеного в наведеному фрагменті методу класу.
\vspace{-3mm}
\begin{verbatim}class A{
   ...
   A(A&&);
};\end{verbatim}
\vspace{-3mm}
%enditem
\par
\vspace{2mm}
%beginitem
3. Вкажіть, що не так з наведеним кодом.
\vspace{-3mm}
\begin{verbatim}int a, *p=&a; delete a;\end{verbatim}
\vspace{-3mm}
%enditem
\par
\vspace{2mm}
%beginitem
4. За наявності означень
\vspace{-3mm}
\begin{verbatim}
class A{ };
class B{ };
class C:public A{B b;A a;};\end{verbatim}
\vspace{-3mm}
вкажіть послідовність викликів деструкторів за виконання коду \code{\{C c;\}} 

(Якщо деструктор викликається кілька разів поспіль, то має зазначатися кожен його виклик, наприклад, $\sim$A, $\sim$A, $\sim$A.)
%enditem

\vspace{2mm}
%beginitem
5. Яку часову оцінку складності має алгоритм швидкого сортування?
%enditem
\par
\vspace{2mm}
%beginitem
6. Модифікувати клас точок простору 

\qquad\code{class Point\{ double x, y, z;\};}


додавши до нього оператор перетворення на тип \code{bool}: точка є істинною тоді і тільки тоді, коли містить хоча б одну ненульову координату. Означення оператора теж має бути записане.%enditem
\par
\vspace{2mm}
%beginitem
7. 
Задано клас A (структуру, оператори < та == десь коректно означено) та клас вузла списку Node.

struct A\{

\quad  якісь коректні означення 

\quad  bool operator<(const A\&); 

\quad A\& operator =(const A\&)=delete; 

\quad  A\& operator =(A\&\&)=delete;\};

struct Node\{Node *prev; Node *next; A a;\};

Написати функцію, яка вилучає зі списку всі вузли, елемент в яких менший елемента в попередньому вузлі списку або дорівнює елементу в наступному вузлі. Список задається вказівником на перший вузол списку. Наприклад, зі списку, що містить послідовність 8, 5, 6, 9, 4, мають бути вилучені вузли, що містять 5 та 4. У цьому прикладі 6 залишається в списку.

Крім правильності та відповідності умові оцінюється якість коду, наявність витоків ресурсів та коректність роботи з пам'яттю. 

%enditem
}
{\smallskip\smallskip \begin {scriptsize} Затверджено на засіданні кафедри теоретичної кібернетики, протокол \No 4  від   25.01.2024 р.\par\smallskip\smallskip\smallskip\smallskip
{\parbox {6.5 cm} {Зав. кафедри \dotfill Ю.В.~Крак}}\hspace {0.7cm} {\hbox to 6.5 cm { Екзаменатор \dotfill Т.О.~Карнаух}} \end{scriptsize}}
%end
\newpage
%begin
{{\begin {scriptsize}\par
{ \center  Київський національний університет імені Тараса Шевченка \par}
{\parbox {5 cm}{Освітня програма \hfill}{\parbox {7 cm}{Інформатика\hfill}}\parbox {6 cm}{\hfill Семестр 2}}\par
{\parbox {5 cm} {Навчальний предмет \hfill}\parbox {7 cm}{Програмування \hfill}\parbox {6cm}{\hfill}\par}\vspace{-15pt} \end{scriptsize}}{\center Екзаменаційний білет \No \textbf { 82 }\par}}
{%beginitem
1. Вкажіть основну відмінність між класами (class) та структурами (struct).%enditem
\par
\vspace{2mm}
%beginitem
2. Максимально змістовно повно сформулюйте назву оголошеного в наведеному фрагменті методу класу.
\vspace{-3mm}
\begin{verbatim}class A{
   ...
   A();
};\end{verbatim}
\vspace{-3mm}
%enditem
\par
\vspace{2mm}
%beginitem
3. Вкажіть, що не так з наведеним кодом.
\vspace{-3mm}
\begin{verbatim}int *pa=new int(2), *pb=new int(3); pa=pb;\end{verbatim}
\vspace{-3mm}
%enditem
\par
\vspace{2mm}
%beginitem
4. За наявності означень
\vspace{-3mm}
\begin{verbatim}
class A{ };
class B{ };
class C:public A{B b;A a;};\end{verbatim}
\vspace{-3mm}
вкажіть послідовність викликів конструкторів за виконання коду \code{\{C c;\}} 

(Якщо конструктор викликається кілька разів поспіль, то має зазначатися кожен його виклик, наприклад, A, A, A.)
%enditem

\vspace{2mm}
%beginitem
5. Яку часову оцінку складності має алгоритм сортування бульбашкою?
%enditem
\par
\vspace{2mm}
%beginitem
6. Модифікувати клас точок простору 

\qquad\code{class Point\{ double x, y, z;\};}


додавши до нього оператор перетворення на рядок типу \code{std::string}, що містить рядкове зображення точки простору (у круглих дужках координати через кому), наприклад, (1.01, 2.35, 0.5), та його означення.%enditem
\par
\vspace{2mm}
%beginitem
7. 
Задано клас A (структуру, оператори < та == десь коректно означено) та клас вузла списку Node.

struct A\{

\quad  якісь коректні означення 

\quad  bool operator<(const A\&); 

\quad  bool operator==(const A\&); 

\quad A\& operator =(const A\&)=delete; 

\quad  A\& operator =(A\&\&)=delete;\};

struct Node\{Node *next; A a;\};

 Написати функцію, яка вилучає зі списку всі вузли, елемент в яких більший елемента в попередньому вузлі списку або дорівнює елементу в наступному вузлі. Список є циклічним та задається вказівником на перший вузол списку. Наприклад, зі списку, що містить послідовність 8, 10, 9, 4, мають бути вилучені вузли, що містять 8 (бо список циклічний і $8>4$) та 10. У цьому прикладі 9 залишається в списку.

Крім правильності та відповідності умові оцінюється якість коду, наявність витоків ресурсів та коректність роботи з пам'яттю. 

%enditem
}
{\smallskip\smallskip \begin {scriptsize} Затверджено на засіданні кафедри теоретичної кібернетики, протокол \No 4  від   25.01.2024 р.\par\smallskip\smallskip\smallskip\smallskip
{\parbox {6.5 cm} {Зав. кафедри \dotfill Ю.В.~Крак}}\hspace {0.7cm} {\hbox to 6.5 cm { Екзаменатор \dotfill Т.О.~Карнаух}} \end{scriptsize}}
%end
\newpage
%begin
{{\begin {scriptsize}\par
{ \center  Київський національний університет імені Тараса Шевченка \par}
{\parbox {5 cm}{Освітня програма \hfill}{\parbox {7 cm}{Інформатика\hfill}}\parbox {6 cm}{\hfill Семестр 2}}\par
{\parbox {5 cm} {Навчальний предмет \hfill}\parbox {7 cm}{Програмування \hfill}\parbox {6cm}{\hfill}\par}\vspace{-15pt} \end{scriptsize}}{\center Екзаменаційний білет \No \textbf { 83 }\par}}
{%beginitem
1. Що важливо забезпечити при роботі з динамічними змінними?%enditem
\par
\vspace{2mm}
%beginitem
2. Максимально змістовно повно сформулюйте назву оголошеного в наведеному фрагменті методу класу.
\vspace{-3mm}
\begin{verbatim}class A{
   ...
   A& operator =(A&&);
};\end{verbatim}
\vspace{-3mm}
%enditem
\par
\vspace{2mm}
%beginitem
3. Вкажіть, що не так з наведеним кодом.
\vspace{-3mm}
\begin{verbatim}int *pa=new int(2), *pb=*pa; delete pa;\end{verbatim}
\vspace{-3mm}
%enditem
\par
\vspace{2mm}
%beginitem
4. За наявності означень
\vspace{-3mm}
\begin{verbatim}
class A{ };
class B{ };
class C{B b;A a;};\end{verbatim}
\vspace{-3mm}
вкажіть послідовність викликів деструкторів за виконання коду \code{\{C c;\}} 

(Якщо деструктор викликається кілька разів поспіль, то має зазначатися кожен його виклик, наприклад, $\sim$A, $\sim$A, $\sim$A.)
%enditem

\vspace{2mm}
%beginitem
5. Яку часову оцінку складності має алгоритм сортування вставками?
%enditem
\par
\vspace{2mm}
%beginitem
6. Модифікувати клас 

\qquad\code{class Session\{ int DM, MA, Alg, Prog;\};} 


додавши до нього оператор порівняння \code{<}: один об'єкт є меншим за інший тоді і тільки тоді, коли значення всіх атрибутів об'єктів не менші за 60 і сума атрибутів першого менша за суму атрибутів другого. Означення оператора теж має бути записане.%enditem
\par
\vspace{2mm}
%beginitem
7. 
Задано клас A (структуру, оператори > та == десь коректно означено) та клас вузла списку Node.

struct A\{

\quad  якісь коректні означення 

\quad  bool operator>(const A\&); 

\quad  bool operator==(const A\&); 

\quad A\& operator =(const A\&)=delete; 

\quad  A\& operator =(A\&\&)=delete;\};

struct Node\{Node *next; A a;\};

 Написати функцію, яка вилучає зі списку всі вузли, елемент в яких менший елемента в попередньому вузлі списку або дорівнює елементу в наступному вузлі. Список задається вказівником на перший вузол списку. Наприклад, зі списку, що містить послідовність 8, 5, 6, 9, 4, мають бути вилучені вузли, що містять 5 та 4. У цьому прикладі 6 залишається в списку.

Крім правильності та відповідності умові оцінюється якість коду, наявність витоків ресурсів та коректність роботи з пам'яттю. 

%enditem
}
{\smallskip\smallskip \begin {scriptsize} Затверджено на засіданні кафедри теоретичної кібернетики, протокол \No 4  від   25.01.2024 р.\par\smallskip\smallskip\smallskip\smallskip
{\parbox {6.5 cm} {Зав. кафедри \dotfill Ю.В.~Крак}}\hspace {0.7cm} {\hbox to 6.5 cm { Екзаменатор \dotfill Т.О.~Карнаух}} \end{scriptsize}}
%end
\newpage
%begin
{{\begin {scriptsize}\par
{ \center  Київський національний університет імені Тараса Шевченка \par}
{\parbox {5 cm}{Освітня програма \hfill}{\parbox {7 cm}{Інформатика\hfill}}\parbox {6 cm}{\hfill Семестр 2}}\par
{\parbox {5 cm} {Навчальний предмет \hfill}\parbox {7 cm}{Програмування \hfill}\parbox {6cm}{\hfill}\par}\vspace{-15pt} \end{scriptsize}}{\center Екзаменаційний білет \No \textbf { 84 }\par}}
{%beginitem
1. Вкажіть основну \textbf{семантичну} відмінність між конструктором копії та конструктором переміщення.%enditem
\par
\vspace{2mm}
%beginitem
2. Максимально змістовно повно сформулюйте назву оголошеного в наведеному фрагменті методу класу.
\vspace{-3mm}
\begin{verbatim}class A{
   ...
   ~A();
};\end{verbatim}
\vspace{-3mm}
%enditem
\par
\vspace{2mm}
%beginitem
3. Вкажіть, що не так з наведеним кодом.
\vspace{-3mm}
\begin{verbatim}int **ppi=new int*(new int(3)); delete ppi;\end{verbatim}
\vspace{-3mm}
%enditem
\par
\vspace{2mm}
%beginitem
4. За наявності означень
\vspace{-3mm}
\begin{verbatim}
class A{ };
class B{ };
class C:public A{B b;};\end{verbatim}
\vspace{-3mm}
вкажіть послідовність викликів конструкторів за виконання коду \code{\{C c;\}} 

(Якщо конструктор викликається кілька разів поспіль, то має зазначатися кожен його виклик, наприклад, A, A, A.)
%enditem

\vspace{2mm}
%beginitem
5. Яку часову оцінку складності має алгоритм сортування купою (він же деревом, пірамідою)?
%enditem
\par
\vspace{2mm}
%beginitem
6. Модифікувати клас 

\qquad\code{class Session\{ int DM, MA, Alg, Prog;\};} 


додавши до нього оператор перетворення на тип \code{int}: якщо хоча б один атрибут менший за 60, то має повертатись 0, інакше -- сума атрибутів.%enditem
\par
\vspace{2mm}
%beginitem
7. 
Реалізувати операцію злиття множин з повтореннями. Множина з повтореннями задається переліком своїх елементів, записаних за неспаданням у текстовому файлі через білі символи. Результат у такому ж самому форматі (і теж за неспаданням) має бути записаний у текстовий файл. Питома функція має отримувати три шляхи до файлів. Розмір вхідних множин такий, що вони не можуть бути завантажені повністю в пам’ять програми.

Крім правильності та відповідності умові оцінюється якість коду, наявність витоків ресурсів та коректність роботи з пам'яттю. 

%enditem
}
{\smallskip\smallskip \begin {scriptsize} Затверджено на засіданні кафедри теоретичної кібернетики, протокол \No 4  від   25.01.2024 р.\par\smallskip\smallskip\smallskip\smallskip
{\parbox {6.5 cm} {Зав. кафедри \dotfill Ю.В.~Крак}}\hspace {0.7cm} {\hbox to 6.5 cm { Екзаменатор \dotfill Т.О.~Карнаух}} \end{scriptsize}}
%end
\newpage
%begin
{{\begin {scriptsize}\par
{ \center  Київський національний університет імені Тараса Шевченка \par}
{\parbox {5 cm}{Освітня програма \hfill}{\parbox {7 cm}{Інформатика\hfill}}\parbox {6 cm}{\hfill Семестр 2}}\par
{\parbox {5 cm} {Навчальний предмет \hfill}\parbox {7 cm}{Програмування \hfill}\parbox {6cm}{\hfill}\par}\vspace{-15pt} \end{scriptsize}}{\center Екзаменаційний білет \No \textbf { 85 }\par}}
{%beginitem
1. Вкажіть основну \textbf{семантичну} відмінність між конструктором копії та конструктором переміщення.%enditem
\par
\vspace{2mm}
%beginitem
2. Максимально змістовно повно сформулюйте назву оголошеного в наведеному фрагменті методу класу.
\vspace{-3mm}
\begin{verbatim}class A{
   ...
   A(const A&);
};\end{verbatim}
\vspace{-3mm}
%enditem
\par
\vspace{2mm}
%beginitem
3. Вкажіть, що не так з наведеним кодом.
\vspace{-3mm}
\begin{verbatim}int a, *p=&a; delete a;\end{verbatim}
\vspace{-3mm}
%enditem
\par
\vspace{2mm}
%beginitem
4. За наявності означень
\vspace{-3mm}
\begin{verbatim}
class A{ };
class B:public A{ };
class C:public A{B b;};\end{verbatim}
\vspace{-3mm}
вкажіть послідовність викликів конструкторів за виконання коду \code{\{C c;\}} 

(Якщо конструктор викликається кілька разів поспіль, то має зазначатися кожен його виклик, наприклад, A, A, A.)
%enditem

\vspace{2mm}
%beginitem
5. Яку часову оцінку складності має алгоритм сортування бульбашкою?
%enditem
\par
\vspace{2mm}
%beginitem
6. Модифікувати клас 

\qquad\code{class Session\{ int DM, MA, Alg, Prog;\};} 


додавши до нього оператор перетворення на тип \code{bool}: об'єкт класу є істинним тоді і тільки тоді, коли значення всіх атрибутів не менші за 60. Означення оператора теж має бути записане.%enditem
\par
\vspace{2mm}
%beginitem
7. 
Реалізувати операцію симетричної різниці множин з повтореннями. Множина з повтореннями задається переліком своїх елементів, записаних за неспаданням у текстовому файлі через білі символи. Результат у такому ж самому форматі (і теж за неспаданням) має бути записаний у текстовий файл. Питома функція має отримувати три шляхи до файлів. Розмір вхідних множин такий, що вони не можуть бути завантажені повністю в пам’ять програми.

Крім правильності та відповідності умові оцінюється якість коду, наявність витоків ресурсів та коректність роботи з пам'яттю. 

%enditem
}
{\smallskip\smallskip \begin {scriptsize} Затверджено на засіданні кафедри теоретичної кібернетики, протокол \No 4  від   25.01.2024 р.\par\smallskip\smallskip\smallskip\smallskip
{\parbox {6.5 cm} {Зав. кафедри \dotfill Ю.В.~Крак}}\hspace {0.7cm} {\hbox to 6.5 cm { Екзаменатор \dotfill Т.О.~Карнаух}} \end{scriptsize}}
%end
\newpage
%begin
{{\begin {scriptsize}\par
{ \center  Київський національний університет імені Тараса Шевченка \par}
{\parbox {5 cm}{Освітня програма \hfill}{\parbox {7 cm}{Інформатика\hfill}}\parbox {6 cm}{\hfill Семестр 2}}\par
{\parbox {5 cm} {Навчальний предмет \hfill}\parbox {7 cm}{Програмування \hfill}\parbox {6cm}{\hfill}\par}\vspace{-15pt} \end{scriptsize}}{\center Екзаменаційний білет \No \textbf { 86 }\par}}
{%beginitem
1. Вкажіть основну відмінність між відкритими (public) та прихованими (private) методами класу.%enditem
\par
\vspace{2mm}
%beginitem
2. Максимально змістовно повно сформулюйте назву оголошеного в наведеному фрагменті методу класу.
\vspace{-3mm}
\begin{verbatim}class A{
   ...
   A& operator =(A&&);
};\end{verbatim}
\vspace{-3mm}
%enditem
\par
\vspace{2mm}
%beginitem
3. Вкажіть, що не так з наведеним кодом.
\vspace{-3mm}
\begin{verbatim}int *pa=new int(2), *pb=pa; delete pb; delete pa;\end{verbatim}
\vspace{-3mm}
%enditem
\par
\vspace{2mm}
%beginitem
4. За наявності означень
\vspace{-3mm}
\begin{verbatim}
class A{ };
class B:public A{ };
class C:public B{ };\end{verbatim}
\vspace{-3mm}
вкажіть послідовність викликів деструкторів за виконання коду \code{\{C c;\}} 

(Якщо деструктор викликається кілька разів поспіль, то має зазначатися кожен його виклик, наприклад, $\sim$A, $\sim$A, $\sim$A.)
%enditem

\vspace{2mm}
%beginitem
5. Яку часову оцінку складності має алгоритм сортування вставками?
%enditem
\par
\vspace{2mm}
%beginitem
6. Модифікувати клас 

\qquad\code{class Session\{ int DM, MA, Alg, Prog;\};} 


додавши до нього оператор порівняння \code{<}: один об'єкт є меншим за інший тоді і тільки тоді, коли значення всіх атрибутів об'єктів не менші за 60 і сума атрибутів першого менша за суму атрибутів другого. Означення оператора теж має бути записане.%enditem
\par
\vspace{2mm}
%beginitem
7. 
Реалізувати операцію перетину множин з повтореннями. Множина з повтореннями задається переліком своїх елементів, записаних за неспаданням у текстовому файлі через білі символи. Результат у такому ж самому форматі (і теж за неспаданням) має бути записаний у текстовий файл. Питома функція має отримувати три шляхи до файлів. Розмір вхідних множин такий, що вони не можуть бути завантажені повністю в пам’ять програми.

Крім правильності та відповідності умові оцінюється якість коду, наявність витоків ресурсів та коректність роботи з пам'яттю. 

%enditem
}
{\smallskip\smallskip \begin {scriptsize} Затверджено на засіданні кафедри теоретичної кібернетики, протокол \No 4  від   25.01.2024 р.\par\smallskip\smallskip\smallskip\smallskip
{\parbox {6.5 cm} {Зав. кафедри \dotfill Ю.В.~Крак}}\hspace {0.7cm} {\hbox to 6.5 cm { Екзаменатор \dotfill Т.О.~Карнаух}} \end{scriptsize}}
%end
\newpage
%begin
{{\begin {scriptsize}\par
{ \center  Київський національний університет імені Тараса Шевченка \par}
{\parbox {5 cm}{Освітня програма \hfill}{\parbox {7 cm}{Інформатика\hfill}}\parbox {6 cm}{\hfill Семестр 2}}\par
{\parbox {5 cm} {Навчальний предмет \hfill}\parbox {7 cm}{Програмування \hfill}\parbox {6cm}{\hfill}\par}\vspace{-15pt} \end{scriptsize}}{\center Екзаменаційний білет \No \textbf { 87 }\par}}
{%beginitem
1. Вкажіть основну відмінність між класами (class) та структурами (struct).%enditem
\par
\vspace{2mm}
%beginitem
2. Максимально змістовно повно сформулюйте назву оголошеного в наведеному фрагменті методу класу.
\vspace{-3mm}
\begin{verbatim}class A{
   ...
   A(const A&);
};\end{verbatim}
\vspace{-3mm}
%enditem
\par
\vspace{2mm}
%beginitem
3. Вкажіть, що не так з наведеним кодом.
\vspace{-3mm}
\begin{verbatim}int **ppi=new int*(new int(3)); delete ppi;\end{verbatim}
\vspace{-3mm}
%enditem
\par
\vspace{2mm}
%beginitem
4. За наявності означень
\vspace{-3mm}
\begin{verbatim}
class A{ };
class B:public A{ };
class C:public A{B b;};\end{verbatim}
\vspace{-3mm}
вкажіть послідовність викликів деструкторів за виконання коду \code{\{C c;\}} 

(Якщо деструктор викликається кілька разів поспіль, то має зазначатися кожен його виклик, наприклад, $\sim$A, $\sim$A, $\sim$A.)
%enditem

\vspace{2mm}
%beginitem
5. Яку часову оцінку складності має алгоритм сортування злиттям?
%enditem
\par
\vspace{2mm}
%beginitem
6. Модифікувати клас 

\qquad\code{class Session\{ int DM, MA, Alg, Prog;\};} 


додавши до нього оператор перетворення на тип \code{int}: якщо хоча б один атрибут менший за 60, то має повертатись 0, інакше -- сума атрибутів.%enditem
\par
\vspace{2mm}
%beginitem
7. 
Задано клас A (структуру, оператори < та == десь коректно означено) та клас вузла списку Node.

struct A\{

\quad  якісь коректні означення 

\quad  bool operator<(const A\&); 

\quad  bool operator==(const A\&); 

\quad A\& operator =(const A\&)=delete; 

\quad  A\& operator =(A\&\&)=delete;\};

struct Node\{Node *prev; Node *next; A a;\};

Написати функцію, яка вилучає зі списку всі вузли, елемент в яких більший елемента в попередньому вузлі списку або дорівнює елементу в наступному вузлі. Список є циклічним та задається вказівником на перший вузол списку. Наприклад, зі списку, що містить послідовність 8, 10, 9, 4, мають бути вилучені вузли, що містять 8 (бо список циклічний і $8>4$) та 10. У цьому прикладі 9 залишається в списку.

Крім правильності та відповідності умові оцінюється якість коду, наявність витоків ресурсів та коректність роботи з пам'яттю. 

%enditem
}
{\smallskip\smallskip \begin {scriptsize} Затверджено на засіданні кафедри теоретичної кібернетики, протокол \No 4  від   25.01.2024 р.\par\smallskip\smallskip\smallskip\smallskip
{\parbox {6.5 cm} {Зав. кафедри \dotfill Ю.В.~Крак}}\hspace {0.7cm} {\hbox to 6.5 cm { Екзаменатор \dotfill Т.О.~Карнаух}} \end{scriptsize}}
%end
\newpage
%begin
{{\begin {scriptsize}\par
{ \center  Київський національний університет імені Тараса Шевченка \par}
{\parbox {5 cm}{Освітня програма \hfill}{\parbox {7 cm}{Інформатика\hfill}}\parbox {6 cm}{\hfill Семестр 2}}\par
{\parbox {5 cm} {Навчальний предмет \hfill}\parbox {7 cm}{Програмування \hfill}\parbox {6cm}{\hfill}\par}\vspace{-15pt} \end{scriptsize}}{\center Екзаменаційний білет \No \textbf { 88 }\par}}
{%beginitem
1. Вкажіть основну \textbf{синтаксичну} відмінність між конструктором копії та конструктором переміщення.%enditem
\par
\vspace{2mm}
%beginitem
2. Максимально змістовно повно сформулюйте назву оголошеного в наведеному фрагменті методу класу.
\vspace{-3mm}
\begin{verbatim}class A{
   ...
   A(A&&);
};\end{verbatim}
\vspace{-3mm}
%enditem
\par
\vspace{2mm}
%beginitem
3. Вкажіть, що не так з наведеним кодом.
\vspace{-3mm}
\begin{verbatim}int *pa=new int(2), *pb=*pa; delete pa;\end{verbatim}
\vspace{-3mm}
%enditem
\par
\vspace{2mm}
%beginitem
4. За наявності означень
\vspace{-3mm}
\begin{verbatim}
class A{ };
class B{ };
class C{B b;A a;};\end{verbatim}
\vspace{-3mm}
вкажіть послідовність викликів деструкторів за виконання коду \code{\{C c;\}} 

(Якщо деструктор викликається кілька разів поспіль, то має зазначатися кожен його виклик, наприклад, $\sim$A, $\sim$A, $\sim$A.)
%enditem

\vspace{2mm}
%beginitem
5. Яку часову оцінку складності має алгоритм швидкого сортування?
%enditem
\par
\vspace{2mm}
%beginitem
6. Модифікувати клас точок простору 

\qquad\code{class Point\{ double x, y, z;\};}


додавши до нього оператор перетворення на рядок типу \code{std::string}, що містить рядкове зображення точки простору (у круглих дужках координати через кому), наприклад, (1.01, 2.35, 0.5), та його означення.%enditem
\par
\vspace{2mm}
%beginitem
7. 
Реалізувати операцію різниці множин з повтореннями. Множина з повтореннями задається переліком своїх елементів, записаних за неспаданням у текстовому файлі через білі символи. Результат у такому ж самому форматі (і теж за неспаданням) має бути записаний у текстовий файл. Питома функція має отримувати три шляхи до файлів. Розмір вхідних множин такий, що вони не можуть бути завантажені повністю в пам’ять програми.

Крім правильності та відповідності умові оцінюється якість коду, наявність витоків ресурсів та коректність роботи з пам'яттю. 

%enditem
}
{\smallskip\smallskip \begin {scriptsize} Затверджено на засіданні кафедри теоретичної кібернетики, протокол \No 4  від   25.01.2024 р.\par\smallskip\smallskip\smallskip\smallskip
{\parbox {6.5 cm} {Зав. кафедри \dotfill Ю.В.~Крак}}\hspace {0.7cm} {\hbox to 6.5 cm { Екзаменатор \dotfill Т.О.~Карнаух}} \end{scriptsize}}
%end
\newpage
%begin
{{\begin {scriptsize}\par
{ \center  Київський національний університет імені Тараса Шевченка \par}
{\parbox {5 cm}{Освітня програма \hfill}{\parbox {7 cm}{Інформатика\hfill}}\parbox {6 cm}{\hfill Семестр 2}}\par
{\parbox {5 cm} {Навчальний предмет \hfill}\parbox {7 cm}{Програмування \hfill}\parbox {6cm}{\hfill}\par}\vspace{-15pt} \end{scriptsize}}{\center Екзаменаційний білет \No \textbf { 89 }\par}}
{%beginitem
1. Вкажіть основну відмінність між віртуальними методами та тими, що не є віртуальними.%enditem
\par
\vspace{2mm}
%beginitem
2. Максимально змістовно повно сформулюйте назву оголошеного в наведеному фрагменті методу класу.
\vspace{-3mm}
\begin{verbatim}class A{
   ...
   ~A();
};\end{verbatim}
\vspace{-3mm}
%enditem
\par
\vspace{2mm}
%beginitem
3. Вкажіть, що не так з наведеним кодом.
\vspace{-3mm}
\begin{verbatim}int *pa=new int(2), *pb=new int(3); pa=pb;\end{verbatim}
\vspace{-3mm}
%enditem
\par
\vspace{2mm}
%beginitem
4. За наявності означень
\vspace{-3mm}
\begin{verbatim}
class A{ };
class B{ };
class C{B b;A a;};\end{verbatim}
\vspace{-3mm}
вкажіть послідовність викликів конструкторів за виконання коду \code{\{C c;\}} 

(Якщо конструктор викликається кілька разів поспіль, то має зазначатися кожен його виклик, наприклад, A, A, A.)
%enditem

\vspace{2mm}
%beginitem
5. Яку часову оцінку складності має алгоритм сортування купою (він же деревом, пірамідою)?
%enditem
\par
\vspace{2mm}
%beginitem
6. Модифікувати клас точок простору 

\qquad\code{class Point\{ double x, y, z;\};}


додавши до нього оператор перетворення на тип \code{bool}: точка є істинною тоді і тільки тоді, коли містить хоча б одну ненульову координату. Означення оператора теж має бути записане.%enditem
\par
\vspace{2mm}
%beginitem
7. 
Реалізувати операцію симетричної різниці множин з повтореннями. Множина з повтореннями задається переліком своїх елементів, записаних за неспаданням у текстовому файлі через білі символи. Результат у такому ж самому форматі (і теж за неспаданням) має бути записаний у текстовий файл. Питома функція має отримувати три шляхи до файлів. Розмір вхідних множин такий, що вони не можуть бути завантажені повністю в пам’ять програми.

Крім правильності та відповідності умові оцінюється якість коду, наявність витоків ресурсів та коректність роботи з пам'яттю. 

%enditem
}
{\smallskip\smallskip \begin {scriptsize} Затверджено на засіданні кафедри теоретичної кібернетики, протокол \No 4  від   25.01.2024 р.\par\smallskip\smallskip\smallskip\smallskip
{\parbox {6.5 cm} {Зав. кафедри \dotfill Ю.В.~Крак}}\hspace {0.7cm} {\hbox to 6.5 cm { Екзаменатор \dotfill Т.О.~Карнаух}} \end{scriptsize}}
%end
\newpage
%begin
{{\begin {scriptsize}\par
{ \center  Київський національний університет імені Тараса Шевченка \par}
{\parbox {5 cm}{Освітня програма \hfill}{\parbox {7 cm}{Інформатика\hfill}}\parbox {6 cm}{\hfill Семестр 2}}\par
{\parbox {5 cm} {Навчальний предмет \hfill}\parbox {7 cm}{Програмування \hfill}\parbox {6cm}{\hfill}\par}\vspace{-15pt} \end{scriptsize}}{\center Екзаменаційний білет \No \textbf { 90 }\par}}
{%beginitem
1. Що важливо забезпечити при роботі з динамічними змінними?%enditem
\par
\vspace{2mm}
%beginitem
2. Максимально змістовно повно сформулюйте назву оголошеного в наведеному фрагменті методу класу.
\vspace{-3mm}
\begin{verbatim}class A{
   ...
   A();
};\end{verbatim}
\vspace{-3mm}
%enditem
\par
\vspace{2mm}
%beginitem
3. Вкажіть, що не так з наведеним кодом.
\vspace{-3mm}
\begin{verbatim}int *pa=new int(2), *pb=&pa; delete pa;\end{verbatim}
\vspace{-3mm}
%enditem
\par
\vspace{2mm}
%beginitem
4. За наявності означень
\vspace{-3mm}
\begin{verbatim}
class A{ };
class B{ };
class C:public A{B b[2];};\end{verbatim}
\vspace{-3mm}
вкажіть послідовність викликів конструкторів за виконання коду \code{\{C c;\}} 

(Якщо конструктор викликається кілька разів поспіль, то має зазначатися кожен його виклик, наприклад, A, A, A.)
%enditem

\vspace{2mm}
%beginitem
5. Яку часову оцінку складності має алгоритм сортування вибором?
%enditem
\par
\vspace{2mm}
%beginitem
6. Модифікувати клас 

\qquad\code{class Session\{ int DM, MA, Alg, Prog;\};} 


додавши до нього оператор перетворення на тип \code{bool}: об'єкт класу є істинним тоді і тільки тоді, коли значення всіх атрибутів не менші за 60. Означення оператора теж має бути записане.%enditem
\par
\vspace{2mm}
%beginitem
7. 
Реалізувати операцію різниці множин з повтореннями. Множина з повтореннями задається переліком своїх елементів, записаних за неспаданням у текстовому файлі через білі символи. Результат у такому ж самому форматі (і теж за неспаданням) має бути записаний у текстовий файл. Питома функція має отримувати три шляхи до файлів. Розмір вхідних множин такий, що вони не можуть бути завантажені повністю в пам’ять програми.

Крім правильності та відповідності умові оцінюється якість коду, наявність витоків ресурсів та коректність роботи з пам'яттю. 

%enditem
}
{\smallskip\smallskip \begin {scriptsize} Затверджено на засіданні кафедри теоретичної кібернетики, протокол \No 4  від   25.01.2024 р.\par\smallskip\smallskip\smallskip\smallskip
{\parbox {6.5 cm} {Зав. кафедри \dotfill Ю.В.~Крак}}\hspace {0.7cm} {\hbox to 6.5 cm { Екзаменатор \dotfill Т.О.~Карнаух}} \end{scriptsize}}
%end
\newpage
%begin
{{\begin {scriptsize}\par
{ \center  Київський національний університет імені Тараса Шевченка \par}
{\parbox {5 cm}{Освітня програма \hfill}{\parbox {7 cm}{Інформатика\hfill}}\parbox {6 cm}{\hfill Семестр 2}}\par
{\parbox {5 cm} {Навчальний предмет \hfill}\parbox {7 cm}{Програмування \hfill}\parbox {6cm}{\hfill}\par}\vspace{-15pt} \end{scriptsize}}{\center Екзаменаційний білет \No \textbf { 91 }\par}}
{%beginitem
1. Що важливо забезпечити при роботі з динамічними змінними?%enditem
\par
\vspace{2mm}
%beginitem
2. Максимально змістовно повно сформулюйте назву оголошеного в наведеному фрагменті методу класу.
\vspace{-3mm}
\begin{verbatim}class A{
   ...
   A& operator =(A&&);
};\end{verbatim}
\vspace{-3mm}
%enditem
\par
\vspace{2mm}
%beginitem
3. Вкажіть, що не так з наведеним кодом.
\vspace{-3mm}
\begin{verbatim}int *pa=new int(2), *pb=pa; delete pb; *pa=3;\end{verbatim}
\vspace{-3mm}
%enditem
\par
\vspace{2mm}
%beginitem
4. За наявності означень
\vspace{-3mm}
\begin{verbatim}
class A{ };
class B{ };
class C:public A{B b[2];};\end{verbatim}
\vspace{-3mm}
вкажіть послідовність викликів деструкторів за виконання коду \code{\{C c;\}} 

(Якщо деструктор викликається кілька разів поспіль, то має зазначатися кожен його виклик, наприклад, $\sim$A, $\sim$A, $\sim$A.)
%enditem

\vspace{2mm}
%beginitem
5. Яку часову оцінку складності має алгоритм швидкого сортування?
%enditem
\par
\vspace{2mm}
%beginitem
6. Модифікувати клас 

\qquad\code{class Session\{ int DM, MA, Alg, Prog;\};} 


додавши до нього оператор перетворення на тип \code{bool}: об'єкт класу є істинним тоді і тільки тоді, коли значення всіх атрибутів не менші за 60. Означення оператора теж має бути записане.%enditem
\par
\vspace{2mm}
%beginitem
7. 
Задано клас A (структуру, оператори > та == десь коректно означено) та клас вузла списку Node.

struct A\{

\quad  якісь коректні означення 

\quad  bool operator>(const A\&); 

\quad  bool operator==(const A\&); 

\quad A\& operator =(const A\&)=delete; 

\quad  A\& operator =(A\&\&)=delete;\};

struct Node\{Node *next; A a;\};

 Написати функцію, яка вилучає зі списку всі вузли, елемент в яких менший елемента в попередньому вузлі списку або дорівнює елементу в наступному вузлі. Список задається вказівником на перший вузол списку. Наприклад, зі списку, що містить послідовність 8, 5, 6, 9, 4, мають бути вилучені вузли, що містять 5 та 4. У цьому прикладі 6 залишається в списку.

Крім правильності та відповідності умові оцінюється якість коду, наявність витоків ресурсів та коректність роботи з пам'яттю. 

%enditem
}
{\smallskip\smallskip \begin {scriptsize} Затверджено на засіданні кафедри теоретичної кібернетики, протокол \No 4  від   25.01.2024 р.\par\smallskip\smallskip\smallskip\smallskip
{\parbox {6.5 cm} {Зав. кафедри \dotfill Ю.В.~Крак}}\hspace {0.7cm} {\hbox to 6.5 cm { Екзаменатор \dotfill Т.О.~Карнаух}} \end{scriptsize}}
%end
\newpage
%begin
{{\begin {scriptsize}\par
{ \center  Київський національний університет імені Тараса Шевченка \par}
{\parbox {5 cm}{Освітня програма \hfill}{\parbox {7 cm}{Інформатика\hfill}}\parbox {6 cm}{\hfill Семестр 2}}\par
{\parbox {5 cm} {Навчальний предмет \hfill}\parbox {7 cm}{Програмування \hfill}\parbox {6cm}{\hfill}\par}\vspace{-15pt} \end{scriptsize}}{\center Екзаменаційний білет \No \textbf { 92 }\par}}
{%beginitem
1. Вкажіть основну відмінність між відкритими (public) та прихованими (private) методами класу.%enditem
\par
\vspace{2mm}
%beginitem
2. Максимально змістовно повно сформулюйте назву оголошеного в наведеному фрагменті методу класу.
\vspace{-3mm}
\begin{verbatim}class A{
   ...
   ~A();
};\end{verbatim}
\vspace{-3mm}
%enditem
\par
\vspace{2mm}
%beginitem
3. Вкажіть, що не так з наведеним кодом.
\vspace{-3mm}
\begin{verbatim}int *pa=new int(2), *pb=&pa; delete pa;\end{verbatim}
\vspace{-3mm}
%enditem
\par
\vspace{2mm}
%beginitem
4. За наявності означень
\vspace{-3mm}
\begin{verbatim}
class A{ };
class B{ };
class C:public A{B b;};\end{verbatim}
\vspace{-3mm}
вкажіть послідовність викликів конструкторів за виконання коду \code{\{C c;\}} 

(Якщо конструктор викликається кілька разів поспіль, то має зазначатися кожен його виклик, наприклад, A, A, A.)
%enditem

\vspace{2mm}
%beginitem
5. Яку часову оцінку складності має алгоритм сортування бульбашкою?
%enditem
\par
\vspace{2mm}
%beginitem
6. Модифікувати клас 

\qquad\code{class Session\{ int DM, MA, Alg, Prog;\};} 


додавши до нього оператор перетворення на тип \code{int}: якщо хоча б один атрибут менший за 60, то має повертатись 0, інакше -- сума атрибутів.%enditem
\par
\vspace{2mm}
%beginitem
7. 
Задано клас A (структуру, оператори < та == десь коректно означено) та клас вузла списку Node.

struct A\{

\quad  якісь коректні означення 

\quad  bool operator<(const A\&); 

\quad  bool operator==(const A\&); 

\quad A\& operator =(const A\&)=delete; 

\quad  A\& operator =(A\&\&)=delete;\};

struct Node\{Node *prev; Node *next; A a;\};

Написати функцію, яка вилучає зі списку всі вузли, елемент в яких більший елемента в попередньому вузлі списку або дорівнює елементу в наступному вузлі. Список є циклічним та задається вказівником на перший вузол списку. Наприклад, зі списку, що містить послідовність 8, 10, 9, 4, мають бути вилучені вузли, що містять 8 (бо список циклічний і $8>4$) та 10. У цьому прикладі 9 залишається в списку.

Крім правильності та відповідності умові оцінюється якість коду, наявність витоків ресурсів та коректність роботи з пам'яттю. 

%enditem
}
{\smallskip\smallskip \begin {scriptsize} Затверджено на засіданні кафедри теоретичної кібернетики, протокол \No 4  від   25.01.2024 р.\par\smallskip\smallskip\smallskip\smallskip
{\parbox {6.5 cm} {Зав. кафедри \dotfill Ю.В.~Крак}}\hspace {0.7cm} {\hbox to 6.5 cm { Екзаменатор \dotfill Т.О.~Карнаух}} \end{scriptsize}}
%end
\newpage
%begin
{{\begin {scriptsize}\par
{ \center  Київський національний університет імені Тараса Шевченка \par}
{\parbox {5 cm}{Освітня програма \hfill}{\parbox {7 cm}{Інформатика\hfill}}\parbox {6 cm}{\hfill Семестр 2}}\par
{\parbox {5 cm} {Навчальний предмет \hfill}\parbox {7 cm}{Програмування \hfill}\parbox {6cm}{\hfill}\par}\vspace{-15pt} \end{scriptsize}}{\center Екзаменаційний білет \No \textbf { 93 }\par}}
{%beginitem
1. Вкажіть основну відмінність між віртуальними методами та тими, що не є віртуальними.%enditem
\par
\vspace{2mm}
%beginitem
2. Максимально змістовно повно сформулюйте назву оголошеного в наведеному фрагменті методу класу.
\vspace{-3mm}
\begin{verbatim}class A{
   ...
   A();
};\end{verbatim}
\vspace{-3mm}
%enditem
\par
\vspace{2mm}
%beginitem
3. Вкажіть, що не так з наведеним кодом.
\vspace{-3mm}
\begin{verbatim}int *pa=new int(2), *pb=pa; delete pb; *pa=3;\end{verbatim}
\vspace{-3mm}
%enditem
\par
\vspace{2mm}
%beginitem
4. За наявності означень
\vspace{-3mm}
\begin{verbatim}
class A{ };
class B{ };
class C:public A{B b;A a;};\end{verbatim}
\vspace{-3mm}
вкажіть послідовність викликів деструкторів за виконання коду \code{\{C c;\}} 

(Якщо деструктор викликається кілька разів поспіль, то має зазначатися кожен його виклик, наприклад, $\sim$A, $\sim$A, $\sim$A.)
%enditem

\vspace{2mm}
%beginitem
5. Яку часову оцінку складності має алгоритм сортування купою (він же деревом, пірамідою)?
%enditem
\par
\vspace{2mm}
%beginitem
6. Модифікувати клас точок простору 

\qquad\code{class Point\{ double x, y, z;\};}


додавши до нього оператор перетворення на рядок типу \code{std::string}, що містить рядкове зображення точки простору (у круглих дужках координати через кому), наприклад, (1.01, 2.35, 0.5), та його означення.%enditem
\par
\vspace{2mm}
%beginitem
7. 
Задано клас A (структуру, оператори < та == десь коректно означено) та клас вузла списку Node.

struct A\{

\quad  якісь коректні означення 

\quad  bool operator<(const A\&); 

\quad  bool operator==(const A\&); 

\quad A\& operator =(const A\&)=delete; 

\quad  A\& operator =(A\&\&)=delete;\};

struct Node\{Node *next; A a;\};

 Написати функцію, яка вилучає зі списку всі вузли, елемент в яких більший елемента в попередньому вузлі списку або дорівнює елементу в наступному вузлі. Список є циклічним та задається вказівником на перший вузол списку. Наприклад, зі списку, що містить послідовність 8, 10, 9, 4, мають бути вилучені вузли, що містять 8 (бо список циклічний і $8>4$) та 10. У цьому прикладі 9 залишається в списку.

Крім правильності та відповідності умові оцінюється якість коду, наявність витоків ресурсів та коректність роботи з пам'яттю. 

%enditem
}
{\smallskip\smallskip \begin {scriptsize} Затверджено на засіданні кафедри теоретичної кібернетики, протокол \No 4  від   25.01.2024 р.\par\smallskip\smallskip\smallskip\smallskip
{\parbox {6.5 cm} {Зав. кафедри \dotfill Ю.В.~Крак}}\hspace {0.7cm} {\hbox to 6.5 cm { Екзаменатор \dotfill Т.О.~Карнаух}} \end{scriptsize}}
%end
\newpage
%begin
{{\begin {scriptsize}\par
{ \center  Київський національний університет імені Тараса Шевченка \par}
{\parbox {5 cm}{Освітня програма \hfill}{\parbox {7 cm}{Інформатика\hfill}}\parbox {6 cm}{\hfill Семестр 2}}\par
{\parbox {5 cm} {Навчальний предмет \hfill}\parbox {7 cm}{Програмування \hfill}\parbox {6cm}{\hfill}\par}\vspace{-15pt} \end{scriptsize}}{\center Екзаменаційний білет \No \textbf { 94 }\par}}
{%beginitem
1. Вкажіть основну \textbf{семантичну} відмінність між конструктором копії та конструктором переміщення.%enditem
\par
\vspace{2mm}
%beginitem
2. Максимально змістовно повно сформулюйте назву оголошеного в наведеному фрагменті методу класу.
\vspace{-3mm}
\begin{verbatim}class A{
   ...
   A(A&&);
};\end{verbatim}
\vspace{-3mm}
%enditem
\par
\vspace{2mm}
%beginitem
3. Вкажіть, що не так з наведеним кодом.
\vspace{-3mm}
\begin{verbatim}int *pa=new int(2), *pb=new int(3); pa=pb;\end{verbatim}
\vspace{-3mm}
%enditem
\par
\vspace{2mm}
%beginitem
4. За наявності означень
\vspace{-3mm}
\begin{verbatim}
class A{ };
class B{ };
class C:public A{B b;A a;};\end{verbatim}
\vspace{-3mm}
вкажіть послідовність викликів конструкторів за виконання коду \code{\{C c;\}} 

(Якщо конструктор викликається кілька разів поспіль, то має зазначатися кожен його виклик, наприклад, A, A, A.)
%enditem

\vspace{2mm}
%beginitem
5. Яку часову оцінку складності має алгоритм швидкого сортування?
%enditem
\par
\vspace{2mm}
%beginitem
6. Модифікувати клас 

\qquad\code{class Session\{ int DM, MA, Alg, Prog;\};} 


додавши до нього оператор порівняння \code{<}: один об'єкт є меншим за інший тоді і тільки тоді, коли значення всіх атрибутів об'єктів не менші за 60 і сума атрибутів першого менша за суму атрибутів другого. Означення оператора теж має бути записане.%enditem
\par
\vspace{2mm}
%beginitem
7. 
Реалізувати операцію злиття множин з повтореннями. Множина з повтореннями задається переліком своїх елементів, записаних за неспаданням у текстовому файлі через білі символи. Результат у такому ж самому форматі (і теж за неспаданням) має бути записаний у текстовий файл. Питома функція має отримувати три шляхи до файлів. Розмір вхідних множин такий, що вони не можуть бути завантажені повністю в пам’ять програми.

Крім правильності та відповідності умові оцінюється якість коду, наявність витоків ресурсів та коректність роботи з пам'яттю. 

%enditem
}
{\smallskip\smallskip \begin {scriptsize} Затверджено на засіданні кафедри теоретичної кібернетики, протокол \No 4  від   25.01.2024 р.\par\smallskip\smallskip\smallskip\smallskip
{\parbox {6.5 cm} {Зав. кафедри \dotfill Ю.В.~Крак}}\hspace {0.7cm} {\hbox to 6.5 cm { Екзаменатор \dotfill Т.О.~Карнаух}} \end{scriptsize}}
%end
\newpage
%begin
{{\begin {scriptsize}\par
{ \center  Київський національний університет імені Тараса Шевченка \par}
{\parbox {5 cm}{Освітня програма \hfill}{\parbox {7 cm}{Інформатика\hfill}}\parbox {6 cm}{\hfill Семестр 2}}\par
{\parbox {5 cm} {Навчальний предмет \hfill}\parbox {7 cm}{Програмування \hfill}\parbox {6cm}{\hfill}\par}\vspace{-15pt} \end{scriptsize}}{\center Екзаменаційний білет \No \textbf { 95 }\par}}
{%beginitem
1. Вкажіть основну \textbf{синтаксичну} відмінність між конструктором копії та конструктором переміщення.%enditem
\par
\vspace{2mm}
%beginitem
2. Максимально змістовно повно сформулюйте назву оголошеного в наведеному фрагменті методу класу.
\vspace{-3mm}
\begin{verbatim}class A{
   ...
   A(const A&);
};\end{verbatim}
\vspace{-3mm}
%enditem
\par
\vspace{2mm}
%beginitem
3. Вкажіть, що не так з наведеним кодом.
\vspace{-3mm}
\begin{verbatim}int **ppi=new int*(new int(3)); delete ppi;\end{verbatim}
\vspace{-3mm}
%enditem
\par
\vspace{2mm}
%beginitem
4. За наявності означень
\vspace{-3mm}
\begin{verbatim}
class A{ };
class B:public A{ };
class C:public B{ };\end{verbatim}
\vspace{-3mm}
вкажіть послідовність викликів конструкторів за виконання коду \code{\{C c;\}} 

(Якщо конструктор викликається кілька разів поспіль, то має зазначатися кожен його виклик, наприклад, A, A, A.)
%enditem

\vspace{2mm}
%beginitem
5. Яку часову оцінку складності має алгоритм сортування злиттям?
%enditem
\par
\vspace{2mm}
%beginitem
6. Модифікувати клас точок простору 

\qquad\code{class Point\{ double x, y, z;\};}


додавши до нього оператор перетворення на тип \code{bool}: точка є істинною тоді і тільки тоді, коли містить хоча б одну ненульову координату. Означення оператора теж має бути записане.%enditem
\par
\vspace{2mm}
%beginitem
7. 
Реалізувати операцію перетину множин з повтореннями. Множина з повтореннями задається переліком своїх елементів, записаних за неспаданням у текстовому файлі через білі символи. Результат у такому ж самому форматі (і теж за неспаданням) має бути записаний у текстовий файл. Питома функція має отримувати три шляхи до файлів. Розмір вхідних множин такий, що вони не можуть бути завантажені повністю в пам’ять програми.

Крім правильності та відповідності умові оцінюється якість коду, наявність витоків ресурсів та коректність роботи з пам'яттю. 

%enditem
}
{\smallskip\smallskip \begin {scriptsize} Затверджено на засіданні кафедри теоретичної кібернетики, протокол \No 4  від   25.01.2024 р.\par\smallskip\smallskip\smallskip\smallskip
{\parbox {6.5 cm} {Зав. кафедри \dotfill Ю.В.~Крак}}\hspace {0.7cm} {\hbox to 6.5 cm { Екзаменатор \dotfill Т.О.~Карнаух}} \end{scriptsize}}
%end
\newpage
%begin
{{\begin {scriptsize}\par
{ \center  Київський національний університет імені Тараса Шевченка \par}
{\parbox {5 cm}{Освітня програма \hfill}{\parbox {7 cm}{Інформатика\hfill}}\parbox {6 cm}{\hfill Семестр 2}}\par
{\parbox {5 cm} {Навчальний предмет \hfill}\parbox {7 cm}{Програмування \hfill}\parbox {6cm}{\hfill}\par}\vspace{-15pt} \end{scriptsize}}{\center Екзаменаційний білет \No \textbf { 96 }\par}}
{%beginitem
1. Вкажіть основну відмінність між класами (class) та структурами (struct).%enditem
\par
\vspace{2mm}
%beginitem
2. Максимально змістовно повно сформулюйте назву оголошеного в наведеному фрагменті методу класу.
\vspace{-3mm}
\begin{verbatim}class A{
   ...
   A& operator =(A&&);
};\end{verbatim}
\vspace{-3mm}
%enditem
\par
\vspace{2mm}
%beginitem
3. Вкажіть, що не так з наведеним кодом.
\vspace{-3mm}
\begin{verbatim}int *pa=new int(2), *pb=*pa; delete pa;\end{verbatim}
\vspace{-3mm}
%enditem
\par
\vspace{2mm}
%beginitem
4. За наявності означень
\vspace{-3mm}
\begin{verbatim}
class A{ };
class B{ };
class C:public A{B b;};\end{verbatim}
\vspace{-3mm}
вкажіть послідовність викликів деструкторів за виконання коду \code{\{C c;\}} 

(Якщо деструктор викликається кілька разів поспіль, то має зазначатися кожен його виклик, наприклад, $\sim$A, $\sim$A, $\sim$A.)
%enditem

\vspace{2mm}
%beginitem
5. Яку часову оцінку складності має алгоритм сортування вставками?
%enditem
\par
\vspace{2mm}
%beginitem
6. Модифікувати клас 

\qquad\code{class Session\{ int DM, MA, Alg, Prog;\};} 


додавши до нього оператор перетворення на тип \code{bool}: об'єкт класу є істинним тоді і тільки тоді, коли значення всіх атрибутів не менші за 60. Означення оператора теж має бути записане.%enditem
\par
\vspace{2mm}
%beginitem
7. 
Задано клас A (структуру, оператори < та == десь коректно означено) та клас вузла списку Node.

struct A\{

\quad  якісь коректні означення 

\quad  bool operator<(const A\&); 

\quad A\& operator =(const A\&)=delete; 

\quad  A\& operator =(A\&\&)=delete;\};

struct Node\{Node *prev; Node *next; A a;\};

Написати функцію, яка вилучає зі списку всі вузли, елемент в яких менший елемента в попередньому вузлі списку або дорівнює елементу в наступному вузлі. Список задається вказівником на перший вузол списку. Наприклад, зі списку, що містить послідовність 8, 5, 6, 9, 4, мають бути вилучені вузли, що містять 5 та 4. У цьому прикладі 6 залишається в списку.

Крім правильності та відповідності умові оцінюється якість коду, наявність витоків ресурсів та коректність роботи з пам'яттю. 

%enditem
}
{\smallskip\smallskip \begin {scriptsize} Затверджено на засіданні кафедри теоретичної кібернетики, протокол \No 4  від   25.01.2024 р.\par\smallskip\smallskip\smallskip\smallskip
{\parbox {6.5 cm} {Зав. кафедри \dotfill Ю.В.~Крак}}\hspace {0.7cm} {\hbox to 6.5 cm { Екзаменатор \dotfill Т.О.~Карнаух}} \end{scriptsize}}
%end
\newpage
%begin
{{\begin {scriptsize}\par
{ \center  Київський національний університет імені Тараса Шевченка \par}
{\parbox {5 cm}{Освітня програма \hfill}{\parbox {7 cm}{Інформатика\hfill}}\parbox {6 cm}{\hfill Семестр 2}}\par
{\parbox {5 cm} {Навчальний предмет \hfill}\parbox {7 cm}{Програмування \hfill}\parbox {6cm}{\hfill}\par}\vspace{-15pt} \end{scriptsize}}{\center Екзаменаційний білет \No \textbf { 97 }\par}}
{%beginitem
1. Вкажіть основну відмінність між віртуальними методами та тими, що не є віртуальними.%enditem
\par
\vspace{2mm}
%beginitem
2. Максимально змістовно повно сформулюйте назву оголошеного в наведеному фрагменті методу класу.
\vspace{-3mm}
\begin{verbatim}class A{
   ...
   A(A&&);
};\end{verbatim}
\vspace{-3mm}
%enditem
\par
\vspace{2mm}
%beginitem
3. Вкажіть, що не так з наведеним кодом.
\vspace{-3mm}
\begin{verbatim}int a, *p=&a; delete a;\end{verbatim}
\vspace{-3mm}
%enditem
\par
\vspace{2mm}
%beginitem
4. За наявності означень
\vspace{-3mm}
\begin{verbatim}
class A{ };
class B{ };
class C:public A{B b[2];};\end{verbatim}
\vspace{-3mm}
вкажіть послідовність викликів конструкторів за виконання коду \code{\{C c;\}} 

(Якщо конструктор викликається кілька разів поспіль, то має зазначатися кожен його виклик, наприклад, A, A, A.)
%enditem

\vspace{2mm}
%beginitem
5. Яку часову оцінку складності має алгоритм сортування вибором?
%enditem
\par
\vspace{2mm}
%beginitem
6. Модифікувати клас 

\qquad\code{class Session\{ int DM, MA, Alg, Prog;\};} 


додавши до нього оператор порівняння \code{<}: один об'єкт є меншим за інший тоді і тільки тоді, коли значення всіх атрибутів об'єктів не менші за 60 і сума атрибутів першого менша за суму атрибутів другого. Означення оператора теж має бути записане.%enditem
\par
\vspace{2mm}
%beginitem
7. 
Задано клас A (структуру, оператори < та == десь коректно означено) та клас вузла списку Node.

struct A\{

\quad  якісь коректні означення 

\quad  bool operator<(const A\&); 

\quad  bool operator==(const A\&); 

\quad A\& operator =(const A\&)=delete; 

\quad  A\& operator =(A\&\&)=delete;\};

struct Node\{Node *prev; Node *next; A a;\};

Написати функцію, яка вилучає зі списку всі вузли, елемент в яких більший елемента в попередньому вузлі списку або дорівнює елементу в наступному вузлі. Список є циклічним та задається вказівником на перший вузол списку. Наприклад, зі списку, що містить послідовність 8, 10, 9, 4, мають бути вилучені вузли, що містять 8 (бо список циклічний і $8>4$) та 10. У цьому прикладі 9 залишається в списку.

Крім правильності та відповідності умові оцінюється якість коду, наявність витоків ресурсів та коректність роботи з пам'яттю. 

%enditem
}
{\smallskip\smallskip \begin {scriptsize} Затверджено на засіданні кафедри теоретичної кібернетики, протокол \No 4  від   25.01.2024 р.\par\smallskip\smallskip\smallskip\smallskip
{\parbox {6.5 cm} {Зав. кафедри \dotfill Ю.В.~Крак}}\hspace {0.7cm} {\hbox to 6.5 cm { Екзаменатор \dotfill Т.О.~Карнаух}} \end{scriptsize}}
%end
\newpage
%begin
{{\begin {scriptsize}\par
{ \center  Київський національний університет імені Тараса Шевченка \par}
{\parbox {5 cm}{Освітня програма \hfill}{\parbox {7 cm}{Інформатика\hfill}}\parbox {6 cm}{\hfill Семестр 2}}\par
{\parbox {5 cm} {Навчальний предмет \hfill}\parbox {7 cm}{Програмування \hfill}\parbox {6cm}{\hfill}\par}\vspace{-15pt} \end{scriptsize}}{\center Екзаменаційний білет \No \textbf { 98 }\par}}
{%beginitem
1. Вкажіть основну відмінність між відкритими (public) та прихованими (private) методами класу.%enditem
\par
\vspace{2mm}
%beginitem
2. Максимально змістовно повно сформулюйте назву оголошеного в наведеному фрагменті методу класу.
\vspace{-3mm}
\begin{verbatim}class A{
   ...
   ~A();
};\end{verbatim}
\vspace{-3mm}
%enditem
\par
\vspace{2mm}
%beginitem
3. Вкажіть, що не так з наведеним кодом.
\vspace{-3mm}
\begin{verbatim}int *pa=new int(2), *pb=pa; delete pb; delete pa;\end{verbatim}
\vspace{-3mm}
%enditem
\par
\vspace{2mm}
%beginitem
4. За наявності означень
\vspace{-3mm}
\begin{verbatim}
class A{ };
class B{ };
class C:public A{B b;};\end{verbatim}
\vspace{-3mm}
вкажіть послідовність викликів деструкторів за виконання коду \code{\{C c;\}} 

(Якщо деструктор викликається кілька разів поспіль, то має зазначатися кожен його виклик, наприклад, $\sim$A, $\sim$A, $\sim$A.)
%enditem

\vspace{2mm}
%beginitem
5. Яку часову оцінку складності має алгоритм швидкого сортування?
%enditem
\par
\vspace{2mm}
%beginitem
6. Модифікувати клас точок простору 

\qquad\code{class Point\{ double x, y, z;\};}


додавши до нього оператор перетворення на рядок типу \code{std::string}, що містить рядкове зображення точки простору (у круглих дужках координати через кому), наприклад, (1.01, 2.35, 0.5), та його означення.%enditem
\par
\vspace{2mm}
%beginitem
7. 
Реалізувати операцію різниці множин з повтореннями. Множина з повтореннями задається переліком своїх елементів, записаних за неспаданням у текстовому файлі через білі символи. Результат у такому ж самому форматі (і теж за неспаданням) має бути записаний у текстовий файл. Питома функція має отримувати три шляхи до файлів. Розмір вхідних множин такий, що вони не можуть бути завантажені повністю в пам’ять програми.

Крім правильності та відповідності умові оцінюється якість коду, наявність витоків ресурсів та коректність роботи з пам'яттю. 

%enditem
}
{\smallskip\smallskip \begin {scriptsize} Затверджено на засіданні кафедри теоретичної кібернетики, протокол \No 4  від   25.01.2024 р.\par\smallskip\smallskip\smallskip\smallskip
{\parbox {6.5 cm} {Зав. кафедри \dotfill Ю.В.~Крак}}\hspace {0.7cm} {\hbox to 6.5 cm { Екзаменатор \dotfill Т.О.~Карнаух}} \end{scriptsize}}
%end
\newpage
%begin
{{\begin {scriptsize}\par
{ \center  Київський національний університет імені Тараса Шевченка \par}
{\parbox {5 cm}{Освітня програма \hfill}{\parbox {7 cm}{Інформатика\hfill}}\parbox {6 cm}{\hfill Семестр 2}}\par
{\parbox {5 cm} {Навчальний предмет \hfill}\parbox {7 cm}{Програмування \hfill}\parbox {6cm}{\hfill}\par}\vspace{-15pt} \end{scriptsize}}{\center Екзаменаційний білет \No \textbf { 99 }\par}}
{%beginitem
1. Вкажіть основну відмінність між класами (class) та структурами (struct).%enditem
\par
\vspace{2mm}
%beginitem
2. Максимально змістовно повно сформулюйте назву оголошеного в наведеному фрагменті методу класу.
\vspace{-3mm}
\begin{verbatim}class A{
   ...
   A();
};\end{verbatim}
\vspace{-3mm}
%enditem
\par
\vspace{2mm}
%beginitem
3. Вкажіть, що не так з наведеним кодом.
\vspace{-3mm}
\begin{verbatim}int *pa=new int(2), *pb=pa; delete pb; *pa=3;\end{verbatim}
\vspace{-3mm}
%enditem
\par
\vspace{2mm}
%beginitem
4. За наявності означень
\vspace{-3mm}
\begin{verbatim}
class A{ };
class B{ };
class C:public A{B b;};\end{verbatim}
\vspace{-3mm}
вкажіть послідовність викликів конструкторів за виконання коду \code{\{C c;\}} 

(Якщо конструктор викликається кілька разів поспіль, то має зазначатися кожен його виклик, наприклад, A, A, A.)
%enditem

\vspace{2mm}
%beginitem
5. Яку часову оцінку складності має алгоритм швидкого сортування?
%enditem
\par
\vspace{2mm}
%beginitem
6. Модифікувати клас точок простору 

\qquad\code{class Point\{ double x, y, z;\};}


додавши до нього оператор перетворення на тип \code{bool}: точка є істинною тоді і тільки тоді, коли містить хоча б одну ненульову координату. Означення оператора теж має бути записане.%enditem
\par
\vspace{2mm}
%beginitem
7. 
Реалізувати операцію перетину множин з повтореннями. Множина з повтореннями задається переліком своїх елементів, записаних за неспаданням у текстовому файлі через білі символи. Результат у такому ж самому форматі (і теж за неспаданням) має бути записаний у текстовий файл. Питома функція має отримувати три шляхи до файлів. Розмір вхідних множин такий, що вони не можуть бути завантажені повністю в пам’ять програми.

Крім правильності та відповідності умові оцінюється якість коду, наявність витоків ресурсів та коректність роботи з пам'яттю. 

%enditem
}
{\smallskip\smallskip \begin {scriptsize} Затверджено на засіданні кафедри теоретичної кібернетики, протокол \No 4  від   25.01.2024 р.\par\smallskip\smallskip\smallskip\smallskip
{\parbox {6.5 cm} {Зав. кафедри \dotfill Ю.В.~Крак}}\hspace {0.7cm} {\hbox to 6.5 cm { Екзаменатор \dotfill Т.О.~Карнаух}} \end{scriptsize}}
%end
\newpage
%begin
{{\begin {scriptsize}\par
{ \center  Київський національний університет імені Тараса Шевченка \par}
{\parbox {5 cm}{Освітня програма \hfill}{\parbox {7 cm}{Інформатика\hfill}}\parbox {6 cm}{\hfill Семестр 2}}\par
{\parbox {5 cm} {Навчальний предмет \hfill}\parbox {7 cm}{Програмування \hfill}\parbox {6cm}{\hfill}\par}\vspace{-15pt} \end{scriptsize}}{\center Екзаменаційний білет \No \textbf { 100 }\par}}
{%beginitem
1. Вкажіть основну \textbf{синтаксичну} відмінність між конструктором копії та конструктором переміщення.%enditem
\par
\vspace{2mm}
%beginitem
2. Максимально змістовно повно сформулюйте назву оголошеного в наведеному фрагменті методу класу.
\vspace{-3mm}
\begin{verbatim}class A{
   ...
   A(const A&);
};\end{verbatim}
\vspace{-3mm}
%enditem
\par
\vspace{2mm}
%beginitem
3. Вкажіть, що не так з наведеним кодом.
\vspace{-3mm}
\begin{verbatim}int a, *p=&a; delete a;\end{verbatim}
\vspace{-3mm}
%enditem
\par
\vspace{2mm}
%beginitem
4. За наявності означень
\vspace{-3mm}
\begin{verbatim}
class A{ };
class B{ };
class C:public A{B b[2];};\end{verbatim}
\vspace{-3mm}
вкажіть послідовність викликів деструкторів за виконання коду \code{\{C c;\}} 

(Якщо деструктор викликається кілька разів поспіль, то має зазначатися кожен його виклик, наприклад, $\sim$A, $\sim$A, $\sim$A.)
%enditem

\vspace{2mm}
%beginitem
5. Яку часову оцінку складності має алгоритм сортування злиттям?
%enditem
\par
\vspace{2mm}
%beginitem
6. Модифікувати клас 

\qquad\code{class Session\{ int DM, MA, Alg, Prog;\};} 


додавши до нього оператор перетворення на тип \code{int}: якщо хоча б один атрибут менший за 60, то має повертатись 0, інакше -- сума атрибутів.%enditem
\par
\vspace{2mm}
%beginitem
7. 
Задано клас A (структуру, оператори < та == десь коректно означено) та клас вузла списку Node.

struct A\{

\quad  якісь коректні означення 

\quad  bool operator<(const A\&); 

\quad  bool operator==(const A\&); 

\quad A\& operator =(const A\&)=delete; 

\quad  A\& operator =(A\&\&)=delete;\};

struct Node\{Node *next; A a;\};

 Написати функцію, яка вилучає зі списку всі вузли, елемент в яких більший елемента в попередньому вузлі списку або дорівнює елементу в наступному вузлі. Список є циклічним та задається вказівником на перший вузол списку. Наприклад, зі списку, що містить послідовність 8, 10, 9, 4, мають бути вилучені вузли, що містять 8 (бо список циклічний і $8>4$) та 10. У цьому прикладі 9 залишається в списку.

Крім правильності та відповідності умові оцінюється якість коду, наявність витоків ресурсів та коректність роботи з пам'яттю. 

%enditem
}
{\smallskip\smallskip \begin {scriptsize} Затверджено на засіданні кафедри теоретичної кібернетики, протокол \No 4  від   25.01.2024 р.\par\smallskip\smallskip\smallskip\smallskip
{\parbox {6.5 cm} {Зав. кафедри \dotfill Ю.В.~Крак}}\hspace {0.7cm} {\hbox to 6.5 cm { Екзаменатор \dotfill Т.О.~Карнаух}} \end{scriptsize}}
%end
\newpage
%begin
{{\begin {scriptsize}\par
{ \center  Київський національний університет імені Тараса Шевченка \par}
{\parbox {5 cm}{Освітня програма \hfill}{\parbox {7 cm}{Інформатика\hfill}}\parbox {6 cm}{\hfill Семестр 2}}\par
{\parbox {5 cm} {Навчальний предмет \hfill}\parbox {7 cm}{Програмування \hfill}\parbox {6cm}{\hfill}\par}\vspace{-15pt} \end{scriptsize}}{\center Екзаменаційний білет \No \textbf { 101 }\par}}
{%beginitem
1. Вкажіть основну \textbf{семантичну} відмінність між конструктором копії та конструктором переміщення.%enditem
\par
\vspace{2mm}
%beginitem
2. Максимально змістовно повно сформулюйте назву оголошеного в наведеному фрагменті методу класу.
\vspace{-3mm}
\begin{verbatim}class A{
   ...
   ~A();
};\end{verbatim}
\vspace{-3mm}
%enditem
\par
\vspace{2mm}
%beginitem
3. Вкажіть, що не так з наведеним кодом.
\vspace{-3mm}
\begin{verbatim}int *pa=new int(2), *pb=*pa; delete pa;\end{verbatim}
\vspace{-3mm}
%enditem
\par
\vspace{2mm}
%beginitem
4. За наявності означень
\vspace{-3mm}
\begin{verbatim}
class A{ };
class B{ };
class C{B b;A a;};\end{verbatim}
\vspace{-3mm}
вкажіть послідовність викликів конструкторів за виконання коду \code{\{C c;\}} 

(Якщо конструктор викликається кілька разів поспіль, то має зазначатися кожен його виклик, наприклад, A, A, A.)
%enditem

\vspace{2mm}
%beginitem
5. Яку часову оцінку складності має алгоритм швидкого сортування?
%enditem
\par
\vspace{2mm}
%beginitem
6. Модифікувати клас точок простору 

\qquad\code{class Point\{ double x, y, z;\};}


додавши до нього оператор перетворення на рядок типу \code{std::string}, що містить рядкове зображення точки простору (у круглих дужках координати через кому), наприклад, (1.01, 2.35, 0.5), та його означення.%enditem
\par
\vspace{2mm}
%beginitem
7. 
Задано клас A (структуру, оператори < та == десь коректно означено) та клас вузла списку Node.

struct A\{

\quad  якісь коректні означення 

\quad  bool operator<(const A\&); 

\quad A\& operator =(const A\&)=delete; 

\quad  A\& operator =(A\&\&)=delete;\};

struct Node\{Node *prev; Node *next; A a;\};

Написати функцію, яка вилучає зі списку всі вузли, елемент в яких менший елемента в попередньому вузлі списку або дорівнює елементу в наступному вузлі. Список задається вказівником на перший вузол списку. Наприклад, зі списку, що містить послідовність 8, 5, 6, 9, 4, мають бути вилучені вузли, що містять 5 та 4. У цьому прикладі 6 залишається в списку.

Крім правильності та відповідності умові оцінюється якість коду, наявність витоків ресурсів та коректність роботи з пам'яттю. 

%enditem
}
{\smallskip\smallskip \begin {scriptsize} Затверджено на засіданні кафедри теоретичної кібернетики, протокол \No 4  від   25.01.2024 р.\par\smallskip\smallskip\smallskip\smallskip
{\parbox {6.5 cm} {Зав. кафедри \dotfill Ю.В.~Крак}}\hspace {0.7cm} {\hbox to 6.5 cm { Екзаменатор \dotfill Т.О.~Карнаух}} \end{scriptsize}}
%end
\newpage
%begin
{{\begin {scriptsize}\par
{ \center  Київський національний університет імені Тараса Шевченка \par}
{\parbox {5 cm}{Освітня програма \hfill}{\parbox {7 cm}{Інформатика\hfill}}\parbox {6 cm}{\hfill Семестр 2}}\par
{\parbox {5 cm} {Навчальний предмет \hfill}\parbox {7 cm}{Програмування \hfill}\parbox {6cm}{\hfill}\par}\vspace{-15pt} \end{scriptsize}}{\center Екзаменаційний білет \No \textbf { 102 }\par}}
{%beginitem
1. Що важливо забезпечити при роботі з динамічними змінними?%enditem
\par
\vspace{2mm}
%beginitem
2. Максимально змістовно повно сформулюйте назву оголошеного в наведеному фрагменті методу класу.
\vspace{-3mm}
\begin{verbatim}class A{
   ...
   A(const A&);
};\end{verbatim}
\vspace{-3mm}
%enditem
\par
\vspace{2mm}
%beginitem
3. Вкажіть, що не так з наведеним кодом.
\vspace{-3mm}
\begin{verbatim}int *pa=new int(2), *pb=&pa; delete pa;\end{verbatim}
\vspace{-3mm}
%enditem
\par
\vspace{2mm}
%beginitem
4. За наявності означень
\vspace{-3mm}
\begin{verbatim}
class A{ };
class B{ };
class C:public A{B b;A a;};\end{verbatim}
\vspace{-3mm}
вкажіть послідовність викликів деструкторів за виконання коду \code{\{C c;\}} 

(Якщо деструктор викликається кілька разів поспіль, то має зазначатися кожен його виклик, наприклад, $\sim$A, $\sim$A, $\sim$A.)
%enditem

\vspace{2mm}
%beginitem
5. Яку часову оцінку складності має алгоритм сортування бульбашкою?
%enditem
\par
\vspace{2mm}
%beginitem
6. Модифікувати клас 

\qquad\code{class Session\{ int DM, MA, Alg, Prog;\};} 


додавши до нього оператор перетворення на тип \code{bool}: об'єкт класу є істинним тоді і тільки тоді, коли значення всіх атрибутів не менші за 60. Означення оператора теж має бути записане.%enditem
\par
\vspace{2mm}
%beginitem
7. 
Задано клас A (структуру, оператори > та == десь коректно означено) та клас вузла списку Node.

struct A\{

\quad  якісь коректні означення 

\quad  bool operator>(const A\&); 

\quad  bool operator==(const A\&); 

\quad A\& operator =(const A\&)=delete; 

\quad  A\& operator =(A\&\&)=delete;\};

struct Node\{Node *next; A a;\};

 Написати функцію, яка вилучає зі списку всі вузли, елемент в яких менший елемента в попередньому вузлі списку або дорівнює елементу в наступному вузлі. Список задається вказівником на перший вузол списку. Наприклад, зі списку, що містить послідовність 8, 5, 6, 9, 4, мають бути вилучені вузли, що містять 5 та 4. У цьому прикладі 6 залишається в списку.

Крім правильності та відповідності умові оцінюється якість коду, наявність витоків ресурсів та коректність роботи з пам'яттю. 

%enditem
}
{\smallskip\smallskip \begin {scriptsize} Затверджено на засіданні кафедри теоретичної кібернетики, протокол \No 4  від   25.01.2024 р.\par\smallskip\smallskip\smallskip\smallskip
{\parbox {6.5 cm} {Зав. кафедри \dotfill Ю.В.~Крак}}\hspace {0.7cm} {\hbox to 6.5 cm { Екзаменатор \dotfill Т.О.~Карнаух}} \end{scriptsize}}
%end
\newpage
%begin
{{\begin {scriptsize}\par
{ \center  Київський національний університет імені Тараса Шевченка \par}
{\parbox {5 cm}{Освітня програма \hfill}{\parbox {7 cm}{Інформатика\hfill}}\parbox {6 cm}{\hfill Семестр 2}}\par
{\parbox {5 cm} {Навчальний предмет \hfill}\parbox {7 cm}{Програмування \hfill}\parbox {6cm}{\hfill}\par}\vspace{-15pt} \end{scriptsize}}{\center Екзаменаційний білет \No \textbf { 103 }\par}}
{%beginitem
1. Вкажіть основну відмінність між віртуальними методами та тими, що не є віртуальними.%enditem
\par
\vspace{2mm}
%beginitem
2. Максимально змістовно повно сформулюйте назву оголошеного в наведеному фрагменті методу класу.
\vspace{-3mm}
\begin{verbatim}class A{
   ...
   A& operator =(A&&);
};\end{verbatim}
\vspace{-3mm}
%enditem
\par
\vspace{2mm}
%beginitem
3. Вкажіть, що не так з наведеним кодом.
\vspace{-3mm}
\begin{verbatim}int *pa=new int(2), *pb=pa; delete pb; delete pa;\end{verbatim}
\vspace{-3mm}
%enditem
\par
\vspace{2mm}
%beginitem
4. За наявності означень
\vspace{-3mm}
\begin{verbatim}
class A{ };
class B:public A{ };
class C:public A{B b;};\end{verbatim}
\vspace{-3mm}
вкажіть послідовність викликів конструкторів за виконання коду \code{\{C c;\}} 

(Якщо конструктор викликається кілька разів поспіль, то має зазначатися кожен його виклик, наприклад, A, A, A.)
%enditem

\vspace{2mm}
%beginitem
5. Яку часову оцінку складності має алгоритм сортування купою (він же деревом, пірамідою)?
%enditem
\par
\vspace{2mm}
%beginitem
6. Модифікувати клас 

\qquad\code{class Session\{ int DM, MA, Alg, Prog;\};} 


додавши до нього оператор порівняння \code{<}: один об'єкт є меншим за інший тоді і тільки тоді, коли значення всіх атрибутів об'єктів не менші за 60 і сума атрибутів першого менша за суму атрибутів другого. Означення оператора теж має бути записане.%enditem
\par
\vspace{2mm}
%beginitem
7. 
Реалізувати операцію симетричної різниці множин з повтореннями. Множина з повтореннями задається переліком своїх елементів, записаних за неспаданням у текстовому файлі через білі символи. Результат у такому ж самому форматі (і теж за неспаданням) має бути записаний у текстовий файл. Питома функція має отримувати три шляхи до файлів. Розмір вхідних множин такий, що вони не можуть бути завантажені повністю в пам’ять програми.

Крім правильності та відповідності умові оцінюється якість коду, наявність витоків ресурсів та коректність роботи з пам'яттю. 

%enditem
}
{\smallskip\smallskip \begin {scriptsize} Затверджено на засіданні кафедри теоретичної кібернетики, протокол \No 4  від   25.01.2024 р.\par\smallskip\smallskip\smallskip\smallskip
{\parbox {6.5 cm} {Зав. кафедри \dotfill Ю.В.~Крак}}\hspace {0.7cm} {\hbox to 6.5 cm { Екзаменатор \dotfill Т.О.~Карнаух}} \end{scriptsize}}
%end
\newpage
%begin
{{\begin {scriptsize}\par
{ \center  Київський національний університет імені Тараса Шевченка \par}
{\parbox {5 cm}{Освітня програма \hfill}{\parbox {7 cm}{Інформатика\hfill}}\parbox {6 cm}{\hfill Семестр 2}}\par
{\parbox {5 cm} {Навчальний предмет \hfill}\parbox {7 cm}{Програмування \hfill}\parbox {6cm}{\hfill}\par}\vspace{-15pt} \end{scriptsize}}{\center Екзаменаційний білет \No \textbf { 104 }\par}}
{%beginitem
1. Що важливо забезпечити при роботі з динамічними змінними?%enditem
\par
\vspace{2mm}
%beginitem
2. Максимально змістовно повно сформулюйте назву оголошеного в наведеному фрагменті методу класу.
\vspace{-3mm}
\begin{verbatim}class A{
   ...
   A();
};\end{verbatim}
\vspace{-3mm}
%enditem
\par
\vspace{2mm}
%beginitem
3. Вкажіть, що не так з наведеним кодом.
\vspace{-3mm}
\begin{verbatim}int **ppi=new int*(new int(3)); delete ppi;\end{verbatim}
\vspace{-3mm}
%enditem
\par
\vspace{2mm}
%beginitem
4. За наявності означень
\vspace{-3mm}
\begin{verbatim}
class A{ };
class B:public A{ };
class C:public B{ };\end{verbatim}
\vspace{-3mm}
вкажіть послідовність викликів конструкторів за виконання коду \code{\{C c;\}} 

(Якщо конструктор викликається кілька разів поспіль, то має зазначатися кожен його виклик, наприклад, A, A, A.)
%enditem

\vspace{2mm}
%beginitem
5. Яку часову оцінку складності має алгоритм сортування вставками?
%enditem
\par
\vspace{2mm}
%beginitem
6. Модифікувати клас точок простору 

\qquad\code{class Point\{ double x, y, z;\};}


додавши до нього оператор перетворення на тип \code{bool}: точка є істинною тоді і тільки тоді, коли містить хоча б одну ненульову координату. Означення оператора теж має бути записане.%enditem
\par
\vspace{2mm}
%beginitem
7. 
Реалізувати операцію злиття множин з повтореннями. Множина з повтореннями задається переліком своїх елементів, записаних за неспаданням у текстовому файлі через білі символи. Результат у такому ж самому форматі (і теж за неспаданням) має бути записаний у текстовий файл. Питома функція має отримувати три шляхи до файлів. Розмір вхідних множин такий, що вони не можуть бути завантажені повністю в пам’ять програми.

Крім правильності та відповідності умові оцінюється якість коду, наявність витоків ресурсів та коректність роботи з пам'яттю. 

%enditem
}
{\smallskip\smallskip \begin {scriptsize} Затверджено на засіданні кафедри теоретичної кібернетики, протокол \No 4  від   25.01.2024 р.\par\smallskip\smallskip\smallskip\smallskip
{\parbox {6.5 cm} {Зав. кафедри \dotfill Ю.В.~Крак}}\hspace {0.7cm} {\hbox to 6.5 cm { Екзаменатор \dotfill Т.О.~Карнаух}} \end{scriptsize}}
%end
\newpage
%begin
{{\begin {scriptsize}\par
{ \center  Київський національний університет імені Тараса Шевченка \par}
{\parbox {5 cm}{Освітня програма \hfill}{\parbox {7 cm}{Інформатика\hfill}}\parbox {6 cm}{\hfill Семестр 2}}\par
{\parbox {5 cm} {Навчальний предмет \hfill}\parbox {7 cm}{Програмування \hfill}\parbox {6cm}{\hfill}\par}\vspace{-15pt} \end{scriptsize}}{\center Екзаменаційний білет \No \textbf { 105 }\par}}
{%beginitem
1. Вкажіть основну \textbf{семантичну} відмінність між конструктором копії та конструктором переміщення.%enditem
\par
\vspace{2mm}
%beginitem
2. Максимально змістовно повно сформулюйте назву оголошеного в наведеному фрагменті методу класу.
\vspace{-3mm}
\begin{verbatim}class A{
   ...
   A(A&&);
};\end{verbatim}
\vspace{-3mm}
%enditem
\par
\vspace{2mm}
%beginitem
3. Вкажіть, що не так з наведеним кодом.
\vspace{-3mm}
\begin{verbatim}int *pa=new int(2), *pb=new int(3); pa=pb;\end{verbatim}
\vspace{-3mm}
%enditem
\par
\vspace{2mm}
%beginitem
4. За наявності означень
\vspace{-3mm}
\begin{verbatim}
class A{ };
class B{ };
class C:public A{B b;A a;};\end{verbatim}
\vspace{-3mm}
вкажіть послідовність викликів конструкторів за виконання коду \code{\{C c;\}} 

(Якщо конструктор викликається кілька разів поспіль, то має зазначатися кожен його виклик, наприклад, A, A, A.)
%enditem

\vspace{2mm}
%beginitem
5. Яку часову оцінку складності має алгоритм сортування вибором?
%enditem
\par
\vspace{2mm}
%beginitem
6. Модифікувати клас 

\qquad\code{class Session\{ int DM, MA, Alg, Prog;\};} 


додавши до нього оператор перетворення на тип \code{int}: якщо хоча б один атрибут менший за 60, то має повертатись 0, інакше -- сума атрибутів.%enditem
\par
\vspace{2mm}
%beginitem
7. 
Реалізувати операцію злиття множин з повтореннями. Множина з повтореннями задається переліком своїх елементів, записаних за неспаданням у текстовому файлі через білі символи. Результат у такому ж самому форматі (і теж за неспаданням) має бути записаний у текстовий файл. Питома функція має отримувати три шляхи до файлів. Розмір вхідних множин такий, що вони не можуть бути завантажені повністю в пам’ять програми.

Крім правильності та відповідності умові оцінюється якість коду, наявність витоків ресурсів та коректність роботи з пам'яттю. 

%enditem
}
{\smallskip\smallskip \begin {scriptsize} Затверджено на засіданні кафедри теоретичної кібернетики, протокол \No 4  від   25.01.2024 р.\par\smallskip\smallskip\smallskip\smallskip
{\parbox {6.5 cm} {Зав. кафедри \dotfill Ю.В.~Крак}}\hspace {0.7cm} {\hbox to 6.5 cm { Екзаменатор \dotfill Т.О.~Карнаух}} \end{scriptsize}}
%end
\newpage
%begin
{{\begin {scriptsize}\par
{ \center  Київський національний університет імені Тараса Шевченка \par}
{\parbox {5 cm}{Освітня програма \hfill}{\parbox {7 cm}{Інформатика\hfill}}\parbox {6 cm}{\hfill Семестр 2}}\par
{\parbox {5 cm} {Навчальний предмет \hfill}\parbox {7 cm}{Програмування \hfill}\parbox {6cm}{\hfill}\par}\vspace{-15pt} \end{scriptsize}}{\center Екзаменаційний білет \No \textbf { 106 }\par}}
{%beginitem
1. Вкажіть основну відмінність між відкритими (public) та прихованими (private) методами класу.%enditem
\par
\vspace{2mm}
%beginitem
2. Максимально змістовно повно сформулюйте назву оголошеного в наведеному фрагменті методу класу.
\vspace{-3mm}
\begin{verbatim}class A{
   ...
   A();
};\end{verbatim}
\vspace{-3mm}
%enditem
\par
\vspace{2mm}
%beginitem
3. Вкажіть, що не так з наведеним кодом.
\vspace{-3mm}
\begin{verbatim}int *pa=new int(2), *pb=*pa; delete pa;\end{verbatim}
\vspace{-3mm}
%enditem
\par
\vspace{2mm}
%beginitem
4. За наявності означень
\vspace{-3mm}
\begin{verbatim}
class A{ };
class B{ };
class C{B b;A a;};\end{verbatim}
\vspace{-3mm}
вкажіть послідовність викликів деструкторів за виконання коду \code{\{C c;\}} 

(Якщо деструктор викликається кілька разів поспіль, то має зазначатися кожен його виклик, наприклад, $\sim$A, $\sim$A, $\sim$A.)
%enditem

\vspace{2mm}
%beginitem
5. Яку часову оцінку складності має алгоритм швидкого сортування?
%enditem
\par
\vspace{2mm}
%beginitem
6. Модифікувати клас точок простору 

\qquad\code{class Point\{ double x, y, z;\};}


додавши до нього оператор перетворення на рядок типу \code{std::string}, що містить рядкове зображення точки простору (у круглих дужках координати через кому), наприклад, (1.01, 2.35, 0.5), та його означення.%enditem
\par
\vspace{2mm}
%beginitem
7. 
Задано клас A (структуру, оператори < та == десь коректно означено) та клас вузла списку Node.

struct A\{

\quad  якісь коректні означення 

\quad  bool operator<(const A\&); 

\quad  bool operator==(const A\&); 

\quad A\& operator =(const A\&)=delete; 

\quad  A\& operator =(A\&\&)=delete;\};

struct Node\{Node *prev; Node *next; A a;\};

Написати функцію, яка вилучає зі списку всі вузли, елемент в яких більший елемента в попередньому вузлі списку або дорівнює елементу в наступному вузлі. Список є циклічним та задається вказівником на перший вузол списку. Наприклад, зі списку, що містить послідовність 8, 10, 9, 4, мають бути вилучені вузли, що містять 8 (бо список циклічний і $8>4$) та 10. У цьому прикладі 9 залишається в списку.

Крім правильності та відповідності умові оцінюється якість коду, наявність витоків ресурсів та коректність роботи з пам'яттю. 

%enditem
}
{\smallskip\smallskip \begin {scriptsize} Затверджено на засіданні кафедри теоретичної кібернетики, протокол \No 4  від   25.01.2024 р.\par\smallskip\smallskip\smallskip\smallskip
{\parbox {6.5 cm} {Зав. кафедри \dotfill Ю.В.~Крак}}\hspace {0.7cm} {\hbox to 6.5 cm { Екзаменатор \dotfill Т.О.~Карнаух}} \end{scriptsize}}
%end
\newpage
%begin
{{\begin {scriptsize}\par
{ \center  Київський національний університет імені Тараса Шевченка \par}
{\parbox {5 cm}{Освітня програма \hfill}{\parbox {7 cm}{Інформатика\hfill}}\parbox {6 cm}{\hfill Семестр 2}}\par
{\parbox {5 cm} {Навчальний предмет \hfill}\parbox {7 cm}{Програмування \hfill}\parbox {6cm}{\hfill}\par}\vspace{-15pt} \end{scriptsize}}{\center Екзаменаційний білет \No \textbf { 107 }\par}}
{%beginitem
1. Вкажіть основну \textbf{синтаксичну} відмінність між конструктором копії та конструктором переміщення.%enditem
\par
\vspace{2mm}
%beginitem
2. Максимально змістовно повно сформулюйте назву оголошеного в наведеному фрагменті методу класу.
\vspace{-3mm}
\begin{verbatim}class A{
   ...
   A(const A&);
};\end{verbatim}
\vspace{-3mm}
%enditem
\par
\vspace{2mm}
%beginitem
3. Вкажіть, що не так з наведеним кодом.
\vspace{-3mm}
\begin{verbatim}int **ppi=new int*(new int(3)); delete ppi;\end{verbatim}
\vspace{-3mm}
%enditem
\par
\vspace{2mm}
%beginitem
4. За наявності означень
\vspace{-3mm}
\begin{verbatim}
class A{ };
class B:public A{ };
class C:public B{ };\end{verbatim}
\vspace{-3mm}
вкажіть послідовність викликів деструкторів за виконання коду \code{\{C c;\}} 

(Якщо деструктор викликається кілька разів поспіль, то має зазначатися кожен його виклик, наприклад, $\sim$A, $\sim$A, $\sim$A.)
%enditem

\vspace{2mm}
%beginitem
5. Яку часову оцінку складності має алгоритм сортування купою (він же деревом, пірамідою)?
%enditem
\par
\vspace{2mm}
%beginitem
6. Модифікувати клас точок простору 

\qquad\code{class Point\{ double x, y, z;\};}


додавши до нього оператор перетворення на тип \code{bool}: точка є істинною тоді і тільки тоді, коли містить хоча б одну ненульову координату. Означення оператора теж має бути записане.%enditem
\par
\vspace{2mm}
%beginitem
7. 
Задано клас A (структуру, оператори < та == десь коректно означено) та клас вузла списку Node.

struct A\{

\quad  якісь коректні означення 

\quad  bool operator<(const A\&); 

\quad A\& operator =(const A\&)=delete; 

\quad  A\& operator =(A\&\&)=delete;\};

struct Node\{Node *prev; Node *next; A a;\};

Написати функцію, яка вилучає зі списку всі вузли, елемент в яких менший елемента в попередньому вузлі списку або дорівнює елементу в наступному вузлі. Список задається вказівником на перший вузол списку. Наприклад, зі списку, що містить послідовність 8, 5, 6, 9, 4, мають бути вилучені вузли, що містять 5 та 4. У цьому прикладі 6 залишається в списку.

Крім правильності та відповідності умові оцінюється якість коду, наявність витоків ресурсів та коректність роботи з пам'яттю. 

%enditem
}
{\smallskip\smallskip \begin {scriptsize} Затверджено на засіданні кафедри теоретичної кібернетики, протокол \No 4  від   25.01.2024 р.\par\smallskip\smallskip\smallskip\smallskip
{\parbox {6.5 cm} {Зав. кафедри \dotfill Ю.В.~Крак}}\hspace {0.7cm} {\hbox to 6.5 cm { Екзаменатор \dotfill Т.О.~Карнаух}} \end{scriptsize}}
%end
\newpage
\end{document}
